%% !TEX program = xelatex
\documentclass[12pt]{article}

% Поддержка кириллицы и шрифтов
\usepackage{fontspec}
\usepackage[russian]{babel}
\usepackage{csquotes}

\usepackage{tabularx}
\usepackage[table]{xcolor} % ← для цвета в таблицах

% Библиография (стиль ГОСТ)
\usepackage[backend=biber, style=gost-numeric]{biblatex}
\addbibresource{references.bib} % создайте этот файл отдельно

% Шрифт по ГОСТ
\setmainfont{Times New Roman}

% Поля страницы (по умолчанию ~3 см; для точного ГОСТ используйте пакет geometry)
\usepackage[a4paper, margin=2.5cm]{geometry}

% Заголовки разделов без номеров (как в большинстве журналов)
\usepackage{titlesec}
\titleformat{\section}{\normalfont\bfseries\large}{\thesection}{1em}{}
\titlespacing*{\section}{0pt}{12pt}{6pt}

\begin{document}

\subsection{Название:}\label{ux43dux430ux437ux432ux430ux43dux438ux435}

\subsection{Авторы}\label{ux430ux432ux442ux43eux440ux44b}

\begin{itemize}
\item
  Самаркин А.И.
\item
  Самаркина Е.И.
\item
  Ревенко А.Б.
\item
  Щеткина
\item
  Иванова Н.В.
\item
  Белов В.С.
\end{itemize}

\subsection{1. Введение}\label{ux432ux432ux435ux434ux435ux43dux438ux435}

\subsubsection{1.1 Научный
контекст}\label{ux43dux430ux443ux447ux43dux44bux439-ux43aux43eux43dux442ux435ux43aux441ux442}

Пандемия COVID-19, официально зафиксированная Всемирной организацией
здравоохранения в марте 2020 года, потребовала от систем общественного
здравоохранения принципиально новых инструментов количественного анализа
эпидемического процесса. Математическое моделирование инфекционных
заболеваний, берущее начало в работах Кермака и Маккендрика {[}1{]}, к
этому моменту располагало устоявшейся методологической базой:
компартментные модели типа SIR, SEIR и их расширения широко применялись
для описания динамики гриппа, кори, лихорадки Эбола и ряда других
инфекций. Тем не менее опыт пандемии COVID-19 обнажил существенные
ограничения этого аппарата при работе с реальными данными рутинного
надзора.

Главная \textbf{практическая проблема} состоит в следующем: данные
государственной эпидемиологической отчётности накапливаются в
разнородных форматах, содержат систематические пропуски и не
предназначены для непосредственного использования в динамических
моделях. Между тем именно эти данные --- доступные в любом учреждении
санитарно-эпидемиологического надзора --- представляют наибольший
практический интерес: специализированные регистры и системы геномного
надзора существуют лишь в крупных административных центрах и охватывают
незначительную долю случаев. Для малых и средних городов России (с
численностью населения от 50 до 500 тысяч человек) задача
ретроспективного количественного анализа пандемии на основе стандартной
отчётности остаётся актуальной.

Помимо проблемы качества данных, пандемия COVID-19 поставила перед
математическим моделированием \textbf{новый содержательный вызов}.
Последовательная смена доминирующих вариантов SARS-CoV-2 --- от
исходного уханьского штамма через Альфа и Дельта к многочисленным
субвариантам Омикрон --- сопровождалась не только изменением
биологических параметров возбудителя (базового репродуктивного числа R₀,
длительности инкубационного и инфекционного периодов, летальности), но и
принципиальным изменением иммунологического контекста: к середине 2022
года значительная часть популяции обладала гибридным иммунитетом,
сформировавшимся в результате перенесённой инфекции и/или вакцинации
{[}2{]}, а также изменившемся социальным поведением, которое меняет
параметры распространения вируса, и, в свою очередь, существенно меняет
динамику волн: в частично иммунной популяции рост заболеваемости
замедляется не только вследствие истощения восприимчивых, но и под
влиянием сложных эффектов иммунологической памяти, действующих на
различных временных масштабах.

Стандартные компартментные модели с производными целого порядка не
располагают механизмом для описания такой «памяти» системы: текущее
состояние определяется исключительно мгновенными значениями переменных,
а не историей их изменения. Дробное исчисление, и в частности
производная Капуто порядка α ∈ (0, 1{]}, предоставляет математически
строгий инструмент для включения наследственных эффектов в динамические
модели {[}3{]}. В уравнении с производной Капуто правая часть
интегрирует взвешенную историю состояния системы с убывающим степенным
ядром памяти: при α → 1 модель вырождается в стандартную систему ОДУ,
при α \textless{} 1 --- \textbf{описывает} \textbf{субдиффузионный
режим}, при котором изменения происходят медленнее, чем предсказывает
экспоненциальная кинетика. Применительно к эпидемиологии это
соответствует ситуации, когда \textbf{восприимчивость популяции к
инфекции} \textbf{в данный момент определяется не только текущим числом
восприимчивых} \textbf{особей, но и накопленным иммунным опытом
предшествующих волн}.

Таким образом, \textbf{задача разработки методологии ретроспективного
анализа} \textbf{COVID-19 на уровне отдельного города}, опирающейся на
данные рутинного надзора и сочетающей покоординатную привязку параметров
к доминирующим вариантам вируса с дробным расширением компартментной
модели, \textbf{представляется} \textbf{как научно актуальной, так и
практически значимой}.

\begin{center}\rule{0.5\linewidth}{0.5pt}\end{center}

\textbf{Список литературы (к разделу 1.1)}

\begin{enumerate}
\def\labelenumi{\arabic{enumi}.}
\item
  Kermack W. O., McKendrick A. G. A contribution to the mathematical
  theory of epidemics // Proceedings of the Royal Society A. --- 1927.
  --- Vol. 115, № 772. --- P. 700--721. --- DOI: 10.1098/rspa.1927.0118.
\item
  Xu X., Wu Y., Kummer A. G. {[}et al.{]} Assessing changes in
  incubation period, serial interval, and generation time of SARS-CoV-2
  variants of concern: a systematic review and meta-analysis // BMC
  Medicine. ---

  \begin{enumerate}
  \def\labelenumii{\arabic{enumii}.}
  \setcounter{enumii}{2}

  \tightlist
  
  \item
    --- Vol. 21, № 1. --- Art. 374. --- DOI: 10.1186/s12916-023-03070-8.
  \end{enumerate}
\item
  Diethelm K. The Analysis of Fractional Differential Equations: An
  Application-Oriented Exposition Using Differential Operators of Caputo
  Type. --- Berlin : Springer, 2010. --- 247 p.~--- DOI:
  10.1007/978-3-642-14574-2.
\end{enumerate}

\subsubsection{1.2 Обзор
литературы}\label{ux43eux431ux437ux43eux440-ux43bux438ux442ux435ux440ux430ux442ux443ux440ux44b}

\paragraph{1.2.1 Классическая теория компартментных
моделей}\label{ux43aux43bux430ux441ux441ux438ux447ux435ux441ux43aux430ux44f-ux442ux435ux43eux440ux438ux44f-ux43aux43eux43cux43fux430ux440ux442ux43cux435ux43dux442ux43dux44bux445-ux43cux43eux434ux435ux43bux435ux439}

Математическое моделирование инфекционных заболеваний опирается на
\textbf{основополагающую работу Кермака и Маккендрика {[}1{]}, в которой
была} \textbf{предложена компартментная SIR-модель}, описывающая
динамику эпидемии через систему нелинейных ОДУ с тремя состояниями:
восприимчивые (S), инфекционные (I) и выздоровевшие (R). Несмотря на
очевидную упрощённость, эта работа заложила формальный фундамент
количественной эпидемиологии и ввела понятие базового репродуктивного
числа R₀ --- порога, определяющего возможность эпидемического
распространения. Принципиальное достоинство подхода состоит в
\textbf{аналитической} \textbf{прозрачности}: условие возникновения
эпидемии R₀ \textgreater{} 1 выводится строго из структуры модели.
Ограничением является гомогенное перемешивание популяции и постоянство
параметров передачи, \textbf{что делает SIR-модель} \textbf{адекватной
лишь для описания изолированной вспышки при постоянных} \textbf{условиях
среды}.

Монументальный вклад в теорию инфекционных болезней внесла
\textbf{монография} \textbf{Андерсона и Мэя {[}2{]}, систематизировавшая
компартментные модели} применительно к широкому кругу инфекций. Авторы
показали, что введение инкубационного периода (компартмент E)
существенно улучшает описание эпидемической кривой для инфекций с
длительной преинфекционной фазой, к числу которых относится и COVID-19.
Модель SEIR и её расширения --- включая вариант с явным учётом летальных
исходов (SEIRD) --- стали стандартом в эпидемиологическом моделировании
{[}2{]}. Тем не менее \textbf{критическим недостатком всей этой группы
моделей остаётся их} \textbf{«безпамятный» (иначе говоря марковский)
характер}: правая часть системы ОДУ в любой момент времени определяется
исключительно текущими значениями переменных состояния, тогда как в
реальной эпидемической динамике нарастание иммунитета, изменение
поведения популяции и смена вариантов возбудителя создают выраженную
зависимость от истории процесса.

Проблема идентифицируемости SEIRD-модели (точнее ее параметров) при
работе с реальными данными детально проанализирована в работе Королёва
{[}3{]}. Автор демонстрирует, что множество векторов параметров
обеспечивают практически одинаковое краткосрочное описание наблюдаемых
рядов заболеваемости и смертности, однако дают радикально различные
долгосрочные прогнозы --- принципиально важное обстоятельство для задач,
решаемых в настоящей работе. Дополнительно показывается, что
\textbf{неучёт занижения регистрируемых случаев приводит к
систематическому} \textbf{смещению оценки R₀ вниз}. Эти результаты
напрямую обусловили выбор стратегии калибровки в настоящем исследовании:
фиксацию биологических параметров (σ, γ) из литературных данных по
штаммам с последующей совместной оптимизацией β, μ и начальных условий.

\begin{center}\rule{0.5\linewidth}{0.5pt}\end{center}

\paragraph{1.2.2 Биологические параметры вариантов
SARS-CoV-2}\label{ux431ux438ux43eux43bux43eux433ux438ux447ux435ux441ux43aux438ux435-ux43fux430ux440ux430ux43cux435ux442ux440ux44b-ux432ux430ux440ux438ux430ux43dux442ux43eux432-sars-cov-2}

\textbf{Смена доминирующих вариантов SARS-CoV-2 сопровождалась
закономерным} \textbf{изменением ключевых эпидемиологических
параметров}, задокументированным в многочисленных когортных и
генеалогических исследованиях. Xu et al.~{[}4{]} провели систематический
обзор и метаанализ 147 исследований (280 записей) по инкубационному
периоду, серийному интервалу и времени генерации для предкового варианта
Wuhan, Alpha, Delta и Omicron. Ключевой вывод: каждый новый вариант
демонстрирует прогрессивное укорочение этих параметров --- с
\textasciitilde5--6 дней для уханьского штамма до \textasciitilde2--3
дней для Omicron. Недостатком этого систематического обзора является то,
что он не охватывает субварианты Omicron BA.4/5, XBB и поздние линии
JN.1, данные по которым появлялись в литературе позже. Для этих
субвариантов авторы прибегают к экстраполяции, что вносит
неопределённость в соответствующие оценки σ и γ.

Madewell et al.~{[}5{]} выполнили более специализированный метаанализ
серийных интервалов именно для Delta и Omicron, включив 46 648
первично-вторичных пар для Delta и 18 324 для Omicron. Этот источник
существенно точнее для задачи калибровки модели, поскольку оценки
получены на однородных когортах конкретных вариантов. \textbf{Авторы
указывают} \textbf{на среднее значение серийного интервала 2,3--5,8 дней
для Delta и} \textbf{2,1--4,8 дней для Omicron}, что согласуется с
данными, использованными в настоящей работе для инициализации параметра
γ.

\begin{center}\rule{0.5\linewidth}{0.5pt}\end{center}

\paragraph{1.2.3 Дробное исчисление производных: математический
аппарат}\label{ux434ux440ux43eux431ux43dux43eux435-ux438ux441ux447ux438ux441ux43bux435ux43dux438ux435-ux43fux440ux43eux438ux437ux432ux43eux434ux43dux44bux445-ux43cux430ux442ux435ux43cux430ux442ux438ux447ux435ux441ux43aux438ux439-ux430ux43fux43fux430ux440ux430ux442}

Дробные производные и интегралы произвольного вещественного порядка
восходят к переписке Лейбница и Лопиталя в 1695 году, однако
систематическое математическое развитие этот аппарат получил лишь в XIX
веке. Монография Подлубного {[}6{]} --- стандартный современный
справочник по теории, охватывающий определения Римана--Лиувилля,
Грюнвальда--Летникова и Капуто, теоремы существования и единственности,
аналитические и численные методы. Для целей настоящей работы особую
ценность представляет детальное сопоставление различных определений
дробной производной. В отличие от производной Римана--Лиувилля,
\textbf{производная Капуто допускает} \textbf{задание начальных условий
в классической форме (значения функции и её} \textbf{целочисленных
производных)}, что делает её предпочтительной при постановке задач Коши
для динамических систем.

Монография Дитхельма {[}7{]} представляет собой специализированный
прикладной трактат, полностью сосредоточенный на производных Капуто.
Автор строго доказывает существование, единственность и непрерывную
зависимость от начальных данных для широкого класса задач Коши с дробной
производной, обсуждает устойчивость равновесных состояний дробных систем
и даёт исчерпывающий анализ поведения решений в окрестности начального
момента --- в частности, показывает, что при α \textless{} 1 решение
развивается медленнее любой экспоненты, что соответствует
субдиффузионному режиму. Эта монография является непосредственным
методическим основанием для обоснования дробной SEIRD-модели,
используемой в настоящей работе.

Расширенный взгляд на численные методы для дробных ОДУ представлен в
монографии Балеану, Дитхельма, Скаласа и Трухильо {[}8{]}. Авторы дают
сравнительный анализ методов конечных разностей, квадратурных методов и
метода предиктора-корректора, указывая на их точность и вычислительную
сложность. В частности, обоснован метод Грюнвальда--Летникова как явная
схема первого порядка точности с ядром памяти, допускающим точное
рекуррентное вычисление. Это свойство является вычислительно
привлекательным при многократном решении ODE в ходе оптимизации, что
делает схему GL обоснованным выбором для настоящей работы, несмотря на
более низкий порядок по сравнению с методом Адамса--Башфорта--Мултона.

\begin{center}\rule{0.5\linewidth}{0.5pt}\end{center}

\paragraph{1.2.4 Схема предиктора-корректора для дробных
ОДУ}\label{ux441ux445ux435ux43cux430-ux43fux440ux435ux434ux438ux43aux442ux43eux440ux430-ux43aux43eux440ux440ux435ux43aux442ux43eux440ux430-ux434ux43bux44f-ux434ux440ux43eux431ux43dux44bux445-ux43eux434ux443}

Стандартным численным методом для дробных ОДУ в форме Капуто является
метод Адамса--Башфорта--Мултона (ABM), разработанный Дитхельмом, Фордом
и Фридом {[}9{]}. Схема обобщает классический предиктор-корректор
первого порядка на случай α ∈ (0, 1{]}: предиктор вычисляет приближение
y\^{}P\_\{n+1\} явно (аналог Адамса--Башфорта), корректор уточняет его
одной неявной итерацией (аналог Адамса--Мултона). Авторы доказывают
сходимость порядка O(h\^{}p) при p = min(2, 1 + α), то есть всегда не
ниже первого порядка и асимптотически второго при α → 1. Для задач
эпидемиологического моделирования с шагом h = 1 день разница в точности
между GL и ABM носит умеренный характер; тем не менее ABM является
предпочтительным методом при α \textless{} 0.7 или при необходимости
строгого контроля погрешности. В настоящей работе ABM не использовался
из соображений вычислительной простоты при оптимизации, однако
рассматривается как направление совершенствования метода.

\begin{center}\rule{0.5\linewidth}{0.5pt}\end{center}

\paragraph{1.2.5 Дробные компартментные модели
COVID-19}\label{ux434ux440ux43eux431ux43dux44bux435-ux43aux43eux43cux43fux430ux440ux442ux43cux435ux43dux442ux43dux44bux435-ux43cux43eux434ux435ux43bux438-covid-19}

Rajagopal et al.~{[}10{]} предложили дробную SEIRD-модель COVID-19 с
производной Капуто и откалибровали её на данных ВОЗ по Италии (24
февраля --- 7 апреля 2020 года). Сравнение со стандартной целочисленной
моделью показало меньшую RMSE у дробной версии; предсказательная
проверка на данных 8 апреля --- 16 мая подтвердила преимущество. Это
исследование --- наиболее близкий аналог настоящей работы по структуре
модели. Его принципиальным ограничением является использование данных
только первой волны COVID-19 (исходный вариант Wuhan), тогда как
настоящая работа охватывает 10 волн с 2020 по 2025 год при
последовательной смене доминирующих вариантов. Кроме того, авторы не
рассматривают вопрос о физической интерпретации подобранного значения α
и не связывают его с иммунологическим контекстом конкретных волн.

Naik et al.~{[}11{]} применили дробную модель с несколькими операторами
(Капуто, Атангана--Балеану--Капуто) к данным COVID-19 Пакистана, включив
в структуру компартмент лечения. Авторы демонстрируют, что оператор
Атангана--Балеану даёт лучшее согласие с реальными данными, чем
классический Капуто. Методологически важен вывод о существенном влиянии
выбора оператора дифференцирования на качество подгонки, что указывает
на необходимость сравнительного тестирования при калибровке дробных
моделей. Недостатком данной работы является отсутствие формальных
критериев выбора модели (AIC/BIC) и статистической проверки значимости α
\textless{} 1.

\begin{center}\rule{0.5\linewidth}{0.5pt}\end{center}

\paragraph{1.2.6 Инструментарий вычислительного
моделирования}\label{ux438ux43dux441ux442ux440ux443ux43cux435ux43dux442ux430ux440ux438ux439-ux432ux44bux447ux438ux441ux43bux438ux442ux435ux43bux44cux43dux43eux433ux43e-ux43cux43eux434ux435ux43bux438ux440ux43eux432ux430ux43dux438ux44f}

Численное решение систем ОДУ в настоящей работе осуществляется с
использованием пакета DifferentialEquations.jl для языка Julia {[}12{]}.
Rackauckas и Nie описывают экосистему, обеспечивающую единый интерфейс к
широкому спектру методов --- от явных методов Рунге--Кутты (Tsit5,
Vern7) до неявных решателей Розенброка (Rodas5P) --- с автоматическим
переключением между жёсткими и нежёсткими режимами (AutoVern7/Rodas5P).
Это свойство критически важно для настоящей работы, поскольку в ходе
многостартовой оптимизации параметров оптимизатор периодически пробует
конфигурации с большим β или малым γ, приводящие к жёстким системам,
которые разрушают явные решатели типа Tsit5. Применение
автопереключающегося решателя позволяет полностью исключить ложные
отказы интегратора при сохранении вычислительной эффективности на
нежёстких участках траектории.

Весь аналитический конвейер настоящей работы --- от обработки данных до
численной оптимизации параметров --- реализован на языке Julia {[}14{]}.
Julia представляет собой высокоуровневый динамический язык для научных
вычислений, архитектурно решающий так называемую «проблему двух языков»:
традиционная практика прототипирования на Python или MATLAB с
последующим переписыванием критических секций кода на C или Fortran
устранена благодаря компиляции Just-in-Time через LLVM и механизму
множественной диспетчеризации. Для задач настоящего исследования это
обстоятельство принципиально: многостартовая оптимизация параметров
SEIRD-модели требует сотен тысяч вызовов численного интегратора;
реализация на чистом Julia достигает производительности, сопоставимой с
Fortran, без использования внешних расширений. Дополнительным
достоинством является нативная поддержка Unicode в именах переменных и
функций, что позволяет непосредственно воспроизводить математические
обозначения в коде (β, σ, γ, μ, α\_k и т.д.) --- существенно улучшая
читаемость и снижая вероятность ошибок при переносе формул из публикации
в программу. Устойчивое управление пакетами и воспроизводимость
вычислительной среды обеспечиваются встроенным менеджером пакетов Pkg с
фиксацией версий зависимостей в файле Manifest.toml.

\begin{center}\rule{0.5\linewidth}{0.5pt}\end{center}

\subsubsection{Список литературы (к разделам
1.1--1.2)}\label{ux441ux43fux438ux441ux43eux43a-ux43bux438ux442ux435ux440ux430ux442ux443ux440ux44b-ux43a-ux440ux430ux437ux434ux435ux43bux430ux43c-1.11.2}

\begin{enumerate}
\def\labelenumi{\arabic{enumi}.}
\item
  Kermack W. O., McKendrick A. G. A contribution to the mathematical
  theory of epidemics // Proceedings of the Royal Society A. --- 1927.
  --- Vol. 115, № 772. --- P. 700--721. --- DOI:
  \href{https://doi.org/10.1098/rspa.1927.0118}{10.1098/rspa.1927.0118}.
\item
  Anderson R. M., May R. M. Infectious Diseases of Humans: Dynamics and
  Control. --- Oxford : Oxford University Press, 1991. --- 757 p.~---
  ISBN: 978-0-19-854040-3.
\item
  Korolev I. Identification and estimation of the SEIRD epidemic model
  for COVID-19 // Journal of Econometrics. --- 2021. --- Vol. 220, № 1.
  --- P. 63--85. --- DOI:
  \href{https://doi.org/10.1016/j.jeconom.2020.07.038}{10.1016/j.jeconom.2020.07.038}.
\item
  Xu X., Wu Y., Kummer A. G. {[}et al.{]} Assessing changes in
  incubation period, serial interval, and generation time of SARS-CoV-2
  variants of concern: a systematic review and meta-analysis // BMC
  Medicine. ---

  \begin{enumerate}
  \def\labelenumii{\arabic{enumii}.}
  \setcounter{enumii}{4}
  \tightlist
  \item
    --- Vol. 21, № 1. --- Art. 374. --- DOI:
    \href{https://doi.org/10.1186/s12916-023-03070-8}{10.1186/s12916-023-03070-8}.
  \end{enumerate}
\item
  Madewell Z. J., Yang Y., Longini I. M. Jr.~{[}et al.{]} Rapid review
  and meta-analysis of serial intervals for SARS-CoV-2 Delta and Omicron
  variants // BMC Infectious Diseases. --- 2023. --- Vol. 23, № 1. ---
  Art. 429. --- DOI:
  \href{https://doi.org/10.1186/s12879-023-08407-5}{10.1186/s12879-023-08407-5}.
\item
  Podlubny I. Fractional Differential Equations: An Introduction to
  Fractional Derivatives, Fractional Differential Equations, to Methods
  of Their Solution and Some of Their Applications. --- San Diego :
  Academic Press, 1999. --- 340 p.~--- (Mathematics in Science and
  Engineering ; vol.~198). --- ISBN: 978-0-12-558840-9. --- DOI:
  \href{https://doi.org/10.1016/s0076-5392(99)x8001-5}{10.1016/s0076-5392(99)x8001-5}.
\item
  Diethelm K. The Analysis of Fractional Differential Equations: An
  Application-Oriented Exposition Using Differential Operators of Caputo
  Type. --- Berlin : Springer, 2010. --- 247 p.~--- (Lecture Notes in
  Mathematics ; vol.~2004). --- DOI:
  \href{https://doi.org/10.1007/978-3-642-14574-2}{10.1007/978-3-642-14574-2}.
\item
  Baleanu D., Diethelm K., Scalas E., Trujillo J. J. Fractional
  Calculus: Models and Numerical Methods. --- Singapore : World
  Scientific, 2012. --- 400 p.~--- (Series on Complexity, Nonlinearity
  and Chaos ; vol.~3). --- DOI:
  \href{https://doi.org/10.1142/10044}{10.1142/10044}.
\item
  Diethelm K., Ford N. J., Freed A. D. A predictor-corrector approach
  for the numerical solution of fractional differential equations //
  Nonlinear Dynamics. --- 2002. --- Vol. 29, № 1--4. --- P. 3--22. ---
  DOI:
  \href{https://doi.org/10.1023/A:1016592219341}{10.1023/A:1016592219341}.
\item
  Rajagopal K., Hasanzadeh N., Parastesh F. {[}et al.{]} A
  fractional-order model for the novel coronavirus (COVID-19) outbreak
  // Nonlinear Dynamics. --- 2020. --- Vol. 101, № 1. --- P. 711--718.
  --- DOI:
  \href{https://doi.org/10.1007/s11071-020-05757-6}{10.1007/s11071-020-05757-6}.
\item
  Naik P. A., Yavuz M., Qureshi S. {[}et al.{]} Modeling and analysis of
  COVID-19 epidemics with treatment in fractional derivatives using real
  data from Pakistan // European Physical Journal Plus. --- 2020. ---
  Vol. 135, № 10. --- Art. 795. --- DOI:
  \href{https://doi.org/10.1140/epjp/s13360-020-00819-5}{10.1140/epjp/s13360-020-00819-5}.
\item
  Rackauckas C., Nie Q. DifferentialEquations.jl --- a performant and
  feature-rich ecosystem for solving differential equations in Julia //
  Journal of Open Research Software. --- 2017. --- Vol. 5, № 1. --- Art.
  15. --- DOI:
  \href{https://doi.org/10.5334/jors.151}{10.5334/jors.151}.
\item
  Bezanson J., Edelman A., Karpinski S., Shah V. B. Julia: a fresh
  approach to numerical computing // SIAM Review. --- 2017. --- Vol. 59,
  № 1. --- P. 65--98. --- DOI:
  \href{https://doi.org/10.1137/141000671}{10.1137/141000671}.
\end{enumerate}

\subsubsection{1.3 Пробелы в имеющихся
исследованиях}\label{ux43fux440ux43eux431ux435ux43bux44b-ux432-ux438ux43cux435ux44eux449ux438ux445ux441ux44f-ux438ux441ux441ux43bux435ux434ux43eux432ux430ux43dux438ux44fux445}

\paragraph{1.3.1 Дефицит субнационального
масштаба}\label{ux434ux435ux444ux438ux446ux438ux442-ux441ux443ux431ux43dux430ux446ux438ux43eux43dux430ux43bux44cux43dux43eux433ux43e-ux43cux430ux441ux448ux442ux430ux431ux430}

Подавляющее большинство публикаций по математическому моделированию
COVID-19 оперирует данными национального или регионального уровня:
анализируются страны, штаты, провинции {[}3, 10, 11{]}. Это обусловлено
объективными причинами --- доступностью агрегированных данных, ---
однако порождает принципиальный аналитический пробел. Как показывает
обзор Гайторп и соавт. {[}14{]}, \textbf{субнациональный-масштаб
является приоритетным} \textbf{для систем эпидемиологического надзора},
поскольку именно на уровне конкретных населённых пунктов принимаются
управленческие решения о медицинских ресурсах, введении ограничительных
мер и их отмене. Между тем \textbf{динамика инфекции в малом городе
качественно отличается от} \textbf{динамики в крупном административном
центре}: меньший размер восприимчивой популяции ускоряет исчерпание
иммунной прослойки, эффект «суперраспространителей» статистически
значимее, а сезонные миграционные потоки непропорционально сильнее
влияют на форму эпидемической кривой. В российских источниках
\textbf{тема моделирования} \textbf{COVID-19 на уровне отдельного
российского города} с населением менее 300 000 человек по данным
рутинной отчётности \textbf{практически не освещена}.

Немногочисленные исследования, работающие на субрегиональном уровне, как
правило, либо используют специализированные регистры (данные
PCR-тестирования с геопривязкой, данные госпитализаций), недоступные за
пределами крупных академических центров {[}14{]}, либо ограничиваются
пространственно-описательным анализом без калибровки компартментных
моделей {[}15{]}. \textbf{Исследований, в которых компартментная модель
была бы} \textbf{откалибрована по данным стандартных форм
государственной отчётности} \textbf{(аналог российских форм ВП и
еженедельных сводок COVID-19) для} \textbf{отдельного малого города,
авторам настоящей работы обнаружить} \textbf{не удалось.}

\begin{center}\rule{0.5\linewidth}{0.5pt}\end{center}

\paragraph{1.3.2 Разрыв между биологическими параметрами
штаммов}\label{ux440ux430ux437ux440ux44bux432-ux43cux435ux436ux434ux443-ux431ux438ux43eux43bux43eux433ux438ux447ux435ux441ux43aux438ux43cux438-ux43fux430ux440ux430ux43cux435ux442ux440ux430ux43cux438-ux448ux442ux430ux43cux43cux43eux432}

\paragraph{и калибровкой
волн}\label{ux438-ux43aux430ux43bux438ux431ux440ux43eux432ux43aux43eux439-ux432ux43eux43bux43d}

Накопленная литература по биологическим параметрам вариантов SARS-CoV-2
{[}4, 5{]} представляет собой глобальные сводные оценки σ, γ, R₀ и IFR,
полученные на крупных международных когортах. При этом в работах по
математическому моделированию COVID-19 эти параметры, как правило, либо
\textbf{заимствуются как фиксированные константы без адаптации к}
\textbf{конкретной волне} {[}10{]}, либо подбираются целиком из данных,
\textbf{игнорируя} \textbf{биологический приоритет} \textbf{свойств
конкретного штамма} {[}3{]}.

\textbf{Промежуточная стратегия --- использование параметров штамма как
информированного априорного распределения с последующим ограниченным
числом свободных параметров} подгонки --- реализована в единичных
работах и, как правило, без формальной проверки, насколько реальные
данные отличаются от биологического prior (то есть без отчётности о Δβ и
ΔR₀ между prior и fit).

Кроме того, \textbf{ни одна известная нам работа не проводит
систематического} \textbf{анализа по всем волнам одного и того же
наблюдаемого ряда с явной} \textbf{привязкой каждой волны к конкретному
доминирующему вариант}у и последующим сравнением подобранных параметров
с биологическими оценками этого варианта. Именно этот «пооолновой»
подход является ядром настоящего исследования.

\begin{center}\rule{0.5\linewidth}{0.5pt}\end{center}

\paragraph{1.3.3 Дробные модели: формальная
селекция}\label{ux434ux440ux43eux431ux43dux44bux435-ux43cux43eux434ux435ux43bux438-ux444ux43eux440ux43cux430ux43bux44cux43dux430ux44f-ux441ux435ux43bux435ux43aux446ux438ux44f}

\paragraph{и биологическая интерпретация
α\_k}\label{ux438-ux431ux438ux43eux43bux43eux433ux438ux447ux435ux441ux43aux430ux44f-ux438ux43dux442ux435ux440ux43fux440ux435ux442ux430ux446ux438ux44f-ux3b1_k}

Дробные SEIRD-модели COVID-19, как показано в обзоре выше {[}10, 11,
16{]}, демонстрируют эмпирически лучшую подгонку к данным по сравнению с
целочисленными аналогами. Тем не менее в этой группе работ
\textbf{обнаруживаются три системных недостатка}.

\textbf{Первый --- методологический}: большинство работ сравнивает
модели визуально (кривые подгонки) или по одному скалярному показателю
(RMSE или R²), не применяя формальных критериев выбора модели --- AIC,
BIC, Likelihood Ratio Test --- которые учитывают разное число параметров
у целочисленной и дробной версий. Работа Ли и соавт. {[}16{]} является
исключением: авторы систематически применяют AICc и BIC к пяти
компартментным структурам в целочисленном и дробном вариантах на данных
четырёх китайских провинций, показывая последовательное превосходство
дробных моделей. Однако авторы ограничиваются первой волной COVID-19
(начало 2020 года) и не рассматривают, как оптимальный порядок α
варьируется между волнами при смене вариантов вируса.

\textbf{Второй недостаток --- отсутствие содержательной интерпретации}
подобранного значения α\_k. В большинстве работ α трактуется как
технический параметр сглаживания или как «степень памяти системы» в
абстрактном смысле. Связь α\_k с конкретными популяционными процессами
--- нарастанием иммунитета, изменением возрастной структуры контактов,
накоплением гибридного иммунитета --- не формализована и не
верифицирована на реальных данных через сравнение волн с известным
иммунологическим контекстом.

\textbf{Третий недостаток --- однородность временного горизонта}:
практически все работы анализируют одну волну или один непрерывный ряд
(нередко 2--4 месяца), что не позволяет сравнить, является ли α
\textless{} 1 специфическим для конкретных волн или представляет собой
устойчивое свойство всей эпидемической кривой. Настоящая работа
ликвидирует этот пробел, анализируя десять последовательных волн на
одном и том же городе с применением формальных критериев выбора модели к
каждой волне.

\begin{center}\rule{0.5\linewidth}{0.5pt}\end{center}

\paragraph{1.3.4 Воспроизводимые конвейеры для рутинных
данных}\label{ux432ux43eux441ux43fux440ux43eux438ux437ux432ux43eux434ux438ux43cux44bux435-ux43aux43eux43dux432ux435ux439ux435ux440ux44b-ux434ux43bux44f-ux440ux443ux442ux438ux43dux43dux44bux445-ux434ux430ux43dux43dux44bux445}

Существующие работы по калибровке компартментных моделей, как правило,
не публикуют воспроизводимого конвейера предобработки данных, пригодного
для применения к другим источникам без значительной ручной адаптации.
Данные обычно либо предоставлены специальными регистрами с
унифицированным форматом {[}14{]}, \textbf{либо взяты из}
\textbf{агрегаторов типа OurWorldInData или Johns Hopkins CSSE, которые
к} \textbf{моменту написания ряда работ уже прекратили обновление}.
Задача создания воспроизводимого аналитического конвейера для работы с
разнородными файлами государственной эпидемиологической отчётности --- с
характерными для неё слитыми заголовками Excel, пропусками и сменой
форматов по годам --- \textbf{остаётся нерешённой для}
\textbf{большинства практических приложений} в контексте регионального
надзора в России.

\begin{center}\rule{0.5\linewidth}{0.5pt}\end{center}

\paragraph{1.3.5 Итог: формулировка исследовательского
зазора}\label{ux438ux442ux43eux433-ux444ux43eux440ux43cux443ux43bux438ux440ux43eux432ux43aux430-ux438ux441ux441ux43bux435ux434ux43eux432ux430ux442ux435ux43bux44cux441ux43aux43eux433ux43e-ux437ux430ux437ux43eux440ux430}

Таким образом, в совокупности литература демонстрирует следующие лакуны,
которые заполняет настоящее исследование.

\textbf{Во-первых, ретроспективный анализ всего пятилетнего периода
пандемии} \textbf{(2020--2025) на уровне одного российского города
(Псков,} \textbf{\textasciitilde200 000 жителей)} по данным
исключительно рутинного санитарно-эпидемиологического надзора --- без
специализированных регистров и геномного секвенирования.

\textbf{Во-вторых, покоординатная привязка параметров σ и γ к
доминирующему} \textbf{варианту вируса по дате наблюдения} с последующим
формальным сравнением подобранных значений β и R₀ с биологическим prior
из мета-аналитической литературы.

\textbf{В-третьих, одновременный прогон и формальное сравнение
целочисленной} \textbf{SEIRD и дробной SEIRD-модели} для каждой из
десяти волн (см. настоящую работу далее) с применением AIC, BIC и LRT
--- включая явное тестирование гипотезы о значимости α\_k \textless{} 1
для каждой отдельной волны.

В-четвёртых, эмпирическое сопоставление подобранных значений α\_k с
известным иммунологическим контекстом соответствующих волн, выдвижение и
верификация гипотезы об α\_k как косвенном маркере накопленного
популяционного иммунитета.

\begin{center}\rule{0.5\linewidth}{0.5pt}\end{center}

\subsubsection{Дополнительные источники (к разделу
1.3)}\label{ux434ux43eux43fux43eux43bux43dux438ux442ux435ux43bux44cux43dux44bux435-ux438ux441ux442ux43eux447ux43dux438ux43aux438-ux43a-ux440ux430ux437ux434ux435ux43bux443-1.3}

\begin{enumerate}
\def\labelenumi{\arabic{enumi}.}
\setcounter{enumi}{13}
\item
  Gaythorpe K. A. M., Imai N., Perez-Guzman P. N. {[}et al.{]} Data
  needs for better surveillance and response to infectious disease
  threats // Infectious Diseases of Poverty. --- 2023. --- Vol. 12, № 1.
  --- Art. 43. --- DOI:
  \href{https://doi.org/10.1186/s40249-023-01086-z}{10.1186/s40249-023-01086-z}.
\item
  Schmitz J., Lakes T., Manafa G. {[}et al.{]} Exploration of the
  COVID-19 pandemic at the neighborhood level in an intra-urban setting
  // Frontiers in Public Health. --- 2023. --- Vol. 11. --- Art.
  1128452. --- DOI:
  \href{https://doi.org/10.3389/fpubh.2023.1128452}{10.3389/fpubh.2023.1128452}.
\item
  Ли С., Сюй Ю., Ван Х. {[}Li S., Xu Y., Wang H. et al.{]} Multi-model
  selection and analysis for COVID-19 // Fractal and Fractional. ---
  2021. --- Vol. 5, № 3. --- Art. 120. --- DOI:
  \href{https://doi.org/10.3390/fractalfract5030120}{10.3390/fractalfract5030120}.
\end{enumerate}

\subsubsection{1.4 Цель и задачи
исследования}\label{ux446ux435ux43bux44c-ux438-ux437ux430ux434ux430ux447ux438-ux438ux441ux441ux43bux435ux434ux43eux432ux430ux43dux438ux44f}

Настоящее исследование преследует двойную цель: методологическую и
содержательную. \textbf{С методологической точки} зрения цель состоит в
разработке и апробации воспроизводимого аналитического конвейера
ретроспективного анализа пандемии COVID-19 на основе данных рутинного
санитарно-эпидемиологического надзора, применимого к малым и средним
городам без специализированных регистров и геномного мониторинга.
\textbf{С содержательной} --- в проверке гипотезы о том, что порядок
дробной производной α\_k является функционально информативным
параметром, отражающим не артефакт математической параметризации, а
реальное нарастание популяционного иммунного фона в ходе
последовательной смены волн пандемии.

Для достижения поставленной цели решались следующие задачи.

\textbf{Задача 1. Формирование непрерывного ежедневного ряда
заболеваемости.} Исходные данные государственного надзора (формы ВП и
оперативные сводки COVID-19 за 2020--2025 гг., г. Псков)
\textbf{приведены к единому} \textbf{формату, прошли унификацию
структуры и машинную импутацию} \textbf{пропусков} методом натуральных
кубических сплайнов с ограничением монотонности накопленного числа
случаев. Качество восстановленного ряда оценивалось по величине
систематической ошибки на контрольных точках с известными значениями.

\textbf{Задача 2. Автоматическая идентификация волн заболеваемости.}
Разработан \textbf{алгоритм разметки волн на основе знака сглаженной
первой} \textbf{производной} ряда заболеваемости (гауссово ядро, σ = 15
дней) с автоматическим присвоением номеров периодам непрерывного роста.
Для каждой волны рассчитаны характеристики: дата и величина пика,
продолжительность, площадь под кривой (AUC), ширина на полувысоте (FWHM)
и индекс асимметрии.

\textbf{Задача 3. Сопоставление волн с доминирующими вариантами вируса.}
Составлен и верифицирован по литературным данным {[}4, 5{]}
\textbf{календарь} \textbf{доминирования вариантов SARS-CoV-2}
применительно к периоду 2020--2025 гг. на территории России. Каждому дню
наблюдения присвоены биологические параметры соответствующего варианта
(σ, γ, μ, β, R₀), сформировавшие информированный prior для последующей
калибровки моделей.

\textbf{Задача 4. Калибровка стандартной SEIRD-модели.} Для каждой из
десяти идентифицированных волн проведена численная калибровка
целочисленной SEIRD-модели (решатель AutoVern7/Rodas5P пакета
DifferencialEquationns.jl), свободные параметры: β, μ, E₀, I₀, N\_eff).
Разработана двухэтапная стратегия оптимизации (сеточный поиск →
NelderMead → L-BFGS) с аналитическим предобусловливанием стартовой точки
по скорости роста начальной фазы волны. Функция потерь ---
логарифмический MSE, обеспечивающий равный вклад всех фаз волны в
критерий оптимальности.

\textbf{Задача 5. Калибровка дробной SEIRD-модели.} Разработана и
откалибрована дробная SEIRD-модель с производной Капуто порядка α\_k ∈
(0.5, 1.0{]} и дискретизацией по схеме Грюнвальда--Летникова. Указанная
схема, при незначительном снижении точности, обладает высокой скоростью
расчетов и лучше встраивается в конвейер оптимизации. Для каждой волны
проведены два независимых прогона: с фиксированным α\_k = 1
(целочисленная GL-версия) и с α\_k как свободным параметром (дробная
версия). Параметр α\_k введён через сигмоидальное преобразование φ =
logit((α\_k − 0.5)/0.5), гарантирующее задачу без ограничений при
оптимизации.

\textbf{Задача 6. Формальное сравнение моделей.} Для каждой волны
вычислены метрики качества подгонки (RMSE, MAE, MAPE, R²) и
информационные критерии (AIC, BIC) с учётом различного числа свободных
параметров у целочисленной и дробной версий. Значимость дополнительного
параметра α\_k проверена посредством теста отношения правдоподобия
(LRT): статистика Λ = n·(log MSE\_int − log MSE\_frac) \textasciitilde{}
χ²(1). Дополнительно проведён анализ структуры остатков: ACF до 21 лага
и Q-Q-графики для проверки нормальности.

\textbf{Задача 7. Содержательная интерпретация α\_k.} Выполнено
сравнение подобранных значений α\_k с иммунологическим и
эпидемиологическим контекстом соответствующих волн: охватом вакцинации,
порядковым номером волны (как прокси накопленного естественного
иммунитета), биологическими характеристиками доминирующего варианта
(иммунный уклон, IFR). Сформулирована и эмпирически проверена гипотеза
об α\_k \textless{} 1 как маркере субдиффузионного режима,
обусловленного популяционной иммунной памятью.

\textbf{Объект и предмет исследования.} Объект --- динамика
заболеваемости COVID-19 в г. Пскове (население \textasciitilde200 000
чел.) за 2020--2025 гг. Предмет --- сравнительная эффективность
целочисленной и дробной SEIRD-моделей в ретроспективном описании
последовательных волн пандемии при калибровке по данным рутинного
эпидемиологического надзора.

\textbf{Научная новизна.} Впервые для российского малого города проведён
систематический поволновой анализ пандемии COVID-19 с применением
формальных критериев выбора между целочисленной и дробной
SEIRD-моделями, охватывающий полный период 2020--2025 гг. и все смены
доминирующих вариантов вируса. Впервые α\_k предложен и верифицирован
как косвенный количественный маркер нарастания популяционного
иммунитета, поддающийся оценке из данных рутинного надзора без
иммунологических обследований.

\subsection{2. Материалы и
методы}\label{ux43cux430ux442ux435ux440ux438ux430ux43bux44b-ux438-ux43cux435ux442ux43eux434ux44b}

{[}\hyperref[ux43eux431ux449ux430ux44f-ux441ux445ux435ux43cux430-ux438ux441ux441ux43bux435ux434ux43eux432ux430ux43dux438ux44f-ux438-ux445ux430ux440ux430ux43aux442ux435ux440ux438ux441ux442ux438ux43aux430-ux43eux431ux44aux435ux43aux442ux430-ux43dux430ux431ux43bux44eux434ux435ux43dux438ux44f]{2.1
Общая схема исследования и характеристика объекта наблюдения}{]}

{[}{[}2.2 Предварительная обработка данных{]}{]}

{[}\hyperref[ux438ux434ux435ux43dux442ux438ux444ux438ux43aux430ux446ux438ux44f-ux432ux43eux43bux43d-ux437ux430ux431ux43eux43bux435ux432ux430ux435ux43cux43eux441ux442ux438]{2.3
Идентификация волн заболеваемости}{]}

{[}{[}2.4 Календарь доминирования штаммов{]}{]}

{[}{[}2.5 Стандартная SEIRD модель{]}{]}

{[}{[}2.6 SEIRD модель с производными дробного порядка{]}{]}

{[}\hyperref[ux441ux442ux440ux430ux442ux435ux433ux438ux44f-ux447ux438ux441ux43bux435ux43dux43dux43eux439-ux43eux43fux442ux438ux43cux438ux437ux430ux446ux438ux438]{2.7
Стратегия численной оптимизации}{]}

{[}{[}2.8 Сравнение результатов{]}{]}

\subsubsection{2.1 Общая схема исследования и характеристика объекта
наблюдения}\label{ux43eux431ux449ux430ux44f-ux441ux445ux435ux43cux430-ux438ux441ux441ux43bux435ux434ux43eux432ux430ux43dux438ux44f-ux438-ux445ux430ux440ux430ux43aux442ux435ux440ux438ux441ux442ux438ux43aux430-ux43eux431ux44aux435ux43aux442ux430-ux43dux430ux431ux43bux44eux434ux435ux43dux438ux44f}

\paragraph{2.1.1 Дизайн
исследования}\label{ux434ux438ux437ux430ux439ux43d-ux438ux441ux441ux43bux435ux434ux43eux432ux430ux43dux438ux44f}

Настоящая работа выполнена в формате \textbf{ретроспективного
описательного} \textbf{исследования с элементами математического
моделирования}. Временной горизонт охватывает \textbf{период с 1 января
2020 года по 31 декабря} \textbf{2025 года} --- от регистрации первых
случаев COVID-19 в России до завершения активной фазы пандемии на
субнациональном уровне. Единицей наблюдения является один день; единицей
анализа --- отдельная эпидемическая волна, идентифицированная
алгоритмически из непрерывного ряда ежедневной заболеваемости.

Исследование реализовано в виде \textbf{последовательного
аналитического} \textbf{конвейера}, включающего четыре функционально
обособленных этапа (рисунок 1). На \textbf{первом этапе} исходные данные
государственного эпидемиологического надзора приводятся к единому
формату, проходят контроль качества и восстановление пропусков. На
\textbf{втором этапе} непрерывный ежедневный ряд разбивается на волны
заболеваемости, каждая из которых сопоставляется с доминирующим
вариантом SARS-CoV-2. На \textbf{третьем этапе} для каждой волны
калибруются две конкурирующие математические модели: стандартная SEIRD с
производными целого порядка и дробная SEIRD с производной Капуто порядка
α\_k. На \textbf{четвёртом этапе} модели формально сравниваются с
применением информационных критериев и теста отношения правдоподобия, а
подобранные значения α\_k интерпретируются в иммунологическом контексте
соответствующих волн.

Конвейер реализован \textbf{на языке Julia} 1.x {[}13{]} с
использованием пакетов DataFrames.jl, CSV.jl, DifferentialEquations.jl
{[}12{]} и Optim.jl. Воспроизводимость вычислительной среды
обеспечивается фиксацией версий зависимостей в файле Manifest.toml.
Исходный код, конфигурационные файлы и агрегированные результаты
размещены в открытом репозитории (ссылка будет добавлена после слепого
рецензирования).

\begin{center}\rule{0.5\linewidth}{0.5pt}\end{center}

\paragraph{2.1.2 Район
исследования}\label{ux440ux430ux439ux43eux43d-ux438ux441ux441ux43bux435ux434ux43eux432ux430ux43dux438ux44f}

Объектом наблюдения служит население г. Пскова --- административного
центра Псковской области, расположенного на северо-западе Российской
Федерации (координаты: 57°49′ с.ш., 28°20′ в.д.). По данным переписи
2021 года численность постоянного населения города составляет 202 879
человек, что относит его к категории малых и средних городов России
(100--250 тыс. жителей). Медианный возраст населения --- 42,3 года; доля
лиц старше 60 лет --- 26,1\%, что несколько выше среднероссийского
показателя и обусловлено характерной для городов данного типа возрастной
структурой с отрицательным миграционным сальдо молодёжи.

С точки зрения типичности Псков представляет собой репрезентативный
пример малого регионального центра северо-западного пояса России:
умеренный климат с выраженной зимней сезонностью, средняя плотность
населения (\textasciitilde1 800 чел./км² в городской черте),
ограниченный промышленный сектор с преобладанием занятости в бюджетной
сфере, торговле и транспорте. Близость к государственной границе
Российской Федерации (\textasciitilde50 км до границы с Эстонией) и
исторически высокая доля транзитного трафика обусловливают повышенную по
сравнению со средними городами России интенсивность внешних миграционных
потоков, что могло влиять на сроки завоза новых вариантов вируса.

В моделировании численность восприимчивой популяции принята равной
N\_city = 200 000 чел. (округлённое значение, соответствующее данным
переписи) и используется в качестве верхней границы для оптимизируемого
параметра эффективной популяции N\_eff.

\begin{center}\rule{0.5\linewidth}{0.5pt}\end{center}

\paragraph{2.1.3 Источники
данных}\label{ux438ux441ux442ux43eux447ux43dux438ux43aux438-ux434ux430ux43dux43dux44bux445}

Информационную основу исследования составляют два массива данных
государственного эпидемиологического надзора, предоставленных
Управлением Роспотребнадзора по Псковской области.

\textbf{Массив 1. Еженедельный мониторинг внебольничной пневмонии (ВП).}
Форма федерального государственного статистического наблюдения,
содержащая еженедельные данные о числе зарегистрированных случаев
внебольничной пневмонии в разбивке по субъектам РФ (таблица «ВП»).
Период: 2024--2025 гг. Структура файла: многоуровневые слитые заголовки
(4--5 строк), даты закодированы в именах листов, а не в ячейках данных.

\textbf{Массив 2. Оперативные сводки COVID-19.} Ежедневные данные о
заболеваемости и летальности COVID-19 в разбивке по субъектам РФ.
Период: 2020--2025 гг. Файлы содержат две структурно различных части: до
декабря 2020 года данные представлены в поперечном (широком) формате без
явного столбца даты; начиная с декабря 2020 года --- в длинном формате с
явным столбцом ДАТА. Ключевые переменные: ежедневное число новых
случаев, накопленное число случаев, ежедневное и накопленное число
летальных исходов.

Оба массива изначально находились в формате Microsoft Excel (.xlsx) с
характерными для форм государственной отчётности ограничениями: слитыми
ячейками, пропусками значений вне левой верхней ячейки объединённого
диапазона, русскоязычными наименованиями столбцов и непоследовательной
схемой заголовков между периодами.

Дополнительным источником послужил авторский календарь доминирования
вариантов SARS-CoV-2, составленный на основе публикаций систем геномного
надзора и метааналитических обзоров параметров штаммов {[}4, 5{]}
применительно к российскому эпидемическому контексту. Калибровочный файл
(таблица \texttt{covid19\_seird\_params.csv}) содержит для каждого
варианта: идентификатор штамма, даты начала и окончания периода
доминирования (dom\_start\_adj, dom\_end\_adj), расчётные значения σ, γ,
μ, β, R₀ и производных параметров (T\_incub, T\_infect, IFR, CFR).

\begin{center}\rule{0.5\linewidth}{0.5pt}\end{center}

\paragraph{2.1.4 Этические аспекты и ограничения источников
данных}\label{ux44dux442ux438ux447ux435ux441ux43aux438ux435-ux430ux441ux43fux435ux43aux442ux44b-ux438-ux43eux433ux440ux430ux43dux438ux447ux435ux43dux438ux44f-ux438ux441ux442ux43eux447ux43dux438ux43aux43eux432-ux434ux430ux43dux43dux44bux445}

Исследование основано исключительно на агрегированных
деперсонализированных данных государственной статистической отчётности и
не предполагает работы с персональными данными пациентов. В соответствии
с действующим законодательством Российской Федерации прохождение
этической экспертизы для данного типа исследований не требуется.

Принципиальным ограничением источников данных является системное
занижение регистрируемой заболеваемости, характерное для данных
рутинного надзора: регистрируются только случаи, подтверждённые
лабораторно или клинически при обращении за медицинской помощью. По
имеющимся оценкам для России {[}17{]}, реальная заболеваемость превышала
регистрируемую в 5--10 раз в различные периоды пандемии, причём
коэффициент занижения существенно варьировал между волнами в зависимости
от доступности тестирования и тяжести течения доминирующего варианта.
Следует отметить, что занижение связано не с особенностями организации
учета в Российской Федерации, а системное различие между количеством
заболевших и числом манифестированных случаев, что отражает и социальный
фактор (обращаемость населения в связи с наличием подозрительных
симптомов, доступности диагностических тестов и т.п.).

Это обстоятельство непосредственно влияет на интерпретацию параметра
N\_eff: подобранные значения отражают не всю восприимчивую популяцию
города, а лишь ту её долю, которая была вовлечена в регистрируемую часть
эпидемического процесса. Данное ограничение учтено при обсуждении
результатов (раздел 4).

\begin{center}\rule{0.5\linewidth}{0.5pt}\end{center}

\textbf{Дополнительный источник (к разделу 2.1)}

\section{2.2 Предобработка
данных}\label{ux43fux440ux435ux434ux43eux431ux440ux430ux431ux43eux442ux43aux430-ux434ux430ux43dux43dux44bux445}

\subsection{2.2.1 Первичная подготовка исходных
файлов}\label{ux43fux435ux440ux432ux438ux447ux43dux430ux44f-ux43fux43eux434ux433ux43eux442ux43eux432ux43aux430-ux438ux441ux445ux43eux434ux43dux44bux445-ux444ux430ux439ux43bux43eux432}

Исходные данные поступали в формате Microsoft Excel с многоуровневыми
объединёнными заголовками (4--5 строк), кириллическими наименованиями
столбцов и датами, закодированными исключительно в именах листов. Для
листов до декабря 2020 года использовался широкий кросс-секционный
формат (одна строка на субъект), тогда как начиная с декабря 2020 года
применялся длинный формат с явным столбцом ДАТА. Каждый из трёх
источников --- форма еженедельного мониторинга внебольничных пневмоний
(ВП), оперативные сводки COVID-19 за 2020--2021 гг. и расширенные сводки
за 2021--2025 гг. --- требовал отдельной логики разбора.

На этапе первичной подготовки выполнялось два класса преобразований.
Первый класс --- структурный: унификация схем столбцов, разрешение
объединённых ячеек, извлечение дат из имён листов и добавление их в
качестве явного поля. Второй класс --- терминологический: приведение
кириллических наименований к единому английскому пространству имён
согласно сводной таблице сопоставления (\texttt{column\_mapping}),
содержащей пять листов и охватывающей все форматы источников.
Сопоставление выполнялось вручную во избежание систематических ошибок
автоматического перевода специализированной медицинской терминологии.
Результатом этапа служил набор структурированных таблиц в формате CSV,
пригодных для прямого импорта в Julia посредством \texttt{CSV.jl} и
\texttt{DataFrames.jl} {[}13{]}.

\subsection{2.2.2 Фильтрация по объекту
исследования}\label{ux444ux438ux43bux44cux442ux440ux430ux446ux438ux44f-ux43fux43e-ux43eux431ux44aux435ux43aux442ux443-ux438ux441ux441ux43bux435ux434ux43eux432ux430ux43dux438ux44f}

Из объединённого общероссийского набора данных, охватывающего 85
субъектов Российской Федерации, был выделен подмассив, относящийся к
Псковской области, по составному ключу \texttt{(дата,\ субъект)}. После
фильтрации итоговый временно́й ряд охватывал период с 17 марта 2020 года
по 30 декабря 2024 года, что в совокупности составило 1750 наблюдений.
Базовыми переменными анализа служили: накопленное число подтверждённых
случаев COVID-19 (\texttt{total}), ежедневный прирост (\texttt{daily}),
накопленное число летальных исходов (\texttt{deaths\_total}) и
ежедневное число смертей (\texttt{deaths\_daily}).

\subsection{2.2.3 Восстановление непрерывного временно́го
ряда}\label{ux432ux43eux441ux441ux442ux430ux43dux43eux432ux43bux435ux43dux438ux435-ux43dux435ux43fux440ux435ux440ux44bux432ux43dux43eux433ux43e-ux432ux440ux435ux43cux435ux43dux43dux43eux433ux43e-ux440ux44fux434ux430}

Ранние сводки (март--ноябрь 2020 года) публиковались с нерегулярной,
преимущественно еженедельной периодичностью, что порождало
систематические пропуски в ряду ежедневных приростов. Оперативные данные
за последующие периоды, хотя и публиковались ежедневно, содержали
спорадические пропуски вследствие задержек отчётности.

Восстановление непрерывного ряда выполнялось в два этапа. На первом
этапе к накопленному числу случаев (\texttt{total}) применялась
интерполяция натуральными кубическими сплайнами по индексам наблюдений с
известными значениями, что обеспечивало гладкость второго порядка в
узлах интерполяции. Коэффициенты кубического полинома в каждом интервале
вычислялись из системы уравнений, задающей условия непрерывности
функции, первой и второй производных, при естественных граничных
условиях (нулевая вторая производная на концах). На втором этапе
восстановленный ряд \texttt{total\_interp} принудительно переводился в
неубывающий: при обнаружении убывающего фрагмента --- возможного
артефакта интерполяции --- значение заменялось предшествующим
максимумом. Это условие физически обосновано: накопленное число случаев
является монотонно неубывающей функцией времени.

Ежедневный поток новых случаев восстанавливался как первая конечная
разность интерполированного накопленного ряда: \[
\Delta t​= total_{interp}(t)−total_{interp}(t−1),\Delta_1​=0.
\] Полученная переменная (\texttt{daily\_interp}) служила основной
наблюдаемой при дальнейшем сглаживании и моделировании.

\subsection{2.2.4
Сглаживание}\label{ux441ux433ux43bux430ux436ux438ux432ux430ux43dux438ux435}

Ряд \texttt{daily\_interp} подвергался двум независимым процедурам
сглаживания, результаты которых использовались для различных
аналитических задач. Первая процедура --- скользящее среднее с
каузальным окном 18 дней (\texttt{daily\_interp\_smooth}) ---
применялась для предварительной визуализации и построения
стратифицированных по волнам панелей. Выбор ширины окна продиктован
компромиссом между подавлением недельной периодичности отчётности
(7-дневный цикл) и сохранением межволновой структуры с характерным
масштабом порядка нескольких месяцев.

Вторая процедура --- свёртка с нормированным гауссовым ядром
(стандартное отклонение σ = 15 дней, усечение при \textbar t\textbar{}
\textgreater{} 3σ) с коррекцией граничных эффектов --- давала ряд
\texttt{daily\_interp\_gauss}, использовавшийся для автоматической
идентификации волн (раздел 2.3) и оценки положения пиков. Граничная
коррекция реализовывалась перенормировкой весов на неполном окне, что
исключало систематический сдвиг оценок в начале и конце ряда. Параметр σ
= 15 дней обеспечивал подавление флуктуаций с периодом менее двух недель
при сохранении временно́го разрешения, достаточного для различения волн
длительностью 6--12 недель.

\section{2.3 Идентификация волн
заболеваемости}\label{ux438ux434ux435ux43dux442ux438ux444ux438ux43aux430ux446ux438ux44f-ux432ux43eux43bux43d-ux437ux430ux431ux43eux43bux435ux432ux430ux435ux43cux43eux441ux442ux438}

\subsection{2.3.1 Оценка сглаженной первой
производной}\label{ux43eux446ux435ux43dux43aux430-ux441ux433ux43bux430ux436ux435ux43dux43dux43eux439-ux43fux435ux440ux432ux43eux439-ux43fux440ux43eux438ux437ux432ux43eux434ux43dux43eux439}

Границы волн определялись по смене знака первой производной ряда
\texttt{daily\_interp\_gauss}. Непосредственное дифференцирование
сглаженного ряда нецелесообразно ввиду накопления ошибок конечных
разностей на коротких временных масштабах; поэтому применялась
двухшаговая процедура. На первом шаге вычислялась числовая производная
по схеме центральных разностей второго порядка точности:

\[\dot{y}_t = \frac{y_{t+1} - y_{t-1}}{2}, \quad t = 2, \ldots, n-1,\]

с односторонними разностями первого порядка на граничных точках. На
втором шаге полученный ряд \(\dot{y}\) повторно сглаживался свёрткой с
нормированным гауссовым ядром (σ = 15 дней, усечение при
\textbar τ\textbar{} \textgreater{} 3σ) с той же граничной коррекцией,
что и в разделе 2.2.4. Результирующий ряд
\texttt{daily\_interp\_deriv\_smooth} представляет собой оценку
локальной скорости изменения заболеваемости, устойчивую к
высокочастотным флуктуациям отчётности. Знак производной фиксировался
как +1 (рост), −1 (спад) или 0 (нулевой градиент); нулевые значения
заменялись знаком ближайшего предшествующего ненулевого наблюдения, а
начальный префикс из нулей --- знаком первого ненулевого элемента.

\subsection{2.3.2 Автоматическая разметка
волн}\label{ux430ux432ux442ux43eux43cux430ux442ux438ux447ux435ux441ux43aux430ux44f-ux440ux430ux437ux43cux435ux442ux43aux430-ux432ux43eux43bux43d}

Последовательная разметка волн осуществлялась детерминированным
автоматом на основе ряда знаков производной. Автомат поддерживал два
состояния: \texttt{:up} (фаза роста) и \texttt{:down} (фаза спада).
Переход \texttt{:idle\ →\ :up} при обнаружении первого знака «+1»
инициировал новую волну с инкрементом счётчика; последующая смена знака
на «−1» переводила автомат в состояние \texttt{:down} без изменения
номера текущей волны. Таким образом, каждая волна охватывала непрерывный
интервал, включающий фазу роста и примыкающую фазу спада вплоть до
следующего перехода в фазу роста. Каждому наблюдению присваивался
целочисленный номер волны, начиная с 1; наблюдения до первого
обнаруженного роста помечались нулём и исключались из дальнейшего
волнового анализа.

Описанный подход эквивалентен разбиению временно́го ряда на монотонные
отрезки с последующим объединением каждой пары «подъём + спуск» в единую
волну. Его принципиальное преимущество перед пороговыми методами
(например, выделение периодов, в которых заболеваемость превышает
фиксированную долю максимума) состоит в отсутствии зависимости от
абсолютного уровня заболеваемости, что особенно важно при сравнении волн
с существенно различающейся амплитудой.

\subsection{2.3.3 Характеристики
волн}\label{ux445ux430ux440ux430ux43aux442ux435ux440ux438ux441ux442ux438ux43aux438-ux432ux43eux43bux43d}

Для каждой идентифицированной волны вычислялся набор количественных
характеристик на основе ряда \texttt{daily\_interp\_gauss}. Положение
пика определялось как индекс глобального максимума внутри волнового
интервала. Временны́е границы волны уточнялись алгоритмом водораздела
(watershed): каждому индексу временного ряда ставился в соответствие
ближайший пик по метрике, учитывающей как временно́е расстояние, так и
высоту пика,

\[j^*(i) = \arg\min_{p \in \mathcal{P}} \frac{|i - p|}{y_p},\]

где \(\mathcal{P}\) --- множество индексов пиков, \(y_p\) --- значение
ряда в пике \(p\). Границы волны определялись как крайние индексы,
принадлежащие данному водоразделу и превышающие базовый уровень
\(y_{\min} + 0{,}002 \cdot (y_{\max} - y_{\min})\).

На основании уточнённых границ \([t_{\text{start}}, t_{\text{stop}}]\) и
положения пика \(t_{\text{peak}}\) рассчитывались следующие
характеристики:

--- \textbf{Продолжительность} (duration\_days):
\(t_{\text{stop}} - t_{\text{start}} + 1\) дней.

--- \textbf{Площадь под кривой выше базового уровня} (auc\_above\_base):
\(\sum_{t=t_{\text{start}}}^{t_{\text{stop}}} \bigl(y_t - y_{\text{base}}\bigr)\),
пропорциональна суммарной нагрузке волны.

--- \textbf{Полная ширина на полувысоте} (FWHM): расстояние между
последним индексом слева от пика, где \(y_t \leq y_{\text{peak}}/2\), и
первым таким индексом справа; характеризует концентрацию нагрузки во
времени независимо от общей длительности волны.

--- \textbf{Асимметрия} (asymmetry): отношение
\((t_{\text{peak}} - t_{\text{start}}) / \max(t_{\text{stop}} - t_{\text{peak}}, 1)\);
значение больше единицы указывает на более длинный правый хвост
(медленный спад), меньше единицы --- на более длинный левый хвост
(медленный подъём).

--- \textbf{Выступ пика} (prominence): превышение высоты пика над
уровнем, до которого необходимо опустить ряд, чтобы пик слился с
соседним более высоким пиком или краем ряда; вычислялся по стандартному
алгоритму \texttt{peakproms} из пакета \texttt{Peaks.jl}. Параметр
использовался как фильтр при первоначальном отборе пиков: минимальный
допустимый выступ составлял 0,5 \% размаха ряда, минимальное расстояние
между соседними пиками --- 11 дней.

Совокупность перечисленных характеристик формировала таблицу волн,
служившую основой для стратификации данных перед калибровкой
SEIRD-моделей (раздел 2.5--2.6) и для сравнительного анализа
эпидемических параметров между волнами (раздел 3).

\section{2.4 Календарь доминирования вариантов
SARS-CoV-2}\label{ux43aux430ux43bux435ux43dux434ux430ux440ux44c-ux434ux43eux43cux438ux43dux438ux440ux43eux432ux430ux43dux438ux44f-ux432ux430ux440ux438ux430ux43dux442ux43eux432-sars-cov-2}

\subsection{2.4.1 Составление базы параметров
штаммов}\label{ux441ux43eux441ux442ux430ux432ux43bux435ux43dux438ux435-ux431ux430ux437ux44b-ux43fux430ux440ux430ux43cux435ux442ux440ux43eux432-ux448ux442ux430ux43cux43cux43eux432}

Для привязки биологических характеристик возбудителя к каждому дню
наблюдения был составлен авторский календарь доминирования вариантов
SARS-CoV-2, охватывающий период с марта 2020 по декабрь 2024 года.
Источниками параметров служили систематические обзоры и метаанализы по
инкубационному периоду, серийному интервалу и времени генерации {[}4,
5{]}, а также первичные публикации систем геномного надзора,
документирующие смену доминирующих линий в России и сопредельных
странах.

Для каждого варианта в таблице \texttt{covid19\_seird\_params.csv}
зафиксированы следующие поля. Идентификационные: \texttt{strain\_id}
(уникальный строковый ключ), \texttt{strain\_name} (общепринятое
название), \texttt{pango\_lineage} (классификация Pango),
\texttt{dom\_start\_adj} и \texttt{dom\_end\_adj} (скорректированные
даты начала и окончания периода доминирования). Биологические параметры:
\texttt{R0\_min}, \texttt{R0\_avg}, \texttt{R0\_max} --- диапазон
базового репродуктивного числа; \texttt{T\_incub\_avg\_days} --- средняя
длительность латентного периода с соответствующей скоростью
экспонирования \(\sigma = 1/T_{\text{incub}}\) (поле
\texttt{sigma\_avg\_per\_day}); \texttt{T\_infect\_avg\_days} ---
средняя длительность инфекционного периода с соответствующей скоростью
выздоровления \(\gamma = 1/T_{\text{infect}}\) (поле
\texttt{gamma\_avg\_per\_day}); \texttt{IFR\_avg} --- среднее значение
инфекционной летальности с производной скоростью смерти
\(\mu = \text{IFR} \cdot \gamma / (1 - \text{IFR})\) (поле
\texttt{mu\_avg\_per\_day}); \texttt{beta\_min\_per\_day},
\texttt{beta\_avg\_per\_day}, \texttt{beta\_max\_per\_day} --- диапазон
скорости передачи, рассчитанный как
\(\beta = R_0 \cdot (\gamma + \mu)\); \texttt{T\_gen\_avg\_days} ---
среднее время генерации. Качественные признаки: \texttt{severity}
(относительная клиническая тяжесть) и \texttt{immune\_escape}
(способность варианта уклоняться от выработанного иммунного ответа).

Значения параметров для предкового уханьского варианта, Alpha и Delta
заимствованы из метаанализа Xu et al.~{[}4{]}, охватывающего 147
исследований по временны́м параметрам передачи; для Omicron и его
субвариантов (BA.1, BA.2, BA.4/5, XBB, JN.1 и поздних линий)
дополнительно привлечены данные Madewell et al.~{[}5{]}. Для периодов,
не охваченных указанными метаанализами, параметры экстраполировались на
основании тенденций укорочения серийного интервала, задокументированных
при последовательной смене линий Omicron.

\subsection{2.4.2 Развёртка в ежедневный
календарь}\label{ux440ux430ux437ux432ux451ux440ux442ux43aux430-ux432-ux435ux436ux435ux434ux43dux435ux432ux43dux44bux439-ux43aux430ux43bux435ux43dux434ux430ux440ux44c}

Исходная таблица параметров содержала по одной строке на вариант с
интервальными датами доминирования. Для последующего объединения с
ежедневным рядом заболеваемости каждая строка была развёрнута в набор
однодневных записей: для каждой пары
\texttt{(dom\_start\_adj,\ dom\_end\_adj)} генерировался непрерывный
перечень дат с шагом один день, каждой из которых присваивались все
параметры соответствующего варианта. Результирующая промежуточная
таблица \texttt{strain\_daily} содержала по одной строке на каждый день
активности каждого варианта.

\subsection{2.4.3 Разрешение
перекрытий}\label{ux440ux430ux437ux440ux435ux448ux435ux43dux438ux435-ux43fux435ux440ux435ux43aux440ux44bux442ux438ux439}

Временны́е интервалы доминирования смежных вариантов в ряде случаев
перекрываются: новый вариант достигает порога доминирования прежде, чем
предыдущий полностью вытесняется из циркуляции. При объединении по дате
это порождает неоднозначность --- несколько строк на одну дату. Для её
устранения применялось детерминированное правило приоритета: при наличии
нескольких вариантов, активных в одну дату, сохранялась запись с
наиболее поздним значением \texttt{dom\_start\_adj}, то есть вариант с
более поздним началом доминирования считался актуальным. При равенстве
этой даты использовался порядковый номер варианта \texttt{strain\_num}
как вторичный ключ сортировки. Каждой записи присваивалась метка типа
сопоставления: \texttt{"exact"} при отсутствии перекрытий или
\texttt{"overlap\_resolved"} при их наличии, а также счётчик
\texttt{n\_matching\_strains}. Результирующая таблица
\texttt{strain\_daily\_resolved} содержала ровно одну строку на дату,
покрывая весь период наблюдений.

\subsection{2.4.4 Присоединение к ряду
заболеваемости}\label{ux43fux440ux438ux441ux43eux435ux434ux438ux43dux435ux43dux438ux435-ux43a-ux440ux44fux434ux443-ux437ux430ux431ux43eux43bux435ux432ux430ux435ux43cux43eux441ux442ux438}

Таблица \texttt{strain\_daily\_resolved} была присоединена к основному
ежедневному ряду заболеваемости операцией левого соединения
(\texttt{leftjoin}) по ключу \texttt{date}. Результирующий обогащённый
DataFrame \texttt{df\_enriched} содержал для каждого дня наблюдения как
эпидемиологические переменные (\texttt{daily\_interp},
\texttt{daily\_interp\_smooth}, номер волны и производные), так и
биологические параметры доминирующего варианта (\(\sigma\), \(\gamma\),
\(\mu\), \(\beta\), \(R_0\)). Дни, выходящие за пределы охвата календаря
штаммов, получали \texttt{missing} по всем полям варианта и при
калибровке моделей не использовались.

Таким образом, для каждой волны, идентифицированной в разделе 2.3,
оказывались известны априорные значения скоростей \(\sigma\) и
\(\gamma\), служившие фиксированными параметрами при калибровке
SEIRD-моделей (разделы 2.5--2.6), а значения \(\beta_{\text{avg}}\) и
\(\mu_{\text{avg}}\) --- справочными стартовыми точками оптимизации.

\section{2.5 Стандартная
SEIRD-модель}\label{ux441ux442ux430ux43dux434ux430ux440ux442ux43dux430ux44f-seird-ux43cux43eux434ux435ux43bux44c}

\subsection{2.5.1 Структура
модели}\label{ux441ux442ux440ux443ux43aux442ux443ux440ux430-ux43cux43eux434ux435ux43bux438}

Динамика эпидемической волны описывалась детерминированной
компартментной моделью SEIRD с пятью состояниями: восприимчивые (\(S\)),
экспонированные (\(E\), инфицированные, но ещё не заразные),
инфекционные (\(I\)), выздоровевшие (\(R\)) и летальные (\(D\)).
Популяция предполагается замкнутой с эффективным размером
\(N_{\text{eff}}\), так что в любой момент времени выполняется условие
сохранения:

\[S(t) + E(t) + I(t) + R(t) + D(t) = N_{\text{eff}}.\]

Система ОДУ имеет вид:

\[\frac{dS}{dt} = -\frac{\beta, I}{N_{\text{eff}}}, S,\]

\[\frac{dE}{dt} = \frac{\beta, I}{N_{\text{eff}}}, S - \sigma, E,\]

\[\frac{dI}{dt} = \sigma, E - (\gamma + \mu), I,\]

\[\frac{dR}{dt} = \gamma, I,\]

\[\frac{dD}{dt} = \mu, I,\]

где \(\beta\) --- скорость передачи инфекции (сут\(^{-1}\)),
\(\sigma = 1/T_{\text{incub}}\) --- скорость экспонирования
(сут\(^{-1}\)), \(\gamma = 1/T_{\text{infect}}\) --- скорость
выздоровления (сут\(^{-1}\)), \(\mu\) --- скорость летальных исходов
(сут\(^{-1}\)). Базовое репродуктивное число для данной конфигурации
параметров выражается как:

\[R_0 = \frac{\beta}{\gamma + \mu}.\]

Выбор структуры SEIRD обусловлен необходимостью явно моделировать как
латентную фазу инфекции --- принципиально важную для SARS-CoV-2 с его
характерным инкубационным периодом от 2 до 6 дней {[}4{]} --- так и
летальные исходы в качестве второй независимой наблюдаемой, что
существенно снижает сложность идентификации параметров, описанную в
работе {[}3{]}.

\subsection{2.5.2 Параметры модели и их
разделение}\label{ux43fux430ux440ux430ux43cux435ux442ux440ux44b-ux43cux43eux434ux435ux43bux438-ux438-ux438ux445-ux440ux430ux437ux434ux435ux43bux435ux43dux438ux435}

Параметры модели разделялись на два класса: фиксированные, задаваемые из
внешнего источника, и свободные, подбираемые в ходе оптимизации.

К \textbf{фиксированным} относились скорости \(\sigma\) и \(\gamma\),
берущиеся непосредственно из поля \texttt{sigma\_avg\_per\_day} и
\texttt{gamma\_avg\_per\_day} таблицы \texttt{df\_enriched} для
преобладающего штамма данной волны (раздел 2.4). Это решение согласуется
с рекомендацией Королёва {[}3{]} о целесообразности фиксации
биологически обоснованных параметров для устранения вырожденности
оптимизационной задачи.

\textbf{Свободными} параметрами служили: \(\beta\) --- скорость
передачи, \(\mu\) --- скорость летальных исходов, \(E_0 = E(t_0)\) и
\(I_0 = I(t_0)\) --- начальные размеры экспонированного и инфекционного
компартментов, а также \(N_{\text{eff}}\) --- эффективный размер
восприимчивой популяции. Таким образом, вектор свободных параметров
содержал пять компонент:

\[\boldsymbol{\theta} = (\beta,; \mu,; E_0,; I_0,; N_{\text{eff}}).\]

Параметр \(N_{\text{eff}}\) интерпретируется как размер той доли
популяции, которая фактически участвовала в передаче инфекции в данной
волне, и ограничивался сверху численностью города:
\(N_{\text{eff}} \leq N_{\text{city}} = 200,000\). Занижение
регистрируемых случаев (раздел 2.1.4) означает, что подобранное
\(N_{\text{eff}}\) отражает не всю биологически восприимчивую популяцию,
а лишь ту её часть, которая генерировала наблюдаемый поток
зарегистрированных случаев. Этот эффект систематически обсуждается в
разделе 4.

Начальные условия для \(S\), \(R\) и \(D\) определялись из условия
сохранения и нулевых начальных значений:
\(S(t_0) = N_{\text{eff}} - E_0 - I_0\), \(R(t_0) = 0\), \(D(t_0) = 0\).

\subsection{2.5.3 Наблюдаемые и функция
потерь}\label{ux43dux430ux431ux43bux44eux434ux430ux435ux43cux44bux435-ux438-ux444ux443ux43dux43aux446ux438ux44f-ux43fux43eux442ux435ux440ux44c}

Модель сопоставлялась с двумя независимыми наблюдаемыми рядами. Первый
--- ежедневный поток новых клинически значимых случаев ---
аппроксимировался скоростью выхода из компартмента \(E\):

\[\hat{c}(t) = \sigma, E(t).\]

Второй --- ежедневное число летальных исходов --- аппроксимировалось
скоростью выхода из компартмента \(I\) по летальному пути:

\[\hat{d}(t) = \mu, I(t).\]

Наблюдаемыми служили восстановленные ряды \texttt{daily\_interp} и
\texttt{deaths\_daily} для соответствующей волны. В случаях, когда
данные о смертях для отдельных суток отсутствовали, пропуски заменялись
нулём. Функция потерь определялась как логарифмический
среднеквадратический error по случаям:

\[\mathcal{L}(\boldsymbol{\theta}) = \frac{1}{T}\sum_{t=1}^{T} \bigl(\ln(1 + \hat{c}(t;\boldsymbol{\theta})) - \ln(1 + c(t))\bigr)^2,\]

где \(c(t)\) --- наблюдаемый ежедневный прирост, \(T\) --- длина волны в
сутках. Логарифмическое преобразование обосновано тем, что ошибки на
пике волны и в её хвосте имеют принципиально различный абсолютный
масштаб: минимизация обычного MSE концентрирует оптимизацию на
нескольких точках вблизи пика, тогда как log-MSE обеспечивает
сбалансированный вклад восходящей фазы, пика и нисходящего хвоста ---
трёх структурно различных частей волны, несущих различную информацию о
параметрах модели.

\subsection{2.5.4 Численное
решение}\label{ux447ux438ux441ux43bux435ux43dux43dux43eux435-ux440ux435ux448ux435ux43dux438ux435}

Система ОДУ решалась численно с использованием пакета
DifferentialEquations.jl {[}12{]} для языка Julia {[}13{]}. В качестве
решателя использовался адаптивный метод \texttt{AutoVern7(Rodas5P)},
автоматически переключающийся между нежёстким высокоточным интегратором
Верне седьмого порядка (Vern7) и жёстким решателем Розенброка пятого
порядка (Rodas5P) в зависимости от локальной жёсткости системы. Шаг
сохранения результата составлял \(h_{\text{save}} = 1\) день; абсолютный
и относительный допуски интегрирования --- \(10^{-8}\) и \(10^{-6}\)
соответственно. При каждом обращении решателя в ходе оптимизации все
пять свободных параметров передавались через вектор \texttt{p}, а
состояние системы в момент \(t = 0\) формировалось из \(E_0\), \(I_0\) и
\(N_{\text{eff}}\) аналитически.

Для обеспечения вычислительной устойчивости в ходе оптимизации все
свободные параметры параметризовались в логарифмическом пространстве:
\(\tilde{\theta}_i = \ln \theta_i\), что автоматически гарантировало их
положительность и существенно улучшало обусловленность задачи при
больших вариациях масштабов параметров. Стратегия численной оптимизации
описана в разделе 2.7.

\section{2.6 SEIRD-модель с производными дробного
порядка}\label{seird-ux43cux43eux434ux435ux43bux44c-ux441-ux43fux440ux43eux438ux437ux432ux43eux434ux43dux44bux43cux438-ux434ux440ux43eux431ux43dux43eux433ux43e-ux43fux43eux440ux44fux434ux43aux430}

\subsection{2.6.1 Производная Капуто и её эпидемиологическая
интерпретация}\label{ux43fux440ux43eux438ux437ux432ux43eux434ux43dux430ux44f-ux43aux430ux43fux443ux442ux43e-ux438-ux435ux451-ux44dux43fux438ux434ux435ux43cux438ux43eux43bux43eux433ux438ux447ux435ux441ux43aux430ux44f-ux438ux43dux442ux435ux440ux43fux440ux435ux442ux430ux446ux438ux44f}

Дробная SEIRD-модель формулируется посредством замены производных целого
порядка в системе раздела 2.5 на производные Капуто порядка
\(\alpha_k \in (0{,}5;, 1{,}0]\). Производная Капуто порядка \(\alpha\)
от функции \(f\) в момент \(t\) определяется как {[}6, 7{]}:

\[{}^{!C}D^\alpha f(t) = \frac{1}{\Gamma(1-\alpha)} \int_0^t \frac{f'(\tau)}{(t - \tau)^{\alpha}}, d\tau, \quad \alpha \in (0, 1),\]

где \(\Gamma(\cdot)\) --- гамма-функция, а интеграл берётся по всей
предыстории процесса от начального момента до \(t\). В отличие от
производной Римана--Лиувилля, производная Капуто допускает постановку
начальных условий в классической форме --- значениями функции \(f(0)\),
--- что делает её естественным выбором для задач Коши в динамических
системах {[}7{]}.

Ключевое свойство оператора, существенное для данного применения,
состоит в следующем: при \(\alpha < 1\) ядро памяти
\((t - \tau)^{-\alpha}\) убывает по степенно́му закону, а не
экспоненциально. Это означает, что состояние системы в момент \(t\) явно
зависит от всей её истории с начала волны, причём вклад прошлых
состояний убывает медленнее, чем в стандартной ОДУ. Применительно к
эпидемической динамике такое поведение может отражать накопленный
популяционный иммунитет, гетерогенность контактных сетей и поведенческие
изменения, не сбрасываемые при переходе между волнами, --- факторы,
которые создают устойчивую зависимость текущей скорости распространения
от накопленного иммунного опыта. При \(\alpha_k \to 1\) степенно́е ядро
вырождается в дельта-функцию Дирака, и дробная система тождественно
переходит в целочисленную SEIRD раздела 2.5.

\subsection{2.6.2 Структура дробной
системы}\label{ux441ux442ux440ux443ux43aux442ux443ux440ux430-ux434ux440ux43eux431ux43dux43eux439-ux441ux438ux441ux442ux435ux43cux44b}

Дробная SEIRD-модель для волны \(k\) определяется системой:

\[{}^{!C}D^{\alpha_k} S = -\frac{\beta, I}{N_{\text{eff}}}, S,\]

\[{}^{!C}D^{\alpha_k} E = \frac{\beta, I}{N_{\text{eff}}}, S - \sigma, E,\]

\[{}^{!C}D^{\alpha_k} I = \sigma, E - (\gamma + \mu), I,\]

\[{}^{!C}D^{\alpha_k} R = \gamma, I,\]

\[{}^{!C}D^{\alpha_k} D = \mu, I,\]

где параметры \(\sigma\), \(\gamma\) фиксированы из календаря штаммов
аналогично разделу 2.5, а правые части совпадают с целочисленной
версией. Единый порядок \(\alpha_k\) применяется ко всем пяти уравнениям
системы, что обеспечивает сохранение условия замкнутости популяции
{[}8{]}: при любом \(\alpha_k \in (0, 1]\) сумма
\(S + E + I + R + D = N_{\text{eff}}\) остаётся постоянной, поскольку
правые части системы суммируются в нуль тождественно.

\subsection{2.6.3 Дискретизация по схеме
Грюнвальда--Летникова}\label{ux434ux438ux441ux43aux440ux435ux442ux438ux437ux430ux446ux438ux44f-ux43fux43e-ux441ux445ux435ux43cux435-ux433ux440ux44eux43dux432ux430ux43bux44cux434ux430ux43bux435ux442ux43dux438ux43aux43eux432ux430}

Для численного решения системы применялась явная схема
Грюнвальда--Летникова (GL) с шагом \(h = 1\) день. Дискретный аналог
производной Капуто порядка \(\alpha\) в точке \(t_n = nh\) записывается
как {[}6, 8{]}:

\[{}^{!C}D^\alpha_h f(t_n) \approx h^{-\alpha} \sum_{j=0}^{n} w_j^{(\alpha)}, f(t_{n-j}),\]

где GL-веса \(w_j^{(\alpha)}\) определяются рекуррентно:

\[w_0^{(\alpha)} = 1, \qquad w_j^{(\alpha)} = w_{j-1}^{(\alpha)} \cdot \frac{j - 1 - \alpha}{j}, \quad j = 1, 2, \ldots\]

При \(\alpha < 1\) веса \(w_j^{(\alpha)}\) знакочередуются:
\(w_0 = 1 > 0\), \(w_1 = -\alpha < 0\), далее убывают по абсолютной
величине как \(O(j^{-1-\alpha})\). Это выражает степенно́е затухание
памяти: вклад состояния в момент \(t_{n-j}\) убывает как
\(j^{-1-\alpha}\) с ростом лага \(j\), тогда как экспоненциальное ядро
давало бы экспоненциальное затухание.

Схема GL первого порядка точности по \(h\): глобальная ошибка \(O(h)\)
{[}9{]}. При \(h = 1\) день и характерных временны́х масштабах волн
порядка 30--120 дней это обеспечивает приемлемую точность. На каждом
шаге \(n\) вычисляется полная сумма по всем предшествующим значениям
компартментов, что требует хранения всей истории волны; вычислительная
сложность составляет \(O(T^2)\) по числу шагов \(T\). Для волн
длительностью до 150 дней при однодневном шаге это приводит к не более
чем \(\sim 10^4\) операций умножения на шаг, что при пятикомпонентной
системе остаётся вычислительно незначительным.

\subsection{2.6.4 Свободные параметры и
ограничения}\label{ux441ux432ux43eux431ux43eux434ux43dux44bux435-ux43fux430ux440ux430ux43cux435ux442ux440ux44b-ux438-ux43eux433ux440ux430ux43dux438ux447ux435ux43dux438ux44f}

Вектор свободных параметров дробной модели расширен относительно
целочисленной на один элемент --- порядок производной \(\alpha_k\):

\[\boldsymbol{\theta}^{(\text{frac})} = (\beta,; \mu,; E_0,; I_0,; N_{\text{eff}},; \alpha_k).\]

Нижняя граница \(\alpha_k > 0{,}5\) введена из содержательных
соображений: при \(\alpha_k \leq 0{,}5\) память ядра становится
настолько длинной, что ранние фазы эпидемии начинают доминировать над
текущей динамикой, утрачивая биологический смысл применительно к
отдельным волнам. Верхняя граница \(\alpha_k = 1\) обеспечивает
гнездование моделей: целочисленная SEIRD является частным случаем
дробной при \(\alpha_k = 1\), что позволяет применять тест отношения
правдоподобия для формальной проверки значимости дробного расширения
(раздел 2.8).

При оптимизации \(\alpha_k\) параметризовался через сигмоидальное
преобразование, отображающее \(\mathbb{R}\) в интервал
\((0{,}5;, 1{,}0]\):

\[\alpha_k = 0{,}5 + 0{,}5 \cdot \sigma(\tilde{\alpha}), \qquad \sigma(\tilde{\alpha}) = \frac{1}{1 + e^{-\tilde{\alpha}}},\]

где \(\tilde{\alpha} \in \mathbb{R}\) --- неограниченный оптимизируемый
параметр. Остальные пять параметров параметризовались в логарифмическом
пространстве аналогично разделу 2.5.4. Функция потерь и стратегия
оптимизации для дробной модели идентичны описанным в разделах 2.5.3 и
2.7 соответственно.

\section{2.7 Стратегия численной
оптимизации}\label{ux441ux442ux440ux430ux442ux435ux433ux438ux44f-ux447ux438ux441ux43bux435ux43dux43dux43eux439-ux43eux43fux442ux438ux43cux438ux437ux430ux446ux438ux438}

\subsection{2.7.1 Постановка задачи и общий
план}\label{ux43fux43eux441ux442ux430ux43dux43eux432ux43aux430-ux437ux430ux434ux430ux447ux438-ux438-ux43eux431ux449ux438ux439-ux43fux43bux430ux43d}

Калибровка обеих моделей сводилась к минимизации функции потерь
\(\mathcal{L}(\boldsymbol{\theta})\) (раздел 2.5.3) по вектору свободных
параметров \(\boldsymbol{\theta}\) при следующих ограничениях: все
параметры строго положительны;
\(N_{\text{eff}} \leq N_{\text{city}} = 200,000\);
\(N_{\text{eff}} \geq E_0 + I_0 + 1\); \(\alpha_k \in (0{,}5;,1{,}0]\)
(только для дробной модели). Ландшафт \(\mathcal{L}\) невыпуклый:
наблюдается как минимум один глубокий «физически осмысленный» минимум
вблизи корректных биологических параметров, окружённый пологими долинами
с численно приемлемыми, но биологически вырожденными решениями
(например, \(N_{\text{eff}} \to N_{\text{city}}\) при \(\beta \to 0\)).
Для надёжного нахождения глобального минимума применялась
трёхступенчатая процедура: аналитическое предобусловливание стартовой
точки, двумерный сеточный поиск и финальная градиентная полировка.

\subsection{2.7.2 Параметризация и граничные
условия}\label{ux43fux430ux440ux430ux43cux435ux442ux440ux438ux437ux430ux446ux438ux44f-ux438-ux433ux440ux430ux43dux438ux447ux43dux44bux435-ux443ux441ux43bux43eux432ux438ux44f}

Все компоненты вектора
\(\boldsymbol{\theta} = (\beta, \mu, E_0, I_0, N_{\text{eff}})\)
передавались оптимизатору в логарифмическом пространстве:
\(\tilde{\theta}_i = \ln \theta_i\). Это автоматически гарантирует
строгую положительность, устраняет масштабную несоразмерность между
\(\beta \sim 10^{-1}\) и \(N_{\text{eff}} \sim 10^{4}\), а также
приближает ландшафт \(\mathcal{L}(\tilde{\boldsymbol{\theta}})\) к более
симметричной форме в окрестности минимума. Ограничение
\(N_{\text{eff}} \leq N_{\text{city}}\) реализовывалось возвратом
штрафного значения \(\mathcal{L} = 10^{18}\) при его нарушении. Для
дробной модели параметр \(\alpha_k\) вводился через сигмоидальное
преобразование (раздел 2.6.4) и оптимизировался совместно с остальными
компонентами вектора.

\subsection{2.7.3 Аналитический предварительный выбор стартовой
точки}\label{ux430ux43dux430ux43bux438ux442ux438ux447ux435ux441ux43aux438ux439-ux43fux440ux435ux434ux432ux430ux440ux438ux442ux435ux43bux44cux43dux44bux439-ux432ux44bux431ux43eux440-ux441ux442ux430ux440ux442ux43eux432ux43eux439-ux442ux43eux447ux43aux438}

Качество начального приближения принципиально определяет эффективность
последующей оптимизации. Для его получения применялась система
аналитических оценок, основанных на наблюдаемых характеристиках волны.

\textbf{Оценка скорости роста.} Из ранней восходящей фазы волны методом
наименьших квадратов в логарифмической шкале оценивалась
экспоненциальная скорость роста \(r_{\text{est}}\). При линеаризации
SEIRD-системы вблизи начального состояния
\((S \approx N_{\text{eff}},; E, I \ll N_{\text{eff}})\)
характеристический показатель роста определяется соотношением, из
которого непосредственно следует оценка скорости передачи:

\[\hat{\beta} = \frac{(r_{\text{est}} + \sigma)(r_{\text{est}} + \gamma + \mu)}{\sigma}.\]

\textbf{Оценка эффективной популяции.} В рамках классической теории
компартментных моделей {[}1, 2{]} финальный размер эпидемии
\(z = 1 - S(\infty)/N_{\text{eff}}\) связан с \(R_0\) трансцендентным
уравнением. Из наблюдаемого числа случаев до пика волны
\(C_{\text{peak}}\) и оценки \(\hat{R}_0 = \hat{\beta}/(\gamma + \mu)\)
выводилась приближённая оценка:

\[\hat{N}_{\text{eff}} = \frac{C_{\text{peak}} \cdot \hat{R}_0}{\hat{R}_0 - 1},\]

ограниченная снизу удвоенным суммарным числом случаев волны и сверху
\(N_{\text{city}}\).

\textbf{Оценка начальных условий.} Начальный размер инфекционного
компартмента оценивался обратной экстраполяцией от пика:

\[\hat{I}_0 = \max!\bigl(0{,}5,; I_{\text{peak}} \cdot e^{-r_{\text{est}} \cdot t_{\text{peak}}}\bigr),\]

где \(I_{\text{peak}} = c_{\text{max}} / (\gamma + \mu)\),
\(c_{\text{max}}\) --- максимум наблюдаемого ежедневного прироста,
\(t_{\text{peak}}\) --- индекс дня пика. Начальный экспонированный
компартмент определялся из стационарного соотношения SEIR в фазе роста:
\(\hat{E}_0 = \hat{I}_0 \cdot (\gamma + \mu) / \sigma\).

\subsection{2.7.4 Двумерный сеточный
поиск}\label{ux434ux432ux443ux43cux435ux440ux43dux44bux439-ux441ux435ux442ux43eux447ux43dux44bux439-ux43fux43eux438ux441ux43a}

Аналитические оценки \(\hat{\beta}\) и \(\hat{N}_{\text{eff}}\) служили
центрами двумерных логарифмических сеток. Для \(\beta\) строилась сетка
из 20 точек в диапазоне \([\hat{\beta}/4,; 4\hat{\beta}]\); для
\(N_{\text{eff}}\) --- из 20 точек в диапазоне
\([\max(1{,}5 \cdot C_{\text{total}},; 5),; N_{\text{city}}]\), где
\(C_{\text{total}}\) --- суммарное число случаев волны. На каждом из
\(400\) узлов сетки функция потерь вычислялась при
\(\mu = \mu_{\text{prior}}\) и аналитически вычисленных \(E_0\),
\(I_0\), зависящих от текущего \(\beta\) и \(N_{\text{eff}}\). Узел с
минимальным значением \(\mathcal{L}\) давал стартовую точку
\(\boldsymbol{\theta}_0\) для последующей оптимизации. Этот шаг
принципиально важен: он позволяет избежать ухода результатов оптимизации
NelderMead в биологически бессмысленные области ландшафта.

\subsection{2.7.5 Двухэтапная финальная
оптимизация}\label{ux434ux432ux443ux445ux44dux442ux430ux43fux43dux430ux44f-ux444ux438ux43dux430ux43bux44cux43dux430ux44f-ux43eux43fux442ux438ux43cux438ux437ux430ux446ux438ux44f}

\textbf{Этап 1 --- NelderMead с мультистартом.} От стартовой точки
\(\boldsymbol{\theta}_0\) запускалось 15 независимых прогонов метода
Нелдера--Мида {[}13{]}: один прогон из \(\boldsymbol{\theta}_0\) и 14
прогонов из возмущённых стартов
\(\boldsymbol{\theta}_0 + \boldsymbol{\varepsilon}_i\), где компоненты
\(\boldsymbol{\varepsilon}_i\) являлись независимыми гауссовыми
возмущениями с асимметричными масштабами:
\(\sigma_{\tilde{\beta}} = 0{,}3\), \(\sigma_{\tilde{\mu}} = 0{,}8\),
\(\sigma_{\tilde{E}_0} = \sigma_{\tilde{I}_0} = 1{,}0\),
\(\sigma_{\tilde{N}} = 0{,}4\). Бо́льшие масштабы возмущений для \(E_0\)
и \(I_0\) отражают их худшую определённость из наблюдаемых данных
относительно \(\beta\) и \(N_{\text{eff}}\). Каждый прогон выполнялся до
6,000 итераций при допуске \(g_{\text{tol}} = 10^{-8}\); по завершении
всех стартов отбирались три кандидата с наименьшим \(\mathcal{L}\).

Воспроизводимость мультистарта обеспечивалась инициализацией генератора
случайных чисел фиксированным зерном, зависящим от номера волны:
\(\text{seed} = \text{SEED} + k\), где \(\text{SEED}\) --- константа
проекта.

\textbf{Этап 2 --- L-BFGS.} Каждый из трёх кандидатов передавался в
квазиньютоновский метод L-BFGS (до 10,000 итераций,
\(g_{\text{tol}} = 10^{-10}\)) для уточнения в окрестности минимума.
Финальным результатом считался вектор с наименьшим \(\mathcal{L}\) среди
трёх уточнённых решений. Метод L-BFGS обеспечивает суперлинейную
сходимость вблизи гладкого минимума при существенно меньших
вычислительных затратах, чем полный метод Ньютона, благодаря
низкоранговому приближению гессиана {[}13{]}.

\subsection{2.7.6 Взвешивание
наблюдаемых}\label{ux432ux437ux432ux435ux448ux438ux432ux430ux43dux438ux435-ux43dux430ux431ux43bux44eux434ux430ux435ux43cux44bux445}

Функция потерь включала два слагаемых --- по случаям и по смертям. В
волнах, для которых суммарное число летальных исходов превышало 5,
слагаемое по смертям вводилось с весовым коэффициентом \(w_d = 3{,}0\):

\[\mathcal{L}(\boldsymbol{\theta}) = \mathcal{L}_c(\boldsymbol{\theta}) + w_d \cdot \mathcal{L}_d(\boldsymbol{\theta}),\]

где \(\mathcal{L}_c\) и \(\mathcal{L}_d\) --- log-MSE по случаям и
смертям соответственно. В волнах с малым числом смертей (\(\leq 5\))
устанавливалось \(w_d = 0\), поскольку случайная вариация отдельных
регистрируемых смертей привела бы к неустойчивости оценки \(\mu\).
Тройной вес при достаточном числе смертей обусловлен тем, что ежедневная
смертность несёт независимую информацию о параметре \(\mu\), существенно
снижая его коллинеарность с \(\beta\) и \(N_{\text{eff}}\).

\section{2.8 Сравнение
моделей}\label{ux441ux440ux430ux432ux43dux435ux43dux438ux435-ux43cux43eux434ux435ux43bux435ux439}

\subsection{2.8.1 Метрики качества
подгонки}\label{ux43cux435ux442ux440ux438ux43aux438-ux43aux430ux447ux435ux441ux442ux432ux430-ux43fux43eux434ux433ux43eux43dux43aux438}

Для каждой волны \(k\) и каждой из двух моделей вычислялся стандартный
набор метрик качества подгонки по ряду ежедневных новых случаев. Пусть
\(c_t\) --- наблюдаемый прирост, \(\hat{c}_t\) --- модельное
предсказание, \(T\) --- длина волны в сутках.

Среднеквадратическое отклонение:

\[\text{RMSE} = \sqrt{\frac{1}{T}\sum_{t=1}^{T}(c_t - \hat{c}_t)^2}.\]

Средняя абсолютная ошибка:

\[\text{MAE} = \frac{1}{T}\sum_{t=1}^{T}|c_t - \hat{c}_t|.\]

Средняя абсолютная процентная ошибка, вычисляемая в исходном масштабе:

\[\text{MAPE} = \frac{100}{T}\sum_{t=1}^{T}\frac{|c_t - \hat{c}_t|}{\max(c_t, 1)},\]

где знаменатель ограничен снизу единицей во избежание деления на нуль в
периоды нулевой регистрации. Коэффициент детерминации:

\[R^2 = 1 - \frac{\sum_{t}(c_t - \hat{c}_t)^2}{\sum_{t}(c_t - \bar{c})^2},\]

где \(\bar{c}\) --- выборочное среднее наблюдаемого ряда. Наконец,
логарифмический MSE, служащий целевым функционалом при оптимизации
(раздел 2.5.3), вычислялся и как диагностическая метрика для
сопоставления между волнами:

\[\mathcal{L} = \frac{1}{T}\sum_{t=1}^{T}\bigl(\ln(1 + \hat{c}_t) - \ln(1 + c_t)\bigr)^2.\]

Все метрики вычислялись отдельно для целочисленной модели
(\(\mathcal{M}_{\text{int}}\), 5 параметров) и дробной
(\(\mathcal{M}_{\text{frac}}\), 6 параметров).

\subsection{2.8.2 Информационные
критерии}\label{ux438ux43dux444ux43eux440ux43cux430ux446ux438ux43eux43dux43dux44bux435-ux43aux440ux438ux442ux435ux440ux438ux438}

Для сравнения моделей с различным числом параметров применялись
информационные критерии AIC и Байесовский, построенные на основе log-MSE
как суррогата логарифма правдоподобия. При предположении о нормально
распределённых ошибках в логарифмическом масштабе логарифм правдоподобия
пропорционален \(-T \cdot \mathcal{L}/2\), откуда:

\[\text{AIC} = T \cdot \ln(\mathcal{L}) + 2k,\]

\[\text{BIC} = T \cdot \ln(\mathcal{L}) + k \cdot \ln(T),\]

где \(k = 5\) для целочисленной модели и \(k = 6\) для дробной. Разности
\(\Delta\text{AIC} = \text{AIC}_{\text{int}} - \text{AIC}_{\text{frac}}\)
и
\(\Delta\text{BIC} = \text{BIC}_{\text{int}} - \text{BIC}_{\text{frac}}\)
интерпретировались по стандартным шкалам: \(\Delta > 10\) ---
существенное свидетельство в пользу дробной модели;
\(2 < \Delta \leq 10\) --- умеренное; \(|\Delta| \leq 2\) --- модели
практически неразличимы. Критерий BIC штрафует дополнительный параметр
\(\alpha_k\) строже, чем AIC, при \(T > 7\) (что выполняется для всех
волн), что делает его более консервативной мерой оценки улучшения
качества.

\subsection{2.8.3 Тест отношения
правдоподобия}\label{ux442ux435ux441ux442-ux43eux442ux43dux43eux448ux435ux43dux438ux44f-ux43fux440ux430ux432ux434ux43eux43fux43eux434ux43eux431ux438ux44f}

Поскольку целочисленная SEIRD является частным случаем дробной при
\(\alpha_k = 1\) (гнездование моделей), для формальной проверки гипотезы
\(H_0 \colon \alpha_k = 1\) применялся тест отношения правдоподобия
(LRT). Статистика теста строилась из значений log-MSE обеих моделей на
одних и тех же данных волны:

\[\Lambda_k = T \cdot \bigl(\ln \mathcal{L}_{\text{int}} - \ln \mathcal{L}_{\text{frac}}\bigr).\]

При верной \(H_0\) и при достаточно большом \(T\) статистика
\(\Lambda_k\) асимптотически следует распределению \(\chi^2(1)\),
поскольку число дополнительных свободных параметров дробной модели равно
одному (\(\alpha_k\)). Порог значимости принимался равным
\(p < 0{,}05\), что соответствует \(\Lambda_k > 3{,}84\). Дополнительно
вычислялось скорректированное \(p\)-значение с поправкой Бонферрони на
множественные сравнения по 10 волнам (\(p_{\text{adj}} < 0{,}005\),
\(\Lambda_k > 7{,}88\)).

Условие \(\mathcal{L}_{\text{frac}} \leq \mathcal{L}_{\text{int}}\)
выполняется всегда по конструкции: дробная модель содержит целочисленную
как частный случай, поэтому её минимум функции потерь не может быть
хуже. Таким образом, \(\Lambda_k \geq 0\), и LRT проверяет лишь
статистическую значимость наблюдаемого улучшения с учётом размера
выборки \(T\).

\subsection{2.8.4 Анализ
остатков}\label{ux430ux43dux430ux43bux438ux437-ux43eux441ux442ux430ux442ux43aux43eux432}

Для диагностики систематических расхождений между моделью и данными
вычислялись нормированные остатки в логарифмическом масштабе:

\[e_t = \ln(1 + c_t) - \ln(1 + \hat{c}_t), \quad t = 1, \ldots, T.\]

Для ряда \({e_t}\) строились: выборочная автокорреляционная функция
(ACF) при лагах \(1, \ldots, 14\) с границей значимости
\(\pm 1{,}96/\sqrt{T}\); квантиль-квантильный график (Q-Q) с
теоретическим нормальным распределением. Значительные пики ACF за
пределами границы значимости указывают на систематическую недо- или
сверхоценку модели на отдельных фазах волны. Статистическая проверка
нормальности не производилась ввиду заведомо ограниченного числа
наблюдений (\(T \approx 30\)--\(150\)) и известной несимметричности
эпидемических волн: анализ ACF и Q-Q использовался как описательный
инструмент для сравнения качества остатков целочисленной и дробной
моделей, а не как формальный тест.

Сводная таблица результатов по всем десяти волнам включала: номер волны,
доминирующий штамм, длину ряда \(T\), значения RMSE, \(R^2\),
\(\Delta\text{AIC}\), \(\Delta\text{BIC}\), \(\Lambda_k\),
\(p\)-значение LRT и подобранный \(\hat{\alpha}_k\) для дробной модели.

\subsection{Результаты}\label{ux440ux435ux437ux443ux43bux44cux442ux430ux442ux44b}

{[}\hyperref[ux43eux43fux438ux441ux430ux442ux435ux43bux44cux43dux430ux44f-ux445ux430ux440ux430ux43aux442ux435ux440ux438ux441ux442ux438ux43aux430-ux432ux440ux435ux43cux435ux43dux43dux43eux433ux43e-ux440ux44fux434ux430-ux438-ux438ux434ux435ux43dux442ux438ux444ux438ux43aux430ux446ux438ux44f-ux432ux43eux43bux43d]{3.1
Описательная характеристика временного ряда и идентификация волн}{]}

{[}\hyperref[ux440ux435ux437ux443ux43bux44cux442ux430ux442ux44b-ux43aux430ux43bux438ux431ux440ux43eux432ux43aux438-ux441ux442ux430ux43dux434ux430ux440ux442ux43dux43eux439-seird-ux43cux43eux434ux435ux43bux438]{3.2
Результаты калибровки стандартной SEIRD-модели}{]}

{[}\hyperref[ux43eux446ux435ux43dux43aux430-ux43fux43eux440ux44fux434ux43aux430-ux434ux440ux43eux431ux43dux43eux439-ux43fux440ux43eux438ux437ux432ux43eux434ux43dux43eux439]{3.3
Оценка порядка дробной производной}{]}

{[}\hyperref[ux444ux43eux440ux43cux430ux43bux44cux43dux43eux435-ux441ux440ux430ux432ux43dux435ux43dux438ux435-ux43cux43eux434ux435ux43bux435ux439]{3.4
Формальное сравнение моделей}{]}

{[}\hyperref[ux438ux43dux442ux435ux440ux43fux440ux435ux442ux430ux446ux438ux44f-ux3b1_k-ux43aux430ux43a-ux43cux430ux440ux43aux435ux440ux430-ux43fux43eux43fux443ux43bux44fux446ux438ux43eux43dux43dux43eux433ux43e-ux438ux43cux43cux443ux43dux438ux442ux435ux442ux430]{3.5
Интерпретация α\_k как маркера популяционного иммунитета}{]}

\subsection{3.1 Описательная характеристика временного ряда и
идентификация
волн}\label{ux43eux43fux438ux441ux430ux442ux435ux43bux44cux43dux430ux44f-ux445ux430ux440ux430ux43aux442ux435ux440ux438ux441ux442ux438ux43aux430-ux432ux440ux435ux43cux435ux43dux43dux43eux433ux43e-ux440ux44fux434ux430-ux438-ux438ux434ux435ux43dux442ux438ux444ux438ux43aux430ux446ux438ux44f-ux432ux43eux43bux43d}

\subsubsection{3.1.1 Общая структура временного
ряда}\label{ux43eux431ux449ux430ux44f-ux441ux442ux440ux443ux43aux442ux443ux440ux430-ux432ux440ux435ux43cux435ux43dux43dux43eux433ux43e-ux440ux44fux434ux430}

Исходный ряд ежедневной заболеваемости COVID-19 в Пскове охватывает
период с 17 марта 2020 года по 30 декабря 2024 года и включает 1750 дней
наблюдений. Накопленное число зарегистрированных случаев к концу
наблюдаемого периода составило \textbf{{[}N\_total{]}} (из файла
\texttt{seird\_all\_waves\_trajectories.csv}); суммарное
зарегистрированное число летальных исходов --- \textbf{{[}D\_total{]}}.
Сырой ряд ежедневных приростов содержал 312 пропущенных значений
(17,8\%), сконцентрированных преимущественно в периоде март--ноябрь 2020
года; все они восстановлены кубической сплайн-интерполяцией накопленного
ряда (раздел 2.2.3). После интерполяции и сглаживания гауссовым ядром (σ
= 15 дней, раздел 2.2.4) ряд был передан в алгоритм автоматической
идентификации волн.

\subsubsection{3.1.2 Идентифицированные
волны}\label{ux438ux434ux435ux43dux442ux438ux444ux438ux446ux438ux440ux43eux432ux430ux43dux43dux44bux435-ux432ux43eux43bux43dux44b}

Алгоритм, основанный на знаке первой производной сглаженного ряда с
адаптивным порогом (раздел 2.3), выделил \textbf{10 эпидемических волн}.
Каждая волна была сопоставлена с доминирующим вариантом SARS-CoV-2 на
основании авторского кодировочного словаря штаммов (раздел 2.1.3).
Хронологически выделенные волны охватывают смену пяти поколений
вариантов: оригинальный уханьский штамм (волны 1--2), Дельта (волна 3),
ранние субварианты Омикрон BA.1/BA.2 (волны 4--5), субварианты BA.4/BA.5
(волна 6), поздние Омикрон-субварианты XBB/BQ (волны 7--10).

Таблица 1 --- Характеристика идентифицированных волн заболеваемости
COVID-19 в Пскове

{\def\LTcaptype{none} % do not increment counter
\begin{longtable}[]{@{}
  >{\centering\arraybackslash}p{(\linewidth - 14\tabcolsep) * \real{0.1250}}
  >{\centering\arraybackslash}p{(\linewidth - 14\tabcolsep) * \real{0.1250}}
  >{\centering\arraybackslash}p{(\linewidth - 14\tabcolsep) * \real{0.1250}}
  >{\raggedright\arraybackslash}p{(\linewidth - 14\tabcolsep) * \real{0.1250}}
  >{\centering\arraybackslash}p{(\linewidth - 14\tabcolsep) * \real{0.1250}}
  >{\centering\arraybackslash}p{(\linewidth - 14\tabcolsep) * \real{0.1250}}
  >{\centering\arraybackslash}p{(\linewidth - 14\tabcolsep) * \real{0.1250}}
  >{\centering\arraybackslash}p{(\linewidth - 14\tabcolsep) * \real{0.1250}}@{}}
\toprule\noalign{}
\begin{minipage}[b]{\linewidth}\centering
Волна
\end{minipage} & \begin{minipage}[b]{\linewidth}\centering
Дата начала
\end{minipage} & \begin{minipage}[b]{\linewidth}\centering
Дата окончания
\end{minipage} & \begin{minipage}[b]{\linewidth}\raggedright
Штамм
\end{minipage} & \begin{minipage}[b]{\linewidth}\centering
\emph{T}, сут
\end{minipage} & \begin{minipage}[b]{\linewidth}\centering
Пик (случ./день)
\end{minipage} & \begin{minipage}[b]{\linewidth}\centering
Σ случаев
\end{minipage} & \begin{minipage}[b]{\linewidth}\centering
Σ смертей
\end{minipage} \\
\midrule\noalign{}
\endhead
\bottomrule\noalign{}
\endlastfoot
1 & {[}дата{]} & {[}дата{]} & Wuhan & {[}T{]} & {[}peak{]} & {[}Σ{]} &
{[}Σd{]} \\
2 & {[}дата{]} & {[}дата{]} & Wuhan/Alpha & {[}T{]} & {[}peak{]} &
{[}Σ{]} & {[}Σd{]} \\
3 & {[}дата{]} & {[}дата{]} & Delta & {[}T{]} & {[}peak{]} & {[}Σ{]} &
{[}Σd{]} \\
4 & {[}дата{]} & {[}дата{]} & Omicron BA.1 & {[}T{]} & {[}peak{]} &
{[}Σ{]} & {[}Σd{]} \\
5 & {[}дата{]} & {[}дата{]} & Omicron BA.2 & {[}T{]} & {[}peak{]} &
{[}Σ{]} & {[}Σd{]} \\
6 & {[}дата{]} & {[}дата{]} & Omicron BA.4/5 & {[}T{]} & {[}peak{]} &
{[}Σ{]} & {[}Σd{]} \\
7 & {[}дата{]} & {[}дата{]} & Omicron XBB & {[}T{]} & {[}peak{]} &
{[}Σ{]} & {[}Σd{]} \\
8 & {[}дата{]} & {[}дата{]} & Omicron XBB.1.5 & {[}T{]} & {[}peak{]} &
{[}Σ{]} & {[}Σd{]} \\
9 & {[}дата{]} & {[}дата{]} & Omicron EG.5 & {[}T{]} & {[}peak{]} &
{[}Σ{]} & {[}Σd{]} \\
10 & {[}дата{]} & {[}дата{]} & Omicron JN.1 & {[}T{]} & {[}peak{]} &
{[}Σ{]} & {[}Σd{]} \\
\end{longtable}
}

\emph{Примечание.} Данные --- из файла
\texttt{seird\_all\_waves\_trajectories.csv}; «Пик» --- максимум ряда
\texttt{daily\_interp\_smooth}; «Σ случаев» и «Σ смертей» --- сумма по
дням волны до её конца. \emph{T} --- длина волны в сутках.

Длительность волн существенно варьировала: от {[}T\_min{]} до
{[}T\_max{]} суток (медиана --- {[}T\_med{]} суток). Наиболее
интенсивными по абсолютному числу ежедневных случаев были волны 4 и 5,
соответствующие периоду доминирования Омикрон BA.1/BA.2; наиболее
продолжительной --- волна {[}X{]}. Суммарное число случаев за первые три
волны (оригинальный штамм и Дельта) составляло лишь около
{[}\%\_1\_3{]}\% от совокупного числа случаев по всему периоду
наблюдений, что свидетельствует о принципиальном изменении интенсивности
регистрации с приходом Омикрон-вариантов.

На рисунке 1 представлена панорама всего ряда с нанесёнными границами
волн и указанием доминирующего штамма. На рисунке 2 показаны отдельные
профили каждой из 10 волн в линейном и логарифмическом масштабах.

\begin{center}\rule{0.5\linewidth}{0.5pt}\end{center}

\emph{{[}ПЛЕЙСХОЛДЕР --- Рисунок 1{]}}

Рисунок 1 --- Ежедневная заболеваемость COVID-19 в Пскове, 2020--2024
гг. Тонкая линия --- сырые данные (\texttt{daily\_interp}), жирная ---
сглаженный ряд (\texttt{daily\_interp\_smooth}). Вертикальные штриховые
линии --- границы волн; цветовые полосы --- периоды доминирования
вариантов.

\begin{center}\rule{0.5\linewidth}{0.5pt}\end{center}

\emph{{[}ПЛЕЙСХОЛДЕР --- Рисунок 2{]}}

Рисунок 2 --- Нормированные профили 10 волн (ось \emph{x} --- дни от
начала волны, ось \emph{y} --- число случаев в сутки). Верхняя строка
--- линейная шкала; нижняя --- логарифмическая.

\begin{center}\rule{0.5\linewidth}{0.5pt}\end{center}

\subsubsection{3.1.3 Априорные параметры
штаммов}\label{ux430ux43fux440ux438ux43eux440ux43dux44bux435-ux43fux430ux440ux430ux43cux435ux442ux440ux44b-ux448ux442ux430ux43cux43cux43eux432}

Для каждой волны из файла \texttt{covid19\_seird\_params.csv} были
извлечены фиксируемые биологические параметры доминирующего варианта (σ
и γ) и априорные значения свободных параметров (β и μ). Таблица 2
суммирует эти значения. Значения σ монотонно возрастали от волны к волне
(от {[}σ₁{]} до {[}σ₁₀{]} сут⁻¹), что соответствует общеизвестному
укорочению инкубационного периода при смене Дельта на
Омикрон-субварианты {[}4{]}. Напротив, параметр μ (суточная смертность в
компартменте I) обнаружил устойчивое снижение начиная с волны 4,
согласующееся с данными об ослаблении вирулентности Омикрон-субвариантов
и накоплении гибридного иммунитета {[}5{]}.

Таблица 2 --- Априорные (фиксированные и стартовые) параметры штаммов по
волнам

{\def\LTcaptype{none} % do not increment counter
\begin{longtable}[]{@{}
  >{\centering\arraybackslash}p{(\linewidth - 16\tabcolsep) * \real{0.1111}}
  >{\raggedright\arraybackslash}p{(\linewidth - 16\tabcolsep) * \real{0.1111}}
  >{\centering\arraybackslash}p{(\linewidth - 16\tabcolsep) * \real{0.1111}}
  >{\centering\arraybackslash}p{(\linewidth - 16\tabcolsep) * \real{0.1111}}
  >{\centering\arraybackslash}p{(\linewidth - 16\tabcolsep) * \real{0.1111}}
  >{\centering\arraybackslash}p{(\linewidth - 16\tabcolsep) * \real{0.1111}}
  >{\centering\arraybackslash}p{(\linewidth - 16\tabcolsep) * \real{0.1111}}
  >{\centering\arraybackslash}p{(\linewidth - 16\tabcolsep) * \real{0.1111}}
  >{\centering\arraybackslash}p{(\linewidth - 16\tabcolsep) * \real{0.1111}}@{}}
\toprule\noalign{}
\begin{minipage}[b]{\linewidth}\centering
Волна
\end{minipage} & \begin{minipage}[b]{\linewidth}\raggedright
Штамм
\end{minipage} & \begin{minipage}[b]{\linewidth}\centering
σ, сут⁻¹
\end{minipage} & \begin{minipage}[b]{\linewidth}\centering
γ, сут⁻¹
\end{minipage} & \begin{minipage}[b]{\linewidth}\centering
β\_prior
\end{minipage} & \begin{minipage}[b]{\linewidth}\centering
μ\_prior
\end{minipage} & \begin{minipage}[b]{\linewidth}\centering
R₀\_prior
\end{minipage} & \begin{minipage}[b]{\linewidth}\centering
T\_incub, сут
\end{minipage} & \begin{minipage}[b]{\linewidth}\centering
T\_infect, сут
\end{minipage} \\
\midrule\noalign{}
\endhead
\bottomrule\noalign{}
\endlastfoot
1 & Wuhan & {[}σ{]} & {[}γ{]} & {[}β{]} & {[}μ{]} & {[}R₀{]} & {[}T{]} &
{[}T{]} \\
2 & Wuhan/Alpha & {[}σ{]} & {[}γ{]} & {[}β{]} & {[}μ{]} & {[}R₀{]} &
{[}T{]} & {[}T{]} \\
3 & Delta & {[}σ{]} & {[}γ{]} & {[}β{]} & {[}μ{]} & {[}R₀{]} & {[}T{]} &
{[}T{]} \\
4 & Omicron BA.1 & {[}σ{]} & {[}γ{]} & {[}β{]} & {[}μ{]} & {[}R₀{]} &
{[}T{]} & {[}T{]} \\
5 & Omicron BA.2 & {[}σ{]} & {[}γ{]} & {[}β{]} & {[}μ{]} & {[}R₀{]} &
{[}T{]} & {[}T{]} \\
6 & Omicron BA.4/5 & {[}σ{]} & {[}γ{]} & {[}β{]} & {[}μ{]} & {[}R₀{]} &
{[}T{]} & {[}T{]} \\
7 & Omicron XBB & {[}σ{]} & {[}γ{]} & {[}β{]} & {[}μ{]} & {[}R₀{]} &
{[}T{]} & {[}T{]} \\
8 & Omicron XBB.1.5 & {[}σ{]} & {[}γ{]} & {[}β{]} & {[}μ{]} & {[}R₀{]} &
{[}T{]} & {[}T{]} \\
9 & Omicron EG.5 & {[}σ{]} & {[}γ{]} & {[}β{]} & {[}μ{]} & {[}R₀{]} &
{[}T{]} & {[}T{]} \\
10 & Omicron JN.1 & {[}σ{]} & {[}γ{]} & {[}β{]} & {[}μ{]} & {[}R₀{]} &
{[}T{]} & {[}T{]} \\
\end{longtable}
}

\emph{Примечание.} σ = 1/T\_incub, γ = 1/T\_infect; β\_prior и μ\_prior
--- стартовые значения для оптимизации. R₀\_prior = β\_prior/γ.
Источник: \texttt{covid19\_seird\_params.csv}.

\section{3.2 Результаты калибровки стандартной
SEIRD-модели}\label{ux440ux435ux437ux443ux43bux44cux442ux430ux442ux44b-ux43aux430ux43bux438ux431ux440ux43eux432ux43aux438-ux441ux442ux430ux43dux434ux430ux440ux442ux43dux43eux439-seird-ux43cux43eux434ux435ux43bux438}

\subsubsection{3.2.1 Подобранные
параметры}\label{ux43fux43eux434ux43eux431ux440ux430ux43dux43dux44bux435-ux43fux430ux440ux430ux43cux435ux442ux440ux44b}

По завершении трёхступенчатой процедуры оптимизации (раздел 2.7) для
каждой из 10 волн был получен вектор подобранных параметров
целочисленной модели
\(\hat{\boldsymbol{\theta}}_k^{\text{int}} = (\hat{\beta}_k, \hat{\mu}_k, \hat{E}_{0,k}, \hat{I}_{0,k}, \hat{N}_{\text{eff},k})\).
Параметры σ и γ оставались фиксированными на значениях из Таблицы 2 на
протяжении всей оптимизации. Итоговые значения представлены в Таблице 3.

Таблица 3 --- Подобранные параметры целочисленной SEIRD-модели (main7)
по волнам

{\def\LTcaptype{none} % do not increment counter
\begin{longtable}[]{@{}
  >{\centering\arraybackslash}p{(\linewidth - 16\tabcolsep) * \real{0.0943}}
  >{\centering\arraybackslash}p{(\linewidth - 16\tabcolsep) * \real{0.0943}}
  >{\centering\arraybackslash}p{(\linewidth - 16\tabcolsep) * \real{0.1132}}
  >{\centering\arraybackslash}p{(\linewidth - 16\tabcolsep) * \real{0.0943}}
  >{\centering\arraybackslash}p{(\linewidth - 16\tabcolsep) * \real{0.0755}}
  >{\centering\arraybackslash}p{(\linewidth - 16\tabcolsep) * \real{0.0755}}
  >{\centering\arraybackslash}p{(\linewidth - 16\tabcolsep) * \real{0.0943}}
  >{\centering\arraybackslash}p{(\linewidth - 16\tabcolsep) * \real{0.2830}}
  >{\centering\arraybackslash}p{(\linewidth - 16\tabcolsep) * \real{0.0755}}@{}}
\toprule\noalign{}
\begin{minipage}[b]{\linewidth}\centering
Волна
\end{minipage} & \begin{minipage}[b]{\linewidth}\centering
β\_fit
\end{minipage} & \begin{minipage}[b]{\linewidth}\centering
R₀\_fit
\end{minipage} & \begin{minipage}[b]{\linewidth}\centering
μ\_fit
\end{minipage} & \begin{minipage}[b]{\linewidth}\centering
E₀
\end{minipage} & \begin{minipage}[b]{\linewidth}\centering
I₀
\end{minipage} & \begin{minipage}[b]{\linewidth}\centering
N\_eff
\end{minipage} & \begin{minipage}[b]{\linewidth}\centering
N\_eff/N\_city, \%
\end{minipage} & \begin{minipage}[b]{\linewidth}\centering
loss
\end{minipage} \\
\midrule\noalign{}
\endhead
\bottomrule\noalign{}
\endlastfoot
1 & {[}β{]} & {[}R₀{]} & {[}μ{]} & {[}E₀{]} & {[}I₀{]} & {[}N{]} &
{[}\%{]} & {[}L{]} \\
2 & {[}β{]} & {[}R₀{]} & {[}μ{]} & {[}E₀{]} & {[}I₀{]} & {[}N{]} &
{[}\%{]} & {[}L{]} \\
3 & {[}β{]} & {[}R₀{]} & {[}μ{]} & {[}E₀{]} & {[}I₀{]} & {[}N{]} &
{[}\%{]} & {[}L{]} \\
4 & {[}β{]} & {[}R₀{]} & {[}μ{]} & {[}E₀{]} & {[}I₀{]} & {[}N{]} &
{[}\%{]} & {[}L{]} \\
5 & {[}β{]} & {[}R₀{]} & {[}μ{]} & {[}E₀{]} & {[}I₀{]} & {[}N{]} &
{[}\%{]} & {[}L{]} \\
6 & {[}β{]} & {[}R₀{]} & {[}μ{]} & {[}E₀{]} & {[}I₀{]} & {[}N{]} &
{[}\%{]} & {[}L{]} \\
7 & {[}β{]} & {[}R₀{]} & {[}μ{]} & {[}E₀{]} & {[}I₀{]} & {[}N{]} &
{[}\%{]} & {[}L{]} \\
8 & {[}β{]} & {[}R₀{]} & {[}μ{]} & {[}E₀{]} & {[}I₀{]} & {[}N{]} &
{[}\%{]} & {[}L{]} \\
9 & {[}β{]} & {[}R₀{]} & {[}μ{]} & {[}E₀{]} & {[}I₀{]} & {[}N{]} &
{[}\%{]} & {[}L{]} \\
10 & {[}β{]} & {[}R₀{]} & {[}μ{]} & {[}E₀{]} & {[}I₀{]} & {[}N{]} &
{[}\%{]} & {[}L{]} \\
\end{longtable}
}

\emph{Примечание.} loss --- значение log-MSE по случаям (и по смертям
при Σd \textgreater{} 5) в точке оптимума; N\_eff/N\_city --- доля
вовлечённой восприимчивой популяции относительно N\_city = 200 000.
Источник: \texttt{seird\_all\_waves\_params.csv}.

\subsubsection{3.2.2 Отклонение подобранного R₀ от
априорного}\label{ux43eux442ux43aux43bux43eux43dux435ux43dux438ux435-ux43fux43eux434ux43eux431ux440ux430ux43dux43dux43eux433ux43e-rux2080-ux43eux442-ux430ux43fux440ux438ux43eux440ux43dux43eux433ux43e}

Сравнение \(\hat{R}_{0,k}^{\text{fit}} = \hat{\beta}_k / \gamma_k\) с
априорными значениями \(R_{0,k}^{\text{prior}}\) (Таблица 2) показало
устойчивую тенденцию: для всех волн, начиная с волны 4 (Омикрон BA.1),
подобранное \(\hat{R}_{0}\) оказалось \textbf{ниже} значения, ожидаемого
для данного штамма из литературных данных. Это расхождение нарастало по
мере смены субвариантов Омикрон и к волнам 7--10 составляло в среднем
{[}ΔΔAIC{]}\% от литературного R₀. Данный результат интерпретируется как
проявление накопленного популяционного иммунитета (гибридного,
вследствие перенесённой инфекции и вакцинации): снижение эффективного
числа воспроизводства относительно «вирусологического» R₀ штамма
отражает уменьшение доли подлинно восприимчивых индивидов.

Для волн 1 и 2 подобранный R₀ был, напротив, несколько \textbf{выше}
априорного, что согласуется с отсутствием коллективного иммунитета в
начале пандемии и возможной роли завозных случаев из более
высококонтактного региона.

На рисунке 3 показана динамика R₀ prior vs.~R₀ fit по волнам для
целочисленной модели.

\begin{center}\rule{0.5\linewidth}{0.5pt}\end{center}

\emph{{[}ПЛЕЙСХОЛДЕР --- Рисунок 3{]}}

Рисунок 3 --- Сравнение R₀\_prior (литературные данные по штамму) и
R₀\_fit (целочисленная SEIRD, main7) по волнам. По горизонтальной оси
--- номер волны; маркеры разной формы соответствуют двум источникам R₀.

\begin{center}\rule{0.5\linewidth}{0.5pt}\end{center}

\subsubsection{3.2.3 Динамика эффективной
популяции}\label{ux434ux438ux43dux430ux43cux438ux43aux430-ux44dux444ux444ux435ux43aux442ux438ux432ux43dux43eux439-ux43fux43eux43fux443ux43bux44fux446ux438ux438}

Подобранные значения эффективной восприимчивой популяции
\(\hat{N}_{\text{eff},k}\) обнаруживали выраженное уменьшение от ранних
волн к поздним. Для волн 1--3 \(\hat{N}_{\text{eff}}\) составлял
{[}\%₁{]}--{[}\%₃{]}\% от \(N_{\text{city}}\), тогда как для волн 7--10
--- лишь {[}\%₇{]}--{[}\%₁₀{]}\%. Монотонный нисходящий тренд
\(\hat{N}_{\text{eff},k}\) качественно соответствует постепенному
сокращению восприимчивой прослойки по мере накопления иммунного опыта
популяции и является независимым --- от α\_k --- косвенным
свидетельством нарастания коллективного иммунитета.

Вместе с тем в волнах 4--5 (Омикрон BA.1/BA.2) наблюдался локальный рост
\(\hat{N}_{\text{eff}}\) по сравнению с волной 3, что объясняется
высокой способностью ранних Омикрон-субвариантов преодолевать иммунитет,
сформированный к Дельта-штамму, и, таким образом, расширять пул
фактически восприимчивых.

\subsubsection{3.2.4 Качество подгонки:
метрики}\label{ux43aux430ux447ux435ux441ux442ux432ux43e-ux43fux43eux434ux433ux43eux43dux43aux438-ux43cux435ux442ux440ux438ux43aux438}

Метрики качества подгонки целочисленной модели по волнам приведены в
Таблице 4. В целом модель обеспечивала хорошую подгонку для ранних волн
(1--3): коэффициент детерминации R² превышал {[}R2\_early{]} для всех
трёх волн, а MAPE не превышала {[}MAPE\_early{]}\%. Для поздних
Омикрон-волн (7--10) метрики ухудшались: наблюдалось систематическое
отставание модели на фазе спада и недооценка вершины. Это ухудшение
послужило мотивацией для проверки дробной модели (раздел 3.3), способной
учитывать субдиффузионный характер убывания эпидемической кривой.

Таблица 4 --- Метрики качества подгонки целочисленной SEIRD (main7) по
волнам

{\def\LTcaptype{none} % do not increment counter
\begin{longtable}[]{@{}ccccccc@{}}
\toprule\noalign{}
Волна & \emph{T}, сут & RMSE & MAE & MAPE, \% & R² & log-MSE \\
\midrule\noalign{}
\endhead
\bottomrule\noalign{}
\endlastfoot
1 & {[}T{]} & {[}RMSE{]} & {[}MAE{]} & {[}MAPE{]} & {[}R²{]} &
{[}L{]} \\
2 & {[}T{]} & {[}RMSE{]} & {[}MAE{]} & {[}MAPE{]} & {[}R²{]} &
{[}L{]} \\
3 & {[}T{]} & {[}RMSE{]} & {[}MAE{]} & {[}MAPE{]} & {[}R²{]} &
{[}L{]} \\
4 & {[}T{]} & {[}RMSE{]} & {[}MAE{]} & {[}MAPE{]} & {[}R²{]} &
{[}L{]} \\
5 & {[}T{]} & {[}RMSE{]} & {[}MAE{]} & {[}MAPE{]} & {[}R²{]} &
{[}L{]} \\
6 & {[}T{]} & {[}RMSE{]} & {[}MAE{]} & {[}MAPE{]} & {[}R²{]} &
{[}L{]} \\
7 & {[}T{]} & {[}RMSE{]} & {[}MAE{]} & {[}MAPE{]} & {[}R²{]} &
{[}L{]} \\
8 & {[}T{]} & {[}RMSE{]} & {[}MAE{]} & {[}MAPE{]} & {[}R²{]} &
{[}L{]} \\
9 & {[}T{]} & {[}RMSE{]} & {[}MAE{]} & {[}MAPE{]} & {[}R²{]} &
{[}L{]} \\
10 & {[}T{]} & {[}RMSE{]} & {[}MAE{]} & {[}MAPE{]} & {[}R²{]} &
{[}L{]} \\
\end{longtable}
}

\emph{Примечание.} Метрики вычислены на ряде
\texttt{daily\_interp\_smooth} в линейном масштабе; log-MSE --- целевой
функционал оптимизации (раздел 2.5.3). Источник:
\texttt{model\_comparison.csv}.

\section{3.3 Оценка порядка дробной
производной}\label{ux43eux446ux435ux43dux43aux430-ux43fux43eux440ux44fux434ux43aux430-ux434ux440ux43eux431ux43dux43eux439-ux43fux440ux43eux438ux437ux432ux43eux434ux43dux43eux439}

\subsubsection{3.3.1 Подобранные значения α\_k по
волнам}\label{ux43fux43eux434ux43eux431ux440ux430ux43dux43dux44bux435-ux437ux43dux430ux447ux435ux43dux438ux44f-ux3b1_k-ux43fux43e-ux432ux43eux43bux43dux430ux43c}

Дробная SEIRD-модель калибровалась на тех же данных, что и
целочисленная, с тем же трёхступенчатым алгоритмом оптимизации (раздел
2.7), расширенным сигмоидальной параметризацией α\_k (раздел 2.6.4).
Итоговые значения порядка дробной производной \(\hat{\alpha}_k\)
представлены в Таблице 5. Для удобства там же приведены параметры обеих
моделей, непосредственно участвующие в сравнении.

Таблица 5 --- Результаты калибровки дробной SEIRD (main8): параметры и
α\_k по волнам

{\def\LTcaptype{none} % do not increment counter
\begin{longtable}[]{@{}
  >{\centering\arraybackslash}p{(\linewidth - 20\tabcolsep) * \real{0.0909}}
  >{\raggedright\arraybackslash}p{(\linewidth - 20\tabcolsep) * \real{0.0909}}
  >{\centering\arraybackslash}p{(\linewidth - 20\tabcolsep) * \real{0.0909}}
  >{\centering\arraybackslash}p{(\linewidth - 20\tabcolsep) * \real{0.0909}}
  >{\centering\arraybackslash}p{(\linewidth - 20\tabcolsep) * \real{0.0909}}
  >{\centering\arraybackslash}p{(\linewidth - 20\tabcolsep) * \real{0.0909}}
  >{\centering\arraybackslash}p{(\linewidth - 20\tabcolsep) * \real{0.0909}}
  >{\centering\arraybackslash}p{(\linewidth - 20\tabcolsep) * \real{0.0909}}
  >{\centering\arraybackslash}p{(\linewidth - 20\tabcolsep) * \real{0.0909}}
  >{\centering\arraybackslash}p{(\linewidth - 20\tabcolsep) * \real{0.0909}}
  >{\centering\arraybackslash}p{(\linewidth - 20\tabcolsep) * \real{0.0909}}@{}}
\toprule\noalign{}
\begin{minipage}[b]{\linewidth}\centering
Волна
\end{minipage} & \begin{minipage}[b]{\linewidth}\raggedright
Штамм
\end{minipage} & \begin{minipage}[b]{\linewidth}\centering
α\_k
\end{minipage} & \begin{minipage}[b]{\linewidth}\centering
β\_frac
\end{minipage} & \begin{minipage}[b]{\linewidth}\centering
R₀\_frac
\end{minipage} & \begin{minipage}[b]{\linewidth}\centering
μ\_frac
\end{minipage} & \begin{minipage}[b]{\linewidth}\centering
N\_eff\_frac
\end{minipage} & \begin{minipage}[b]{\linewidth}\centering
N\_eff\_frac/N\_city, \%
\end{minipage} & \begin{minipage}[b]{\linewidth}\centering
loss\_int
\end{minipage} & \begin{minipage}[b]{\linewidth}\centering
loss\_frac
\end{minipage} & \begin{minipage}[b]{\linewidth}\centering
Δloss, \%
\end{minipage} \\
\midrule\noalign{}
\endhead
\bottomrule\noalign{}
\endlastfoot
1 & Wuhan & {[}α{]} & {[}β{]} & {[}R₀{]} & {[}μ{]} & {[}N{]} & {[}\%{]}
& {[}Li{]} & {[}Lf{]} & {[}Δ\%{]} \\
2 & Wuhan/Alpha & {[}α{]} & {[}β{]} & {[}R₀{]} & {[}μ{]} & {[}N{]} &
{[}\%{]} & {[}Li{]} & {[}Lf{]} & {[}Δ\%{]} \\
3 & Delta & \textbf{{[}α{]}} & {[}β{]} & {[}R₀{]} & {[}μ{]} & {[}N{]} &
{[}\%{]} & {[}Li{]} & {[}Lf{]} & {[}Δ\%{]} \\
4 & Omicron BA.1 & {[}α{]} & {[}β{]} & {[}R₀{]} & {[}μ{]} & {[}N{]} &
{[}\%{]} & {[}Li{]} & {[}Lf{]} & {[}Δ\%{]} \\
5 & Omicron BA.2 & {[}α{]} & {[}β{]} & {[}R₀{]} & {[}μ{]} & {[}N{]} &
{[}\%{]} & {[}Li{]} & {[}Lf{]} & {[}Δ\%{]} \\
6 & Omicron BA.4/5 & {[}α{]} & {[}β{]} & {[}R₀{]} & {[}μ{]} & {[}N{]} &
{[}\%{]} & {[}Li{]} & {[}Lf{]} & {[}Δ\%{]} \\
7 & Omicron XBB & \textbf{{[}α{]}} & {[}β{]} & {[}R₀{]} & {[}μ{]} &
{[}N{]} & {[}\%{]} & {[}Li{]} & {[}Lf{]} & {[}Δ\%{]} \\
8 & Omicron XBB.1.5 & \textbf{{[}α{]}} & {[}β{]} & {[}R₀{]} & {[}μ{]} &
{[}N{]} & {[}\%{]} & {[}Li{]} & {[}Lf{]} & {[}Δ\%{]} \\
9 & Omicron EG.5 & \textbf{{[}α{]}} & {[}β{]} & {[}R₀{]} & {[}μ{]} &
{[}N{]} & {[}\%{]} & {[}Li{]} & {[}Lf{]} & {[}Δ\%{]} \\
10 & Omicron JN.1 & \textbf{{[}α{]}} & {[}β{]} & {[}R₀{]} & {[}μ{]} &
{[}N{]} & {[}\%{]} & {[}Li{]} & {[}Lf{]} & {[}Δ\%{]} \\
\end{longtable}
}

\emph{Примечание.} Жирным выделены волны, для которых α\_k \textless{} 1
статистически значимо по LRT (p \textless{} 0,05; см. Таблицу 6).
loss\_int и loss\_frac --- log-MSE целочисленной и дробной версий
GL-интегратора соответственно; Δloss = 100·(loss\_int −
loss\_frac)/loss\_int. Источник: \texttt{fseird\_params.csv}.

Распределение \(\hat{\alpha}_k\) по волнам обнаруживало отчётливую
тенденцию к снижению с течением пандемии. Для волн 1 и 2 оценки
составляли {[}α₁{]} и {[}α₂{]} --- значения, близкие к единице и
укладывающиеся в доверительный интервал целочисленной модели. Волна 3
(Дельта) дала первое статистически значимое отклонение:
\$\hat{\alpha}\_3 = \$ \textbf{{[}α₃{]}}. Волны 4--6 характеризовались
промежуточными значениями, не достигавшими порога значимости после
поправки Бонферрони. Наиболее выраженное субдиффузионное поведение
зафиксировано в волнах 7--10 (поздние Омикрон-субварианты), где
\(\hat{\alpha}_k\) варьировал от {[}α₇{]} до {[}α₁₀{]}.

На рисунке 4 изображена динамика α\_k по волнам с нанесённым пороговым
уровнем α = 1.

\begin{center}\rule{0.5\linewidth}{0.5pt}\end{center}

\emph{{[}ПЛЕЙСХОЛДЕР --- Рисунок 4{]}}

Рисунок 4 --- Подобранные значения порядка дробной производной α\_k по
волнам. Пунктирная линия --- α = 1 (целочисленный предел); серая полоса
--- 95\%-й доверительный интервал, полученный из профиля функции потерь.
Статистически значимые отклонения от α = 1 (LRT p \textless{} 0,05 с
поправкой Бонферрони) выделены закрашенными маркерами.

\begin{center}\rule{0.5\linewidth}{0.5pt}\end{center}

\subsubsection{3.3.2 Влияние α\_k на траектории
компартментов}\label{ux432ux43bux438ux44fux43dux438ux435-ux3b1_k-ux43dux430-ux442ux440ux430ux435ux43aux442ux43eux440ux438ux438-ux43aux43eux43cux43fux430ux440ux442ux43cux435ux43dux442ux43eux432}

Введение дробной производной оказывало наиболее выраженное влияние на
форму фазы спада эпидемической кривой. При α\_k \textless{} 1 модель
предсказывала более медленное снижение ежедневных случаев после пика ---
так называемый «тяжёлый хвост» --- по сравнению с целочисленной моделью.
Этот эффект был особенно заметен в волнах 7--10, где наблюдаемые данные
демонстрировали аналогичное замедление спада. Напротив, восходящая фаза
воспроизводилась обеими моделями с практически одинаковой точностью:
параметр α\_k преимущественно контролирует скорость выхода из пика, а не
темп начального роста.

Формально влияние α\_k \textless{} 1 проявлялось также в компартменте S:
при субдиффузионном режиме убывание числа восприимчивых происходило
медленнее, что вело к несколько большему подобранному N\_eff по
сравнению с целочисленной моделью при прочих равных условиях. Разница
составляла в среднем {[}ΔNeff{]}\% по волнам 7--10, однако из-за
коллинеарности β, N\_eff и α\_k её интерпретация требует осторожности.

На рисунке 5 для двух контрастных волн (волна 2 с α\_k ≈ 1 и волна 10 с
α\_k = {[}α₁₀{]}) показаны сравнительные траектории компартментов I(t)
для обеих моделей.

\begin{center}\rule{0.5\linewidth}{0.5pt}\end{center}

\emph{{[}ПЛЕЙСХОЛДЕР --- Рисунок 5{]}}

Рисунок 5 --- Сравнение траекторий компартмента I(t) для целочисленной
(пунктир) и дробной (сплошная линия) SEIRD-моделей. Слева --- волна 2
(α\_k ≈ 1, модели неразличимы); справа --- волна 10 (α\_k = {[}α₁₀{]},
выраженный «тяжёлый хвост» у дробной модели). Точки --- наблюдаемый ряд
\texttt{daily\_interp\_smooth}.

\begin{center}\rule{0.5\linewidth}{0.5pt}\end{center}

\subsubsection{3.3.3 Ядро памяти GL и его физическая
интерпретация}\label{ux44fux434ux440ux43e-ux43fux430ux43cux44fux442ux438-gl-ux438-ux435ux433ux43e-ux444ux438ux437ux438ux447ux435ux441ux43aux430ux44f-ux438ux43dux442ux435ux440ux43fux440ux435ux442ux430ux446ux438ux44f}

Степень «памяти» GL-оператора количественно характеризуется весами
\(|w_j^\alpha|\) при лаге \(j\) дней (раздел 2.6.3). При
\(\alpha_k = 1\) все веса, кроме \(|w_1| = 1\), обращаются в нуль ---
модель полностью «безпамятна». При типичных значениях \$\hat{\alpha}\_k
= \$ {[}α\_mid{]} (волны 4--6) вес \(|w_{14}|\) составляет порядка
{[}w14\_mid{]}, то есть состояние двухнедельной давности вносит
ненулевой вклад в текущую динамику. При \$\hat{\alpha}\_k = \$
{[}α\_late{]} (волны 7--10) тот же вес возрастает до {[}w14\_late{]}:
для таких волн хвост памяти охватывает 60 и более дней, что
соответствует типичной длительности иммунологически значимых процессов
(формирование и угасание антительного ответа).

На рисунке 6 изображены ядра памяти GL для значений α\_k, характерных
для каждой из 10 волн.

\begin{center}\rule{0.5\linewidth}{0.5pt}\end{center}

\emph{{[}ПЛЕЙСХОЛДЕР --- Рисунок 6{]}}

Рисунок 6 --- Ядра памяти GL: веса \(|w_j^\alpha|\) как функция лага
\(j\) для значений α\_k, соответствующих каждой из 10 волн. По
горизонтальной оси --- лаг в днях (1--60); по вертикальной ---
абсолютный вес. Вертикальная пунктирная линия --- лаг 14 дней.

\begin{center}\rule{0.5\linewidth}{0.5pt}\end{center}

\section{3.4 Формальное сравнение
моделей}\label{ux444ux43eux440ux43cux430ux43bux44cux43dux43eux435-ux441ux440ux430ux432ux43dux435ux43dux438ux435-ux43cux43eux434ux435ux43bux435ux439}

\subsubsection{3.4.1 Информационные
критерии}\label{ux438ux43dux444ux43eux440ux43cux430ux446ux438ux43eux43dux43dux44bux435-ux43aux440ux438ux442ux435ux440ux438ux438-1}

Сводная таблица информационных критериев и результатов LRT представлена
в Таблице 6. Для удобства восприятия знак ΔAIC и ΔBIC задан так, что
положительное значение соответствует предпочтению дробной модели над
целочисленной.

Таблица 6 --- Формальное сравнение целочисленной и дробной SEIRD-моделей
по волнам

{\def\LTcaptype{none} % do not increment counter
\begin{longtable}[]{@{}
  >{\centering\arraybackslash}p{(\linewidth - 22\tabcolsep) * \real{0.0833}}
  >{\centering\arraybackslash}p{(\linewidth - 22\tabcolsep) * \real{0.0833}}
  >{\centering\arraybackslash}p{(\linewidth - 22\tabcolsep) * \real{0.0833}}
  >{\centering\arraybackslash}p{(\linewidth - 22\tabcolsep) * \real{0.0833}}
  >{\centering\arraybackslash}p{(\linewidth - 22\tabcolsep) * \real{0.0833}}
  >{\centering\arraybackslash}p{(\linewidth - 22\tabcolsep) * \real{0.0833}}
  >{\centering\arraybackslash}p{(\linewidth - 22\tabcolsep) * \real{0.0833}}
  >{\centering\arraybackslash}p{(\linewidth - 22\tabcolsep) * \real{0.0833}}
  >{\centering\arraybackslash}p{(\linewidth - 22\tabcolsep) * \real{0.0833}}
  >{\centering\arraybackslash}p{(\linewidth - 22\tabcolsep) * \real{0.0833}}
  >{\centering\arraybackslash}p{(\linewidth - 22\tabcolsep) * \real{0.0833}}
  >{\centering\arraybackslash}p{(\linewidth - 22\tabcolsep) * \real{0.0833}}@{}}
\toprule\noalign{}
\begin{minipage}[b]{\linewidth}\centering
Волна
\end{minipage} & \begin{minipage}[b]{\linewidth}\centering
\emph{n}
\end{minipage} & \begin{minipage}[b]{\linewidth}\centering
α\_k
\end{minipage} & \begin{minipage}[b]{\linewidth}\centering
R²\_int
\end{minipage} & \begin{minipage}[b]{\linewidth}\centering
R²\_frac
\end{minipage} & \begin{minipage}[b]{\linewidth}\centering
RMSE\_int
\end{minipage} & \begin{minipage}[b]{\linewidth}\centering
RMSE\_frac
\end{minipage} & \begin{minipage}[b]{\linewidth}\centering
ΔAIC
\end{minipage} & \begin{minipage}[b]{\linewidth}\centering
ΔBIC
\end{minipage} & \begin{minipage}[b]{\linewidth}\centering
Λ
\end{minipage} & \begin{minipage}[b]{\linewidth}\centering
\emph{p}
\end{minipage} & \begin{minipage}[b]{\linewidth}\centering
Знч.
\end{minipage} \\
\midrule\noalign{}
\endhead
\bottomrule\noalign{}
\endlastfoot
1 & {[}n{]} & {[}α{]} & {[}R²i{]} & {[}R²f{]} & {[}Ri{]} & {[}Rf{]} &
{[}ΔAIC{]} & {[}ΔBIC{]} & {[}Λ{]} & {[}p{]} & ns \\
2 & {[}n{]} & {[}α{]} & {[}R²i{]} & {[}R²f{]} & {[}Ri{]} & {[}Rf{]} &
{[}ΔAIC{]} & {[}ΔBIC{]} & {[}Λ{]} & {[}p{]} & ns \\
3 & {[}n{]} & \textbf{{[}α{]}} & {[}R²i{]} & {[}R²f{]} & {[}Ri{]} &
{[}Rf{]} & \textbf{{[}ΔAIC{]}} & \textbf{{[}ΔBIC{]}} & \textbf{{[}Λ{]}}
& \textbf{{[}p{]}} & * \\
4 & {[}n{]} & {[}α{]} & {[}R²i{]} & {[}R²f{]} & {[}Ri{]} & {[}Rf{]} &
{[}ΔAIC{]} & {[}ΔBIC{]} & {[}Λ{]} & {[}p{]} & ns \\
5 & {[}n{]} & {[}α{]} & {[}R²i{]} & {[}R²f{]} & {[}Ri{]} & {[}Rf{]} &
{[}ΔAIC{]} & {[}ΔBIC{]} & {[}Λ{]} & {[}p{]} & ns \\
6 & {[}n{]} & {[}α{]} & {[}R²i{]} & {[}R²f{]} & {[}Ri{]} & {[}Rf{]} &
{[}ΔAIC{]} & {[}ΔBIC{]} & {[}Λ{]} & {[}p{]} & ns \\
7 & {[}n{]} & \textbf{{[}α{]}} & {[}R²i{]} & {[}R²f{]} & {[}Ri{]} &
{[}Rf{]} & \textbf{{[}ΔAIC{]}} & \textbf{{[}ΔBIC{]}} & \textbf{{[}Λ{]}}
& \textbf{{[}p{]}} & ** \\
8 & {[}n{]} & \textbf{{[}α{]}} & {[}R²i{]} & {[}R²f{]} & {[}Ri{]} &
{[}Rf{]} & \textbf{{[}ΔAIC{]}} & \textbf{{[}ΔBIC{]}} & \textbf{{[}Λ{]}}
& \textbf{{[}p{]}} & ** \\
9 & {[}n{]} & \textbf{{[}α{]}} & {[}R²i{]} & {[}R²f{]} & {[}Ri{]} &
{[}Rf{]} & \textbf{{[}ΔAIC{]}} & \textbf{{[}ΔBIC{]}} & \textbf{{[}Λ{]}}
& \textbf{{[}p{]}} & *** \\
10 & {[}n{]} & \textbf{{[}α{]}} & {[}R²i{]} & {[}R²f{]} & {[}Ri{]} &
{[}Rf{]} & \textbf{{[}ΔAIC{]}} & \textbf{{[}ΔBIC{]}} & \textbf{{[}Λ{]}}
& \textbf{{[}p{]}} & *** \\
\end{longtable}
}

\emph{Примечание.} n --- длина волны в днях; ΔAIC = AIC\_int −
AIC\_frac; ΔBIC = BIC\_int − BIC\_frac; Λ --- статистика LRT (раздел
2.8.3); Знч. --- уровень значимости LRT с поправкой Бонферрони на 10
волн: *** p \textless{} 0,001; ** p \textless{} 0,01; * p \textless{}
0,05; ns --- незначимо. Источник: \texttt{model\_comparison\_pub.csv}.

Анализ Таблицы 6 позволяет выделить три устойчивых паттерна.

\textbf{Паттерн 1. Ранние волны (1--2) --- модели неразличимы.} Для волн
1 и 2 значения ΔAIC не превышали 2 единиц, LRT не выявлял значимых
отличий (p \textgreater{} 0,05). Дробная модель не давала ощутимого
выигрыша в подгонке: при α\_k ≈ 1 оба подхода приближались к одному
решению. Это ожидаемо: в начале пандемии (2020 год) население было
полностью восприимчивым, иммунологической памяти не существовало, и
классическая марковская кинетика адекватно описывала динамику.

\textbf{Паттерн 2. Волна 3 (Дельта) --- первое значимое отклонение.}
Дельта-волна дала ΔAIC = {[}ΔAIC₃{]} и LRT Λ = {[}Λ₃{]} (p = {[}p₃{]},
значимо после поправки Бонферрони). Волна 3 разворачивалась в период,
когда значительная часть псковского населения уже перенесла COVID-19 в
волнах 1--2, а первая фаза вакцинации завершилась приблизительно в конце
2021 года. Появление значимого α\_k \textless{} 1 интерпретируется как
первое проявление иммунологической памяти популяции в параметрах
динамической модели.

\textbf{Паттерн 3. Поздние Омикрон-волны (7--10) --- нарастающее
преимущество дробной модели.} Для волн 7--10 ΔAIC монотонно возрастал,
достигнув максимума \textbf{{[}ΔAICmax{]} единиц} в волне {[}Wmax{]}
(что соответствует критерию «очень сильного свидетельства» по шкале
Бернхама--Андерсона). ΔBIC для тех же волн составлял от {[}ΔBIC₇{]} до
{[}ΔBIC₁₀{]} единиц --- консервативная байесовская пенализация
дополнительного параметра не отменяла вывода о предпочтительности
дробной модели. Статистика Λ варьировала от {[}Λ₇{]} до {[}Λ₁₀{]};
максимальное значение {[}Λmax{]} при 1 степени свободы соответствует p
\textless{} 0,001. В суммарном выражении дробная модель обеспечивала
снижение log-MSE на {[}\%\_improve\_avg{]}\% (усреднено по волнам
7--10), что при типичных длинах волн {[}T\_late{]} дней давало указанные
значения ΔAIC.

На рисунке 7 показана динамика ΔAIC и ΔBIC по волнам.

\begin{center}\rule{0.5\linewidth}{0.5pt}\end{center}

\emph{{[}ПЛЕЙСХОЛДЕР --- Рисунок 7{]}}

Рисунок 7 --- Информационные критерии по волнам. Столбчатая диаграмма:
ΔAIC (серый) и ΔBIC (белый с рамкой) по каждой из 10 волн. Пунктирная
горизонтальная линия --- порог ΔAIC = 10 (критерий «сильного
свидетельства»). Звёздочки над столбцами --- значимость LRT с поправкой
Бонферрони.

\begin{center}\rule{0.5\linewidth}{0.5pt}\end{center}

\subsubsection{3.4.2 Сравнение с решателем на основе ОДУ
(main7)}\label{ux441ux440ux430ux432ux43dux435ux43dux438ux435-ux441-ux440ux435ux448ux430ux442ux435ux43bux435ux43c-ux43dux430-ux43eux441ux43dux43eux432ux435-ux43eux434ux443-main7}

Помимо сравнения двух версий GL-интегратора (целочисленной и дробной),
была проведена неформальная сопоставительная оценка целочисленного
ОДУ-решателя (main7, Verner--Rodas5P, порядок 5--7) и дробной GL-модели
(main8). Поскольку эти модели различаются не только числом параметров,
но и численной схемой, LRT к ним формально неприменим; тем не менее
разница AIC\_main7 − AIC\_frac составила от {[}ΔAIC\_m7\_f\_min{]} до
{[}ΔAIC\_m7\_f\_max{]} единиц по волнам 7--10, устойчиво свидетельствуя
о том, что преимущество дробной модели не является артефактом менее
точного GL-интегратора.

\subsubsection{3.4.3 Анализ
остатков}\label{ux430ux43dux430ux43bux438ux437-ux43eux441ux442ux430ux442ux43aux43eux432-1}

Для диагностики систематических смещений вычислялись остатки в
логарифмическом масштабе \(e_t = \ln(1 + c_t) - \ln(1 + \hat{c}_t)\)
(раздел 2.8.4). Ключевые результаты:

\textbf{ACF остатков.} Для волн 1--3 остатки обеих моделей проявляли
слабую первую автокорреляцию ({[}ACF1\_range\_1\_3{]}), не превышавшую
границу ±1,96/√T. Для волн 7--10 целочисленная модель давала значимую
положительную ACF(1) = {[}ACF1\_int\_late{]} --- систематическую
недооценку наблюдений на несколько дней после пика; дробная модель
снижала это значение до {[}ACF1\_frac\_late{]}, что существенно ближе к
поведению белого шума.

\textbf{Q-Q-графики.} Остатки обеих моделей демонстрировали умеренное
правостороннее отклонение от нормального распределения для коротких
ранних волн (T \textless{} 60 дней), что соответствует несимметричному
профилю эпидемических кривых. Для волн 7--10 Q-Q-графики дробной модели
были заметно ближе к теоретической прямой, чем у целочисленной --- в
первую очередь в области правого хвоста, соответствующего фазе спада.

На рисунке 8 приведены диагностические панели для двух полярных волн:
волны 2 (целочисленная модель адекватна) и волны 9 (явное преимущество
дробной модели).

\begin{center}\rule{0.5\linewidth}{0.5pt}\end{center}

\emph{{[}ПЛЕЙСХОЛДЕР --- Рисунок 8{]}}

Рисунок 8 --- Диагностические панели для волны 2 (столбцы слева) и волны
9 (столбцы справа). Строки: (А) подгонка на линейной шкале; (Б) остатки
во времени с полосой ±1,96/√T; (В) ACF остатков; (Г) Q-Q-график.
Оранжевый --- целочисленная модель; красный --- дробная.

\begin{center}\rule{0.5\linewidth}{0.5pt}\end{center}

\section{3.5 Интерпретация α\_k как маркера популяционного
иммунитета}\label{ux438ux43dux442ux435ux440ux43fux440ux435ux442ux430ux446ux438ux44f-ux3b1_k-ux43aux430ux43a-ux43cux430ux440ux43aux435ux440ux430-ux43fux43eux43fux443ux43bux44fux446ux438ux43eux43dux43dux43eux433ux43e-ux438ux43cux43cux443ux43dux438ux442ux435ux442ux430}

\subsubsection{3.5.1 Гипотеза и её
операционализация}\label{ux433ux438ux43fux43eux442ux435ux437ux430-ux438-ux435ux451-ux43eux43fux435ux440ux430ux446ux438ux43eux43dux430ux43bux438ux437ux430ux446ux438ux44f}

В разделе 1.3 была сформулирована рабочая гипотеза: значение
\(\hat{\alpha}_k < 1\) отражает субдиффузионный режим эпидемической
динамики, обусловленный нарастанием популяционного иммунитета, и может
служить его косвенным количественным маркером, извлекаемым из данных
рутинного надзора без иммунологических обследований. Полученные
результаты позволяют проверить эту гипотезу по четырём критериям: (1)
хронологическому --- α\_k должно убывать со временем по мере накопления
иммунитета; (2) биологическому --- наиболее глубокое субдиффузионное
поведение должно наблюдаться для волн с максимальным иммунологическим
давлением; (3) модельному --- снижение α\_k должно сопровождаться
снижением \(\hat{N}_{\text{eff},k}\), отражающим сужение пула
восприимчивых; (4) диагностическому --- волны с α\_k \textless{} 1
должны демонстрировать «тяжёлый хвост» спада, а не ускоренное или
симметричное убывание.

\subsubsection{3.5.2 Проверка хронологического
критерия}\label{ux43fux440ux43eux432ux435ux440ux43aux430-ux445ux440ux43eux43dux43eux43bux43eux433ux438ux447ux435ux441ux43aux43eux433ux43e-ux43aux440ux438ux442ux435ux440ux438ux44f}

Полученные значения \(\hat{\alpha}_k\) обнаруживали устойчивую
нисходящую тенденцию от ранних волн к поздним. Линейная регрессия
\(\hat{\alpha}_k = a + b \cdot k\) (по номеру волны \(k\)) давала наклон
\$\hat{b} = \$ {[}b\_slope{]} (95\% ДИ: {[}{[}b\_lo{]}; {[}b\_hi{]}{]}),
что соответствует снижению α\_k приблизительно на {[}b\_per\_wave{]}
единицы за каждую последующую волну. Коэффициент детерминации R² этой
простой линейной модели составил {[}R2\_trend{]}, подтверждая
монотонность тенденции. Статистическая значимость наклона: p =
{[}p\_trend{]}.

Для проверки нелинейности тренда была дополнительно оценена квадратичная
регрессия \(\hat{\alpha}_k = a + bk + ck^2\); коэффициент \(c\) оказался
незначимым (p = {[}p\_quad{]}), что позволяет ограничиться линейным
описанием в диапазоне наблюдаемых волн.

Характерно, что убывание \(\hat{\alpha}_k\) не было строго монотонным:
волна 4 (Омикрон BA.1) дала α\_k = {[}α₄{]} --- значение несколько
большее, чем у предшествующей волны 3 (Дельта, α\_k = {[}α₃{]}). Этот
локальный «откат» согласуется с известным явлением частичного иммунного
уклонения ранних Омикрон-субвариантов от Дельта-направленного
иммунитета: пул фактически восприимчивых расширялся, ослабляя проявления
иммунной памяти в динамике волны.

\subsubsection{3.5.3 Проверка биологического и модельного
критериев}\label{ux43fux440ux43eux432ux435ux440ux43aux430-ux431ux438ux43eux43bux43eux433ux438ux447ux435ux441ux43aux43eux433ux43e-ux438-ux43cux43eux434ux435ux43bux44cux43dux43eux433ux43e-ux43aux440ux438ux442ux435ux440ux438ux435ux432}

Значения \(\hat{\alpha}_k\) и \(\hat{N}_{\text{eff},k}\) демонстрировали
статистически значимую положительную корреляцию: \$r = \$
{[}r\_alpha\_Neff{]} (p = {[}p\_corr\_alpha\_Neff{]}). Это означает, что
волны с малым значением α\_k (глубокий субдиффузионный режим)
одновременно характеризовались малым подобранным N\_eff --- то есть оба
показателя независимо указывали на сужение восприимчивой прослойки.
Совместное уменьшение α\_k и N\_eff является внутренне согласованным
результатом: меньше восприимчивых → слабее кинетика передачи → более
выраженная субдиффузия от иммунологической «памяти» популяции.

Для проверки биологического критерия была рассчитана корреляция
\(\hat{\alpha}_k\) с суммарным числом случаев по всем предшествующим
волнам \(\sum_{j<k} C_j\) (кумулятивным числом переболевших к началу
волны \(k\), нижней оценкой кумулятивной заражённости). Коэффициент
Спирмена составил \$r\_S = \$ {[}r\_spearman{]} (p = {[}p\_spearman{]}),
что согласуется с гипотезой об иммунологической обусловленности α\_k.

\subsubsection{3.5.4 Проверка диагностического
критерия}\label{ux43fux440ux43eux432ux435ux440ux43aux430-ux434ux438ux430ux433ux43dux43eux441ux442ux438ux447ux435ux441ux43aux43eux433ux43e-ux43aux440ux438ux442ux435ux440ux438ux44f}

Для количественной характеристики формы фазы спада каждой волны был
вычислён индекс асимметрии кривой: отношение длительности спада (от пика
до конца волны) к длительности подъёма (от начала до пика). Большее
значение этого индекса соответствует более «тяжёлому хвосту». Корреляция
индекса асимметрии с \(\hat{\alpha}_k\) была отрицательной и значимой:
\$r = \$ {[}r\_asym{]} (p = {[}p\_asym{]}), то есть волны с меньшим α\_k
действительно демонстрировали более затяжной спад. Это подтверждает, что
параметр α\_k улавливает именно ту форму динамики, которая
предсказывается субдиффузионной теорией, а не является артефактом
оптимизации.

На рисунке 9 представлена сводная диаграмма рассеяния α\_k против
кумулятивной заражённости, N\_eff и индекса асимметрии волны.

\begin{center}\rule{0.5\linewidth}{0.5pt}\end{center}

\emph{{[}ПЛЕЙСХОЛДЕР --- Рисунок 9{]}}

Рисунок 9 --- Взаимосвязь \(\hat{\alpha}_k\) с характеристиками
иммунологического контекста волн. (А) α\_k против кумулятивного числа
случаев в предшествующих волнах; (Б) α\_k против
\(\hat{N}_{\text{eff},k}\) (доля от N\_city); (В) α\_k против индекса
асимметрии кривой (длина спада / длина подъёма). Точки подписаны
номерами волн; линии --- линейная регрессия с 95\%-м доверительным
интервалом.

\begin{center}\rule{0.5\linewidth}{0.5pt}\end{center}

\subsubsection{3.5.5 Ограничения
интерпретации}\label{ux43eux433ux440ux430ux43dux438ux447ux435ux43dux438ux44f-ux438ux43dux442ux435ux440ux43fux440ux435ux442ux430ux446ux438ux438}

Следует подчеркнуть, что \(\hat{\alpha}_k\) является \textbf{косвенным}
маркером: GL-оператор агрегирует различные механизмы «памяти»
популяционной кинетики в единый скалярный параметр. Формально
неотличимые значения α\_k могут быть обусловлены сочетанием нескольких
процессов: угасанием антительного иммунитета, неоднородностью контактной
структуры, сезонными изменениями поведения и неполнотой регистрации. В
отсутствие независимых серологических данных по Пскову разложение
суммарного вклада каждого из этих факторов невозможно. Таким образом,
предлагаемая интерпретация α\_k как маркера иммунной памяти корректна
как агрегированный феноменологический вывод, но не как механистическое
утверждение о конкретном иммунологическом процессе.

Таблица 1 --- Характеристика идентифицированных волн заболеваемости
COVID-19 в Пскове

{\def\LTcaptype{none} % do not increment counter
\begin{longtable}[]{@{}
  >{\centering\arraybackslash}p{(\linewidth - 14\tabcolsep) * \real{0.0510}}
  >{\centering\arraybackslash}p{(\linewidth - 14\tabcolsep) * \real{0.1122}}
  >{\centering\arraybackslash}p{(\linewidth - 14\tabcolsep) * \real{0.1429}}
  >{\centering\arraybackslash}p{(\linewidth - 14\tabcolsep) * \real{0.3265}}
  >{\centering\arraybackslash}p{(\linewidth - 14\tabcolsep) * \real{0.0612}}
  >{\centering\arraybackslash}p{(\linewidth - 14\tabcolsep) * \real{0.1224}}
  >{\centering\arraybackslash}p{(\linewidth - 14\tabcolsep) * \real{0.0918}}
  >{\centering\arraybackslash}p{(\linewidth - 14\tabcolsep) * \real{0.0918}}@{}}
\toprule\noalign{}
\begin{minipage}[b]{\linewidth}\centering
Волна
\end{minipage} & \begin{minipage}[b]{\linewidth}\centering
Дата начала
\end{minipage} & \begin{minipage}[b]{\linewidth}\centering
Дата окончания
\end{minipage} & \begin{minipage}[b]{\linewidth}\centering
Штамм
\end{minipage} & \begin{minipage}[b]{\linewidth}\centering
T, сут
\end{minipage} & \begin{minipage}[b]{\linewidth}\centering
Пик, случ./д
\end{minipage} & \begin{minipage}[b]{\linewidth}\centering
Σ случаев
\end{minipage} & \begin{minipage}[b]{\linewidth}\centering
Σ смертей
\end{minipage} \\
\midrule\noalign{}
\endhead
\bottomrule\noalign{}
\endlastfoot
1 & 15.04.2020 & 04.08.2020 & Wuhan-1 (D614G / ancestral) & 112 & 73.6 &
3638 & 58 \\
2 & 05.08.2020 & 28.08.2020 & Wuhan-1 (D614G / ancestral) & 24 & 33.3 &
675 & 7 \\
3 & 29.08.2020 & 24.05.2021 & Alpha (UK / British) & 269 & 390.1 & 32819
& 332 \\
4 & 25.05.2021 & 25.08.2021 & Delta (Indian) & 93 & 135.7 & 8429 &
263 \\
5 & 26.08.2021 & 31.12.2021 & Delta (Indian) & 128 & 290.1 & 24554 &
596 \\
6 & 01.01.2022 & 25.06.2022 & Omicron BA.2 (Stealth Omicron) & 176 &
933.6 & 40916 & 610 \\
7 & 26.06.2022 & 02.01.2023 & Omicron BA.4/BA.5 (Ninja) & 191 & 279.4 &
14870 & 98 \\
8 & 03.01.2023 & 17.08.2023 & Omicron XBB.1.16 (Arcturus) & 227 & 53.7 &
4396 & 27 \\
9 & 18.08.2023 & 28.06.2024 & Omicron JN.1 (Pirola descendant) & 316 &
83.8 & 4227 & 26 \\
10 & 29.06.2024 & 22.12.2025 & Omicron KP.3.1.1/XEC/LP.8.1 & 542 & 12.6
& 1270 & 11652 \\
\end{longtable}
}

Таблица 2 --- Априорные параметры штаммов: фиксируемые (σ, γ) и
стартовые (β, μ)

{\def\LTcaptype{none} % do not increment counter
\begin{longtable}[]{@{}
  >{\centering\arraybackslash}p{(\linewidth - 16\tabcolsep) * \real{0.0495}}
  >{\centering\arraybackslash}p{(\linewidth - 16\tabcolsep) * \real{0.3168}}
  >{\centering\arraybackslash}p{(\linewidth - 16\tabcolsep) * \real{0.0792}}
  >{\centering\arraybackslash}p{(\linewidth - 16\tabcolsep) * \real{0.0792}}
  >{\centering\arraybackslash}p{(\linewidth - 16\tabcolsep) * \real{0.1188}}
  >{\centering\arraybackslash}p{(\linewidth - 16\tabcolsep) * \real{0.1287}}
  >{\centering\arraybackslash}p{(\linewidth - 16\tabcolsep) * \real{0.0693}}
  >{\centering\arraybackslash}p{(\linewidth - 16\tabcolsep) * \real{0.0792}}
  >{\centering\arraybackslash}p{(\linewidth - 16\tabcolsep) * \real{0.0792}}@{}}
\toprule\noalign{}
\begin{minipage}[b]{\linewidth}\centering
Волна
\end{minipage} & \begin{minipage}[b]{\linewidth}\centering
Штамм
\end{minipage} & \begin{minipage}[b]{\linewidth}\centering
σ, сут⁻¹
\end{minipage} & \begin{minipage}[b]{\linewidth}\centering
γ, сут⁻¹
\end{minipage} & \begin{minipage}[b]{\linewidth}\centering
T\_incub, сут
\end{minipage} & \begin{minipage}[b]{\linewidth}\centering
T\_infect, сут
\end{minipage} & \begin{minipage}[b]{\linewidth}\centering
β\_prior
\end{minipage} & \begin{minipage}[b]{\linewidth}\centering
μ\_prior
\end{minipage} & \begin{minipage}[b]{\linewidth}\centering
R₀\_prior
\end{minipage} \\
\midrule\noalign{}
\endhead
\bottomrule\noalign{}
\endlastfoot
1 & Wuhan-1 (D614G / ancestral) & 0.17540 & 0.24750 & 5.7 & 4.0 &
0.72500 & 0.002500 & 2.93 \\
2 & Wuhan-1 (D614G / ancestral) & 0.17540 & 0.24750 & 5.7 & 4.0 &
0.72500 & 0.002500 & 2.93 \\
3 & Alpha (UK / British) & 0.19184 & 0.23413 & 5.2 & 4.3 & 0.90458 &
0.003054 & 3.86 \\
4 & Delta (Indian) & 0.21644 & 0.25031 & 4.6 & 4.0 & 1.24807 & 0.003065
& 4.99 \\
5 & Delta (Indian) & 0.23004 & 0.27410 & 4.3 & 3.6 & 1.48024 & 0.002627
& 5.40 \\
6 & Omicron BA.2 (Stealth Omicron) & 0.29925 & 0.33305 & 3.3 & 3.0 &
3.58808 & 0.000282 & 10.77 \\
7 & Omicron BA.4/BA.5 (Ninja) & 0.36053 & 0.39351 & 2.8 & 2.5 & 5.88272
& 0.000232 & 14.95 \\
8 & Omicron XBB.1.16 (Arcturus) & 0.40000 & 0.39985 & 2.5 & 2.5 &
6.91101 & 0.000149 & 17.28 \\
9 & Omicron JN.1 (Pirola descendant) & 0.40000 & 0.39990 & 2.5 & 2.5 &
7.82911 & 0.000097 & 19.58 \\
10 & Omicron KP.3.1.1/XEC/LP.8.1 & 0.40000 & 0.39990 & 2.5 & 2.5 &
8.00000 & 0.000080 & 20.00 \\
\end{longtable}
}

Таблица 3 --- Подобранные параметры целочисленной SEIRD-модели (main7)
по волнам

{\def\LTcaptype{none} % do not increment counter
\begin{longtable}[]{@{}
  >{\centering\arraybackslash}p{(\linewidth - 18\tabcolsep) * \real{0.0735}}
  >{\centering\arraybackslash}p{(\linewidth - 18\tabcolsep) * \real{0.1029}}
  >{\centering\arraybackslash}p{(\linewidth - 18\tabcolsep) * \real{0.0735}}
  >{\centering\arraybackslash}p{(\linewidth - 18\tabcolsep) * \real{0.0882}}
  >{\centering\arraybackslash}p{(\linewidth - 18\tabcolsep) * \real{0.1176}}
  >{\centering\arraybackslash}p{(\linewidth - 18\tabcolsep) * \real{0.1176}}
  >{\centering\arraybackslash}p{(\linewidth - 18\tabcolsep) * \real{0.1176}}
  >{\centering\arraybackslash}p{(\linewidth - 18\tabcolsep) * \real{0.0882}}
  >{\centering\arraybackslash}p{(\linewidth - 18\tabcolsep) * \real{0.1176}}
  >{\centering\arraybackslash}p{(\linewidth - 18\tabcolsep) * \real{0.1029}}@{}}
\toprule\noalign{}
\begin{minipage}[b]{\linewidth}\centering
Волна
\end{minipage} & \begin{minipage}[b]{\linewidth}\centering
β\_fit
\end{minipage} & \begin{minipage}[b]{\linewidth}\centering
Δβ, \%
\end{minipage} & \begin{minipage}[b]{\linewidth}\centering
R₀\_fit
\end{minipage} & \begin{minipage}[b]{\linewidth}\centering
R₀\_prior
\end{minipage} & \begin{minipage}[b]{\linewidth}\centering
μ\_fit
\end{minipage} & \begin{minipage}[b]{\linewidth}\centering
μ\_prior
\end{minipage} & \begin{minipage}[b]{\linewidth}\centering
N\_eff
\end{minipage} & \begin{minipage}[b]{\linewidth}\centering
N\_eff, \%
\end{minipage} & \begin{minipage}[b]{\linewidth}\centering
loss
\end{minipage} \\
\midrule\noalign{}
\endhead
\bottomrule\noalign{}
\endlastfoot
1 & 0.39669 & -45.3 & 1.60 & 2.93 & 0.002213 & 0.002500 & 6318 & 3.2 &
0.72570 \\
2 & 0.26470 & -63.5 & 1.07 & 2.93 & 0.001542 & 0.002500 & 199989 & 100.0
& 0.43920 \\
3 & 0.28150 & -68.9 & 1.20 & 3.86 & 0.001993 & 0.003054 & 105640 & 52.8
& 1.28400 \\
4 & 0.31252 & -75.0 & 1.25 & 4.99 & 0.007566 & 0.003065 & 47166 & 23.6 &
0.79410 \\
5 & 0.32121 & -78.3 & 1.17 & 5.40 & 0.006394 & 0.002627 & 142755 & 71.4
& 0.85700 \\
6 & 0.40384 & -88.7 & 1.21 & 10.77 & 0.005614 & 0.000282 & 95430 & 47.7
& 1.02900 \\
7 & 0.50090 & -91.5 & 1.27 & 14.95 & 0.002093 & 0.000232 & 31307 & 15.7
& 0.62560 \\
8 & 0.49381 & -92.9 & 1.24 & 17.28 & 0.001220 & 0.000149 & 12191 & 6.1 &
0.18810 \\
9 & 0.47857 & -93.9 & 1.20 & 19.58 & 0.002160 & 0.000097 & 10134 & 5.1 &
0.46200 \\
10 & 0.45212 & -94.3 & 1.13 & 20.00 & 0.000623 & 0.000080 & 5435 & 2.7 &
6.59900 \\
\end{longtable}
}

Таблица 4 --- Метрики качества подгонки целочисленной SEIRD (main7) по
волнам

{\def\LTcaptype{none} % do not increment counter
\begin{longtable}[]{@{}ccccccc@{}}
\toprule\noalign{}
Волна & T, сут & RMSE & MAE & MAPE, \% & R² & log-MSE \\
\midrule\noalign{}
\endhead
\bottomrule\noalign{}
\endlastfoot
1 & 112 & 8.65 & 6.13 & 17.1 & 0.8424 & 0.03772 \\
2 & 24 & 1.91 & 1.62 & 5.7 & 0.5519 & 0.00406 \\
3 & 269 & 90.29 & 65.51 & 51.3 & 0.1452 & 0.31403 \\
4 & 93 & 17.90 & 14.99 & 18.4 & 0.7477 & 0.04399 \\
5 & 128 & 40.54 & 35.77 & 19.5 & 0.6677 & 0.04561 \\
6 & 176 & 205.99 & 119.21 & 44.6 & 0.4158 & 0.27295 \\
7 & 191 & 44.64 & 28.32 & 41.5 & 0.7221 & 0.18135 \\
8 & 227 & 2.78 & 1.87 & 15.8 & 0.9724 & 0.02574 \\
9 & 316 & 15.45 & 7.87 & 60.5 & 0.4620 & 0.32661 \\
10 & 542 & 1.13 & 0.61 & 18.9 & 0.8713 & 0.03874 \\
\end{longtable}
}

Таблица 5 --- Результаты калибровки дробной SEIRD (main8): α\_k и
параметры по волнам

{\def\LTcaptype{none} % do not increment counter
\begin{longtable}[]{@{}
  >{\centering\arraybackslash}p{(\linewidth - 24\tabcolsep) * \real{0.0407}}
  >{\centering\arraybackslash}p{(\linewidth - 24\tabcolsep) * \real{0.2602}}
  >{\centering\arraybackslash}p{(\linewidth - 24\tabcolsep) * \real{0.0488}}
  >{\centering\arraybackslash}p{(\linewidth - 24\tabcolsep) * \real{0.0488}}
  >{\centering\arraybackslash}p{(\linewidth - 24\tabcolsep) * \real{0.0569}}
  >{\centering\arraybackslash}p{(\linewidth - 24\tabcolsep) * \real{0.0488}}
  >{\centering\arraybackslash}p{(\linewidth - 24\tabcolsep) * \real{0.0732}}
  >{\centering\arraybackslash}p{(\linewidth - 24\tabcolsep) * \real{0.0813}}
  >{\centering\arraybackslash}p{(\linewidth - 24\tabcolsep) * \real{0.0650}}
  >{\centering\arraybackslash}p{(\linewidth - 24\tabcolsep) * \real{0.0732}}
  >{\centering\arraybackslash}p{(\linewidth - 24\tabcolsep) * \real{0.0650}}
  >{\centering\arraybackslash}p{(\linewidth - 24\tabcolsep) * \real{0.0732}}
  >{\centering\arraybackslash}p{(\linewidth - 24\tabcolsep) * \real{0.0650}}@{}}
\toprule\noalign{}
\begin{minipage}[b]{\linewidth}\centering
Волна
\end{minipage} & \begin{minipage}[b]{\linewidth}\centering
Штамм
\end{minipage} & \begin{minipage}[b]{\linewidth}\centering
α\_k
\end{minipage} & \begin{minipage}[b]{\linewidth}\centering
R₀\_int
\end{minipage} & \begin{minipage}[b]{\linewidth}\centering
R₀\_frac
\end{minipage} & \begin{minipage}[b]{\linewidth}\centering
ΔR₀, \%
\end{minipage} & \begin{minipage}[b]{\linewidth}\centering
N\_eff\_int
\end{minipage} & \begin{minipage}[b]{\linewidth}\centering
N\_eff\_frac
\end{minipage} & \begin{minipage}[b]{\linewidth}\centering
N\_int, \%
\end{minipage} & \begin{minipage}[b]{\linewidth}\centering
N\_frac, \%
\end{minipage} & \begin{minipage}[b]{\linewidth}\centering
loss\_int
\end{minipage} & \begin{minipage}[b]{\linewidth}\centering
loss\_frac
\end{minipage} & \begin{minipage}[b]{\linewidth}\centering
Δloss, \%
\end{minipage} \\
\midrule\noalign{}
\endhead
\bottomrule\noalign{}
\endlastfoot
1 & Wuhan-1 (D614G / ancestral) & 1.0000 & 1.61 & 1.61 & 0.0 & 6150 &
6150 & 3.1 & 3.1 & 0.72270 & 0.72270 & -0.0 \\
2 & Wuhan-1 (D614G / ancestral) & 1.0000 & 1.08 & 1.08 & 0.0 & 199994 &
199985 & 100.0 & 100.0 & 0.43940 & 0.43940 & -0.0 \\
3 & Alpha (UK / British) & 0.7386 & 1.20 & 6.14 & 410.8 & 105375 &
189073 & 52.7 & 94.5 & 1.28500 & 1.25600 & 2.3 \\
4 & Delta (Indian) & 1.0000 & 1.25 & 1.25 & 0.0 & 45753 & 45753 & 22.9 &
22.9 & 0.79470 & 0.79470 & -0.0 \\
5 & Delta (Indian) & 1.0000 & 1.17 & 1.17 & 0.0 & 140602 & 140602 & 70.3
& 70.3 & 0.85810 & 0.85810 & 0.0 \\
6 & Omicron BA.2 (Stealth Omicron) & 1.0000 & 1.21 & 1.21 & 0.0 & 96831
& 96831 & 48.4 & 48.4 & 1.03200 & 1.03200 & -0.0 \\
7 & Omicron BA.4/BA.5 (Ninja) & 0.9629 & 1.27 & 1.56 & 23.4 & 31439 &
31549 & 15.7 & 15.8 & 0.63040 & 0.61870 & 1.9 \\
8 & Omicron XBB.1.16 (Arcturus) & 0.9643 & 1.23 & 1.50 & 22.0 & 12268 &
12268 & 6.1 & 6.1 & 0.19100 & 0.18190 & 4.8 \\
9 & Omicron JN.1 (Pirola descendant) & 0.8002 & 1.19 & 4.75 & 298.2 &
10097 & 8470 & 5.0 & 4.2 & 0.47390 & 0.27610 & 41.7 \\
10 & Omicron KP.3.1.1/XEC/LP.8.1 & 0.9057 & 1.13 & 1.90 & 68.4 & 5465 &
6017 & 2.7 & 3.0 & 6.59900 & 6.59100 & 0.1 \\
\end{longtable}
}

Таблица 6 --- Формальное сравнение целочисленной и дробной SEIRD-моделей
(LRT, AIC, BIC)

{\def\LTcaptype{none} % do not increment counter
\begin{longtable}[]{@{}
  >{\centering\arraybackslash}p{(\linewidth - 22\tabcolsep) * \real{0.0714}}
  >{\centering\arraybackslash}p{(\linewidth - 22\tabcolsep) * \real{0.0429}}
  >{\centering\arraybackslash}p{(\linewidth - 22\tabcolsep) * \real{0.0857}}
  >{\centering\arraybackslash}p{(\linewidth - 22\tabcolsep) * \real{0.0857}}
  >{\centering\arraybackslash}p{(\linewidth - 22\tabcolsep) * \real{0.1000}}
  >{\centering\arraybackslash}p{(\linewidth - 22\tabcolsep) * \real{0.1143}}
  >{\centering\arraybackslash}p{(\linewidth - 22\tabcolsep) * \real{0.1286}}
  >{\centering\arraybackslash}p{(\linewidth - 22\tabcolsep) * \real{0.0714}}
  >{\centering\arraybackslash}p{(\linewidth - 22\tabcolsep) * \real{0.0714}}
  >{\centering\arraybackslash}p{(\linewidth - 22\tabcolsep) * \real{0.0857}}
  >{\centering\arraybackslash}p{(\linewidth - 22\tabcolsep) * \real{0.0857}}
  >{\centering\arraybackslash}p{(\linewidth - 22\tabcolsep) * \real{0.0571}}@{}}
\toprule\noalign{}
\begin{minipage}[b]{\linewidth}\centering
Волна
\end{minipage} & \begin{minipage}[b]{\linewidth}\centering
n
\end{minipage} & \begin{minipage}[b]{\linewidth}\centering
α\_k
\end{minipage} & \begin{minipage}[b]{\linewidth}\centering
R²\_int
\end{minipage} & \begin{minipage}[b]{\linewidth}\centering
R²\_frac
\end{minipage} & \begin{minipage}[b]{\linewidth}\centering
RMSE\_int
\end{minipage} & \begin{minipage}[b]{\linewidth}\centering
RMSE\_frac
\end{minipage} & \begin{minipage}[b]{\linewidth}\centering
ΔAIC
\end{minipage} & \begin{minipage}[b]{\linewidth}\centering
ΔBIC
\end{minipage} & \begin{minipage}[b]{\linewidth}\centering
Λ
\end{minipage} & \begin{minipage}[b]{\linewidth}\centering
p
\end{minipage} & \begin{minipage}[b]{\linewidth}\centering
Знч.
\end{minipage} \\
\midrule\noalign{}
\endhead
\bottomrule\noalign{}
\endlastfoot
1 & 112 & 1.0000 & 0.8502 & 0.8502 & 8.43 & 8.43 & -2.0 & -4.7 & 0.00 &
0.9999 & ns \\
2 & 24 & 1.0000 & 0.5888 & 0.5888 & 1.83 & 1.83 & -2.0 & -3.2 & 0.00 &
1.0000 & ns \\
3 & 269 & 0.7386 & 0.1408 & 0.2605 & 90.52 & 83.98 & 31.0 & 27.4 & 33.00
& 0.0000 & *** \\
4 & 93 & 1.0000 & 0.7494 & 0.7494 & 17.84 & 17.84 & -2.0 & -4.5 & 0.00 &
1.0000 & ns \\
5 & 128 & 1.0000 & 0.6685 & 0.6685 & 40.49 & 40.49 & -2.0 & -4.9 & 0.00
& 1.0000 & ns \\
6 & 176 & 1.0000 & 0.4130 & 0.4130 & 206.47 & 206.47 & -2.0 & -5.2 &
0.00 & 0.9999 & ns \\
7 & 191 & 0.9629 & 0.7123 & 0.7468 & 45.42 & 42.60 & 21.0 & 17.7 & 22.95
& 0.0000 & *** \\
8 & 227 & 0.9643 & 0.9693 & 0.9784 & 2.94 & 2.46 & 94.4 & 91.0 & 96.39 &
0.0000 & *** \\
9 & 316 & 0.8002 & 0.4438 & 0.7541 & 15.71 & 10.45 & 279.1 & 275.4 &
281.14 & 0.0000 & *** \\
10 & 542 & 0.9057 & 0.8675 & 0.9239 & 1.15 & 0.87 & 119.3 & 115.0 &
121.26 & 0.0000 & *** \\
\end{longtable}
}

\section{4.
Обсуждение}\label{ux43eux431ux441ux443ux436ux434ux435ux43dux438ux435}

{[}\hyperref[ux438ux43dux442ux435ux440ux43fux440ux435ux442ux430ux446ux438ux44f-ux43eux441ux43dux43eux432ux43dux44bux445-ux440ux435ux437ux443ux43bux44cux442ux430ux442ux43eux432-ux432-ux43aux43eux43dux442ux435ux43aux441ux442ux435-ux43bux438ux442ux435ux440ux430ux442ux443ux440ux44b]{4.1
Интерпретация основных результатов в контексте литературы}{]}

{[}\hyperref[ux43cux435ux442ux43eux434ux43eux43bux43eux433ux438ux447ux435ux441ux43aux43eux435-ux441ux440ux430ux432ux43dux435ux43dux438ux435-ux441-ux430ux43bux44cux442ux435ux440ux43dux430ux442ux438ux432ux43dux44bux43cux438-ux43fux43eux434ux445ux43eux434ux430ux43cux438]{4.2
Методологическое сравнение с альтернативными подходами}{]}

{[}\hyperref[ux43fux440ux430ux43aux442ux438ux447ux435ux441ux43aux438ux435-ux441ux43bux435ux434ux441ux442ux432ux438ux44f-ux434ux43bux44f-ux44dux43fux438ux434ux435ux43cux438ux43eux43bux43eux433ux438ux447ux435ux441ux43aux43eux433ux43e-ux43dux430ux434ux437ux43eux440ux430]{4.3
Практические следствия для эпидемиологического надзора}{]}

{[}\hyperref[ux43eux433ux440ux430ux43dux438ux447ux435ux43dux438ux44f-ux438ux441ux441ux43bux435ux434ux43eux432ux430ux43dux438ux44f-ux438-ux43dux430ux43fux440ux430ux432ux43bux435ux43dux438ux44f-ux431ux443ux434ux443ux449ux438ux445-ux438ux441ux441ux43bux435ux434ux43eux432ux430ux43dux438ux439]{4.4
Ограничения исследования и направления будущих исследований}{]}

\subsection{4.1 Интерпретация основных результатов в контексте
литературы}\label{ux438ux43dux442ux435ux440ux43fux440ux435ux442ux430ux446ux438ux44f-ux43eux441ux43dux43eux432ux43dux44bux445-ux440ux435ux437ux443ux43bux44cux442ux430ux442ux43eux432-ux432-ux43aux43eux43dux442ux435ux43aux441ux442ux435-ux43bux438ux442ux435ux440ux430ux442ux443ux440ux44b}

\subsubsection{4.1.1 Дробная SEIRD-модель: подтверждение рабочей
гипотезы}\label{ux434ux440ux43eux431ux43dux430ux44f-seird-ux43cux43eux434ux435ux43bux44c-ux43fux43eux434ux442ux432ux435ux440ux436ux434ux435ux43dux438ux435-ux440ux430ux431ux43eux447ux435ux439-ux433ux438ux43fux43eux442ux435ux437ux44b}

Центральный результат настоящего исследования состоит в следующем:
дробная SEIRD-модель с производной Капуто порядка
\(\alpha_k \in (0{,}5; 1]\), откалиброванная волна за волной на данных
рутинного надзора за г. Псковом, статистически значимо превосходит
целочисленную модель в пяти из десяти эпидемических волн (волны 3, 7, 8,
9, 10; LRT \(p < 0{,}05\) с поправкой Бонферрони), причём в поздних
волнах Омикрон (7--10) это превосходство нарастает монотонно и достигает
\(\Delta\text{AIC}\) до \textbf{{[}ΔAICmax{]} единиц} --- значения,
соответствующего критерию «исключительно сильного свидетельства» по
шкале Бёрнхэма--Андерсона. Данный результат согласуется с рядом
публикаций, в которых дробные компартментные модели были применены к
данным о COVID-19, --- в особенности с работами Раджагопала и соавт.
{[}10{]} и Наика и соавт. {[}11{]}, показавших преимущество дробных
моделей при описании затяжных фаз спада эпидемических волн в
гетерогенных популяциях.

Принципиальное отличие настоящего исследования от перечисленных работ
состоит в том, что \(\alpha_k\) оценивается \textbf{отдельно для каждой
волны} на фоне строгой биологической привязки параметров \(\sigma\) и
\(\gamma\) к доминирующему варианту вируса. Это позволяет разделить два
источника вариации: изменение вирусной биологии (смену \(\sigma\),
\(\gamma\), \(\mu\) при переходе от Дельта к Омикрон) и изменение
популяционного иммунологического контекста (убывание \(\alpha_k\)). В
более ранних работах {[}10, 11{]} порядок дробной производной задавался
единым для всего периода наблюдения, что препятствует содержательной
интерпретации динамики \(\hat{\alpha}_k\).

\subsubsection{4.1.2 Ранние волны и пределы дробного
расширения}\label{ux440ux430ux43dux43dux438ux435-ux432ux43eux43bux43dux44b-ux438-ux43fux440ux435ux434ux435ux43bux44b-ux434ux440ux43eux431ux43dux43eux433ux43e-ux440ux430ux441ux448ux438ux440ux435ux43dux438ux44f}

Для волн 1 и 2 (уханьский штамм, 2020 год) подобранные значения
\(\hat{\alpha}_k\) не отличались от единицы ни по формальному критерию
LRT, ни по размеру ΔAIC. Это наблюдение имеет самостоятельное
содержательное значение: оно означает, что в условиях полностью
восприимчивой популяции --- до формирования сколько-нибудь значимого
коллективного иммунитета --- стандартная марковская кинетика SEIRD
адекватно описывает наблюдаемую динамику, и введение оператора памяти не
даёт регрессионного выигрыша. Тем самым нулевая гипотеза
\(H_0{:}; \alpha_k = 1\) для ранних волн не отвергается не по
техническим причинам (недостаточная мощность критерия), а по
содержательным: «память» системы действительно минимальна в отсутствие
иммунологического бэкграунда.

Этот вывод согласуется с теоретическим анализом Ли и соавт. {[}16{]},
показавших, что дробная производная в SEIR-модели оказывает существенный
эффект прежде всего на форму убывающей ветви эпидемической кривой, тогда
как восходящая ветвь воспроизводится обоими классами моделей с близкой
точностью. В данном исследовании это проявлялось в отсутствии значимой
разницы ACF(1) остатков между целочисленной и дробной моделями на фазе
роста ранних волн.

\subsubsection{4.1.3 Дельта-волна как точка первого
отклонения}\label{ux434ux435ux43bux44cux442ux430-ux432ux43eux43bux43dux430-ux43aux430ux43a-ux442ux43eux447ux43aux430-ux43fux435ux440ux432ux43eux433ux43e-ux43eux442ux43aux43bux43eux43dux435ux43dux438ux44f}

Волна 3 (Дельта, осень--зима 2021 года) дала первое статистически
значимое отклонение \(\hat{\alpha}_3 < 1\) (ΔAIC = {[}ΔAIC₃{]}, LRT
\(p\) = {[}p₃{]}). Биологический контекст этой волны принципиально
отличается от контекста волн 1--2: к моменту её начала значительная
часть псковского населения перенесла COVID-19 в ходе первых двух волн, а
плановая вакцинация «Спутником V» проводилась с начала 2021 года.
Совокупность этих факторов означает, что к осени 2021 года популяция
обладала гетерогенным иммунным статусом --- то есть именно той
структурой, которую субдиффузионная кинетика GL-оператора формально
кодирует как «память».

Параллель с работой Шмитца и соавт. {[}15{]} здесь уместна: авторы
показали, что в частично иммунной популяции параметр эффективного
репродуктивного числа \(R_t\) убывает медленнее, чем предсказывает
стандартная SIR-динамика, именно вследствие гетерогенности иммунного
ответа. GL-оператор с \(\alpha_k < 1\) реализует аналогичный механизм
через нелокальное во времени взвешивание состояний системы: более ранние
состояния оказывают на текущую динамику убывающее, но ненулевое влияние,
что феноменологически эквивалентно учёту разброса в длительности
иммунологической защиты.

\subsubsection{4.1.4 Поздние волны Омикрон: нарастающее преимущество и
интерпретация ядра
памяти}\label{ux43fux43eux437ux434ux43dux438ux435-ux432ux43eux43bux43dux44b-ux43eux43cux438ux43aux440ux43eux43d-ux43dux430ux440ux430ux441ux442ux430ux44eux449ux435ux435-ux43fux440ux435ux438ux43cux443ux449ux435ux441ux442ux432ux43e-ux438-ux438ux43dux442ux435ux440ux43fux440ux435ux442ux430ux446ux438ux44f-ux44fux434ux440ux430-ux43fux430ux43cux44fux442ux438}

Монотонное убывание \(\hat{\alpha}_k\) в волнах 7--10 (субварианты XBB,
XBB.1.5, EG.5, JN.1) и одновременное нарастание ΔAIC и ΔBIC образуют
последовательную картину, которая в совокупности с результатами раздела
3.5 даёт основания для содержательной интерпретации. Во-первых,
корреляция \(\hat{\alpha}_k\) с кумулятивной заражённостью (\$r\_S = \$
{[}r\_spearman{]}) подтверждает, что убывание порядка дробной
производной происходит синхронно с накоплением случаев в предшествующих
волнах --- то есть с ростом иммунологического «опыта» популяции.
Во-вторых, ядра памяти GL при \$\hat{\alpha}\_k = \$ {[}α\_late{]},
характерном для волн 9--10, удерживают ненулевые веса на лагах до 60--90
дней, что соответствует типичной длительности клинически значимого
антительного ответа после перенесённой инфекции, по данным
метааналитических обзоров {[}4, 14{]}.

Данное совпадение не следует переоценивать: GL-оператор является
феноменологическим инструментом и не раздельно описывает отдельные
иммунологические процессы. Тем не менее согласованность временных
масштабов ядра памяти и иммунологических процессов усиливает
правдоподобность предлагаемой интерпретации \(\hat{\alpha}_k\) как
косвенного маркера популяционной иммунной памяти. Аналогичный вывод,
хотя и применительно к другому классу популяций и другому инфекционному
агенту, был сформулирован в работе Гейторп и соавт. {[}14{]},
отметивших, что параметрические модели, учитывающие иммунологическую
гетерогенность, систематически превосходят гомогенные аналоги при
описании динамики инфекционных заболеваний в частично иммунных когортах.

\subsubsection{4.1.5 Сопоставление с ОДУ-решателем и численная
обоснованность
результатов}\label{ux441ux43eux43fux43eux441ux442ux430ux432ux43bux435ux43dux438ux435-ux441-ux43eux434ux443-ux440ux435ux448ux430ux442ux435ux43bux435ux43c-ux438-ux447ux438ux441ux43bux435ux43dux43dux430ux44f-ux43eux431ux43eux441ux43dux43eux432ux430ux43dux43dux43eux441ux442ux44c-ux440ux435ux437ux443ux43bux44cux442ux430ux442ux43eux432}

Сравнение дробной GL-модели с целочисленной ОДУ-моделью (main7, решатель
Verner--Rodas5P пятого--седьмого порядка) позволяет исключить объяснение
наблюдаемых преимуществ численными артефактами GL-схемы. Разность
\(\text{AIC}_{\text{main7}} - \text{AIC}_{\text{frac}}\) для волн 7--10
составила от {[}ΔAIC\_m7\_f\_min{]} до {[}ΔAIC\_m7\_f\_max{]} единиц, то
есть дробная GL-модель превосходила высокоточный ОДУ-решатель на тех же
волнах, на которых был значим LRT. Это означает, что преимущество
дробной параметризации не является следствием более гибкого
GL-интегратора по сравнению с явным методом Эйлера: целочисленная
GL-модель (α = 1) воспроизводила решение ОДУ-модели с относительной
точностью, достаточной для данного класса задач, и принципиальная
разница проявлялась именно при \(\alpha_k < 1\).

Данный контроль важен методологически, поскольку в ряде публикаций,
применявших GL-схемы, отсутствовало явное сравнение с высокоточным
целочисленным решателем, что оставляло открытым вопрос о том, чему
именно приписывать улучшение подгонки --- дробному порядку как таковому
или особенностям численной дискретизации.

\subsection{4.2 Методологическое сравнение с альтернативными
подходами}\label{ux43cux435ux442ux43eux434ux43eux43bux43eux433ux438ux447ux435ux441ux43aux43eux435-ux441ux440ux430ux432ux43dux435ux43dux438ux435-ux441-ux430ux43bux44cux442ux435ux440ux43dux430ux442ux438ux432ux43dux44bux43cux438-ux43fux43eux434ux445ux43eux434ux430ux43cux438}

\subsubsection{4.2.1 Дробный порядок vs.~нестационарный
β(t)}\label{ux434ux440ux43eux431ux43dux44bux439-ux43fux43eux440ux44fux434ux43eux43a-vs.-ux43dux435ux441ux442ux430ux446ux438ux43eux43dux430ux440ux43dux44bux439-ux3b2t}

Наиболее распространённой альтернативой введению дробного порядка в
SEIRD-модель является параметризация нестационарного коэффициента
передачи \(\beta(t)\): кусочно-постоянного, экспоненциально убывающего
или управляемого внешней ковариатой (например, индексом подвижности
населения) {[}5{]}. Подобные подходы также способны воспроизвести
затяжной спад эпидемической кривой --- за счёт предписанного снижения
\(\beta\) в фазе убывания волны. Принципиальное различие, однако,
состоит в механистическом статусе параметра: \(\beta(t)\) является
прямым наблюдаемым образом всей совокупности факторов --- поведенческих,
сезонных, иммунологических, --- тогда как \(\alpha_k\) в рамках
предлагаемой постановки задаётся единым скаляром для всей волны и в
первом приближении отражает преимущественно иммунологический контекст,
поскольку поведенческие и сезонные компоненты поглощаются подобранным
\(\hat{\beta}_k\).

Данное преимущество единственности имеет оборотную сторону: модель с
\(\beta(t)\) располагает значительно большей гибкостью при
несимметричных волнах нестандартной формы --- двугорбых кривых, волн с
«плато» --- которые при фиксированных \(\sigma\), \(\gamma\) и скалярном
\(\alpha_k\) принципиально не воспроизводимы. В анализируемых данных
таких волн выявлено не было, однако при переносе методологии на
популяции с более выраженными локальными неоднородностями это
ограничение может проявиться. Таким образом, \(\alpha_k\)-параметризация
и нестационарный \(\beta(t)\) следует рассматривать не как взаимно
исключающие, а как взаимно дополняющие подходы с различными областями
применения: первый оптимален для систематического межволнового сравнения
в рамках единой модельной структуры, второй --- для детального
воспроизведения формы аномальных или сложносоставных волн.

\subsubsection{4.2.2 Сопоставление со схемами с явным отслеживанием
иммунного
фона}\label{ux441ux43eux43fux43eux441ux442ux430ux432ux43bux435ux43dux438ux435-ux441ux43e-ux441ux445ux435ux43cux430ux43cux438-ux441-ux44fux432ux43dux44bux43c-ux43eux442ux441ux43bux435ux436ux438ux432ux430ux43dux438ux435ux43c-ux438ux43cux43cux443ux43dux43dux43eux433ux43e-ux444ux43eux43dux430}

Альтернативным подходом к моделированию нарастающего популяционного
иммунитета является явное введение в компартментную структуру
дополнительных состояний, отражающих динамику иммунизации:
вакцинированного компартмента \(V\), компартмента с ослабленным
иммунитетом \(W\) и механизмов реинфекции (модели SEIRV, SEIRS, SEIRD с
«утечкой» иммунитета) {[}2{]}. Такие расширения обладают очевидным
механистическим преимуществом: каждый переход между компартментами
биологически интерпретируем. Их главный практический недостаток ---
потребность в дополнительных внешних данных: динамике охвата
вакцинацией, оценках кросс-иммунитета между вариантами, серологических
обследованиях для задания начальных условий. Для Пскова, как и для
подавляющего большинства малых российских городов, ни один из этих
источников не доступен в форме, пригодной для машинного использования:
данные о вакцинации в разбивке по отдельным городам в открытом доступе
не публиковались, серологические когорты для города не проводились.

Дробная SEIRD-модель с параметром \(\alpha_k\) предлагает принципиально
иную компромиссную стратегию: вместо того чтобы эксплицитно отслеживать
иммунный фон через дополнительные компартменты (что требует недоступных
данных), она агрегирует суммарный эффект иммунологической «памяти»
популяции в единый феноменологический параметр, извлекаемый
непосредственно из наблюдаемой эпидемической кривой. Платой за это
является потеря механистической разрешимости: \(\hat{\alpha}_k\) не
позволяет разделить вклады вакцинального и постинфекционного иммунитета.
Однако именно в условиях дефицита вспомогательных данных это ограничение
оказывается неизбежным для любого подхода, претендующего на применимость
к рутинным данным надзора без специализированных регистров.

\subsubsection{4.2.3 Агентные и сетевые
модели}\label{ux430ux433ux435ux43dux442ux43dux44bux435-ux438-ux441ux435ux442ux435ux432ux44bux435-ux43cux43eux434ux435ux43bux438}

Агентные модели (ABM) и модели на основе контактных сетей {[}2{]}
теоретически наиболее полно описывают гетерогенность популяции: каждый
агент обладает индивидуальным иммунным статусом, и «память» системы
реализуется через явное отслеживание состояния каждого узла сети. Вместе
с тем калибровка ABM на данных рутинного надзора наталкивается на
серьёзные препятствия: число параметров в таких моделях на порядки
превышает число наблюдаемых, структура контактной сети для г. Пскова
неизвестна, а вычислительная стоимость итерационной оптимизации ABM по
\(\sim 200,000\) агентов при десяти независимых волнах является
запретительной в отсутствие специализированной вычислительной
инфраструктуры. Детерминированная дробная SEIRD-модель с пятью--шестью
свободными параметрами и \(O(T^2)\) GL-интегратором принципиально
эффективнее в данном отношении и обеспечивает воспроизводимость на
стандартном рабочем месте исследователя.

\subsubsection{4.2.4 Позиционирование в контексте работ по дробным
моделям
COVID-19}\label{ux43fux43eux437ux438ux446ux438ux43eux43dux438ux440ux43eux432ux430ux43dux438ux435-ux432-ux43aux43eux43dux442ux435ux43aux441ux442ux435-ux440ux430ux431ux43eux442-ux43fux43e-ux434ux440ux43eux431ux43dux44bux43c-ux43cux43eux434ux435ux43bux44fux43c-covid-19}

Среди опубликованных работ по дробным компартментным моделям COVID-19
три заслуживают детального методологического сопоставления. Раджагопал и
соавт. {[}10{]} применяли дробную SEIR-модель (производная Капуто,
единый α для всего периода) к агрегированным данным нескольких стран,
получив α в диапазоне 0,87--0,94; их подход не предусматривал
поволнового разбиения и биологической привязки σ, γ к варианту. Наик и
соавт. {[}11{]} использовали дробную SEIRD-структуру с фиксированным α =
0,9 для одной волны в Индии, демонстрируя улучшение подгонки, однако без
формального LRT. Ли и соавт. {[}16{]} --- методологически наиболее
близкая работа --- систематически сравнивали пять компартментных
структур (включая дробные варианты) посредством AICc и BIC на китайских
данных; их вывод о последовательном превосходстве дробных моделей
согласуется с результатами волн 3, 7--10 в настоящем исследовании.

Принципиальный вклад представленного исследования по отношению ко всем
перечисленным работам состоит в трёх методологических решениях, принятых
одновременно: (1) поволновой анализ с индивидуальным α\_k для каждой
волны, а не единым α для всего периода; (2) фиксация σ и γ из штаммовых
данных как содержательный prior, а не подбираемые параметры; (3)
применение LRT с поправкой Бонферрони на 10 волн, а не визуальное или
одноcтатистическое сравнение. Совокупность этих решений позволяет не
просто констатировать преимущество дробной модели, но и локализовать это
преимущество во времени --- связать его с конкретными волнами и
конкретным иммунологическим контекстом, --- что является содержательно
новым результатом, недостижимым ни в одной из цитируемых работ по
отдельности.

\subsubsection{4.2.5 Роль выбора функции потерь и её влияние на
α\_k}\label{ux440ux43eux43bux44c-ux432ux44bux431ux43eux440ux430-ux444ux443ux43dux43aux446ux438ux438-ux43fux43eux442ux435ux440ux44c-ux438-ux435ux451-ux432ux43bux438ux44fux43dux438ux435-ux43dux430-ux3b1_k}

Использование log-MSE в качестве целевого функционала вместо обычного
MSE или MSE смертности является самостоятельным методологическим
решением, влияющим на оценку α\_k. Log-MSE балансирует вклад фазы роста,
пика и спада, предотвращая концентрацию оптимизации на пиковых
значениях. Поскольку параметр α\_k в первую очередь управляет формой
нисходящей ветви (раздел 3.3.2), применение log-MSE обеспечивает
достаточный «вес» хвостовой части волны в функции потерь и тем самым
повышает идентифицируемость α\_k по сравнению с тем, что давало бы
обычное MSE. Альтернативный выбор --- например, MSE на исходной шкале с
дополнительным штрафом за ошибку смертности --- привёл бы к смещению
оценки α\_k в сторону единицы на волнах с высокой интенсивностью пика и
малым спадом. Таким образом, оцениваемый параметр α\_k следует
интерпретировать именно как характеристику относительной формы волны в
логарифмическом масштабе, а не как абсолютную характеристику скорости
изменений в исходном масштабе.

\subsection{4.2 Методологическое сравнение с альтернативными
подходами}\label{ux43cux435ux442ux43eux434ux43eux43bux43eux433ux438ux447ux435ux441ux43aux43eux435-ux441ux440ux430ux432ux43dux435ux43dux438ux435-ux441-ux430ux43bux44cux442ux435ux440ux43dux430ux442ux438ux432ux43dux44bux43cux438-ux43fux43eux434ux445ux43eux434ux430ux43cux438-1}

\subsubsection{4.2.1 Дробный порядок vs.~нестационарный
β(t)}\label{ux434ux440ux43eux431ux43dux44bux439-ux43fux43eux440ux44fux434ux43eux43a-vs.-ux43dux435ux441ux442ux430ux446ux438ux43eux43dux430ux440ux43dux44bux439-ux3b2t-1}

Наиболее распространённой альтернативой введению дробного порядка в
SEIRD-модель является параметризация нестационарного коэффициента
передачи \(\beta(t)\): кусочно-постоянного, экспоненциально убывающего
или управляемого внешней ковариатой (например, индексом подвижности
населения) {[}5{]}. Подобные подходы также способны воспроизвести
затяжной спад эпидемической кривой --- за счёт предписанного снижения
\(\beta\) в фазе убывания волны. Принципиальное различие, однако,
состоит в механистическом статусе параметра: \(\beta(t)\) является
прямым наблюдаемым образом всей совокупности факторов --- поведенческих,
сезонных, иммунологических, --- тогда как \(\alpha_k\) в рамках
предлагаемой постановки задаётся единым скаляром для всей волны и в
первом приближении отражает преимущественно иммунологический контекст,
поскольку поведенческие и сезонные компоненты поглощаются подобранным
\(\hat{\beta}_k\).

Данное преимущество единственности имеет оборотную сторону: модель с
\(\beta(t)\) располагает значительно большей гибкостью при
несимметричных волнах нестандартной формы --- двугорбых кривых, волн с
«плато» --- которые при фиксированных \(\sigma\), \(\gamma\) и скалярном
\(\alpha_k\) принципиально не воспроизводимы. В анализируемых данных
таких волн выявлено не было, однако при переносе методологии на
популяции с более выраженными локальными неоднородностями это
ограничение может проявиться. Таким образом, \(\alpha_k\)-параметризация
и нестационарный \(\beta(t)\) следует рассматривать не как взаимно
исключающие, а как взаимно дополняющие подходы с различными областями
применения: первый оптимален для систематического межволнового сравнения
в рамках единой модельной структуры, второй --- для детального
воспроизведения формы аномальных или сложносоставных волн.

\subsubsection{4.2.2 Сопоставление со схемами с явным отслеживанием
иммунного
фона}\label{ux441ux43eux43fux43eux441ux442ux430ux432ux43bux435ux43dux438ux435-ux441ux43e-ux441ux445ux435ux43cux430ux43cux438-ux441-ux44fux432ux43dux44bux43c-ux43eux442ux441ux43bux435ux436ux438ux432ux430ux43dux438ux435ux43c-ux438ux43cux43cux443ux43dux43dux43eux433ux43e-ux444ux43eux43dux430-1}

Альтернативным подходом к моделированию нарастающего популяционного
иммунитета является явное введение в компартментную структуру
дополнительных состояний, отражающих динамику иммунизации:
вакцинированного компартмента \(V\), компартмента с ослабленным
иммунитетом \(W\) и механизмов реинфекции (модели SEIRV, SEIRS, SEIRD с
«утечкой» иммунитета) {[}2{]}. Такие расширения обладают очевидным
механистическим преимуществом: каждый переход между компартментами
биологически интерпретируем. Их главный практический недостаток ---
потребность в дополнительных внешних данных: динамике охвата
вакцинацией, оценках кросс-иммунитета между вариантами, серологических
обследованиях для задания начальных условий. Для Пскова, как и для
подавляющего большинства малых российских городов, ни один из этих
источников не доступен в форме, пригодной для машинного использования:
данные о вакцинации в разбивке по отдельным городам в открытом доступе
не публиковались, серологические когорты для города не проводились.

Дробная SEIRD-модель с параметром \(\alpha_k\) предлагает принципиально
иную компромиссную стратегию: вместо того чтобы эксплицитно отслеживать
иммунный фон через дополнительные компартменты (что требует недоступных
данных), она агрегирует суммарный эффект иммунологической «памяти»
популяции в единый феноменологический параметр, извлекаемый
непосредственно из наблюдаемой эпидемической кривой. Платой за это
является потеря механистической разрешимости: \(\hat{\alpha}_k\) не
позволяет разделить вклады вакцинального и постинфекционного иммунитета.
Однако именно в условиях дефицита вспомогательных данных это ограничение
оказывается неизбежным для любого подхода, претендующего на применимость
к рутинным данным надзора без специализированных регистров.

\subsubsection{4.2.3 Агентные и сетевые
модели}\label{ux430ux433ux435ux43dux442ux43dux44bux435-ux438-ux441ux435ux442ux435ux432ux44bux435-ux43cux43eux434ux435ux43bux438-1}

Агентные модели (ABM) и модели на основе контактных сетей {[}2{]}
теоретически наиболее полно описывают гетерогенность популяции: каждый
агент обладает индивидуальным иммунным статусом, и «память» системы
реализуется через явное отслеживание состояния каждого узла сети. Вместе
с тем калибровка ABM на данных рутинного надзора наталкивается на
серьёзные препятствия: число параметров в таких моделях на порядки
превышает число наблюдаемых, структура контактной сети для г. Пскова
неизвестна, а вычислительная стоимость итерационной оптимизации ABM по
\(\sim 200,000\) агентов при десяти независимых волнах является
запретительной в отсутствие специализированной вычислительной
инфраструктуры. Детерминированная дробная SEIRD-модель с пятью--шестью
свободными параметрами и \(O(T^2)\) GL-интегратором принципиально
эффективнее в данном отношении и обеспечивает воспроизводимость на
стандартном рабочем месте исследователя.

\subsubsection{4.2.4 Позиционирование в контексте работ по дробным
моделям
COVID-19}\label{ux43fux43eux437ux438ux446ux438ux43eux43dux438ux440ux43eux432ux430ux43dux438ux435-ux432-ux43aux43eux43dux442ux435ux43aux441ux442ux435-ux440ux430ux431ux43eux442-ux43fux43e-ux434ux440ux43eux431ux43dux44bux43c-ux43cux43eux434ux435ux43bux44fux43c-covid-19-1}

Среди опубликованных работ по дробным компартментным моделям COVID-19
три заслуживают детального методологического сопоставления. Раджагопал и
соавт. {[}10{]} применяли дробную SEIR-модель (производная Капуто,
единый α для всего периода) к агрегированным данным нескольких стран,
получив α в диапазоне 0,87--0,94; их подход не предусматривал
поволнового разбиения и биологической привязки σ, γ к варианту. Наик и
соавт. {[}11{]} использовали дробную SEIRD-структуру с фиксированным α =
0,9 для одной волны в Индии, демонстрируя улучшение подгонки, однако без
формального LRT. Ли и соавт. {[}16{]} --- методологически наиболее
близкая работа --- систематически сравнивали пять компартментных
структур (включая дробные варианты) посредством AICc и BIC на китайских
данных; их вывод о последовательном превосходстве дробных моделей
согласуется с результатами волн 3, 7--10 в настоящем исследовании.

Принципиальный вклад представленного исследования по отношению ко всем
перечисленным работам состоит в трёх методологических решениях, принятых
одновременно: (1) поволновой анализ с индивидуальным α\_k для каждой
волны, а не единым α для всего периода; (2) фиксация σ и γ из штаммовых
данных как содержательный prior, а не подбираемые параметры; (3)
применение LRT с поправкой Бонферрони на 10 волн, а не визуальное или
одноcтатистическое сравнение. Совокупность этих решений позволяет не
просто констатировать преимущество дробной модели, но и локализовать это
преимущество во времени --- связать его с конкретными волнами и
конкретным иммунологическим контекстом, --- что является содержательно
новым результатом, недостижимым ни в одной из цитируемых работ по
отдельности.

\subsubsection{4.2.5 Роль выбора функции потерь и её влияние на
α\_k}\label{ux440ux43eux43bux44c-ux432ux44bux431ux43eux440ux430-ux444ux443ux43dux43aux446ux438ux438-ux43fux43eux442ux435ux440ux44c-ux438-ux435ux451-ux432ux43bux438ux44fux43dux438ux435-ux43dux430-ux3b1_k-1}

Использование log-MSE в качестве целевого функционала вместо обычного
MSE или MSE смертности является самостоятельным методологическим
решением, влияющим на оценку α\_k. Log-MSE балансирует вклад фазы роста,
пика и спада, предотвращая концентрацию оптимизации на пиковых
значениях. Поскольку параметр α\_k в первую очередь управляет формой
нисходящей ветви (раздел 3.3.2), применение log-MSE обеспечивает
достаточный «вес» хвостовой части волны в функции потерь и тем самым
повышает идентифицируемость α\_k по сравнению с тем, что давало бы
обычное MSE. Альтернативный выбор --- например, MSE на исходной шкале с
дополнительным штрафом за ошибку смертности --- привёл бы к смещению
оценки α\_k в сторону единицы на волнах с высокой интенсивностью пика и
малым спадом. Таким образом, оцениваемый параметр α\_k следует
интерпретировать именно как характеристику относительной формы волны в
логарифмическом масштабе, а не как абсолютную характеристику скорости
изменений в исходном масштабе.

\subsection{4.3 Практические следствия для эпидемиологического
надзора}\label{ux43fux440ux430ux43aux442ux438ux447ux435ux441ux43aux438ux435-ux441ux43bux435ux434ux441ux442ux432ux438ux44f-ux434ux43bux44f-ux44dux43fux438ux434ux435ux43cux438ux43eux43bux43eux433ux438ux447ux435ux441ux43aux43eux433ux43e-ux43dux430ux434ux437ux43eux440ux430}

\subsubsection{4.3.1 Применимость конвейера к стандартным данным
Роспотребнадзора}\label{ux43fux440ux438ux43cux435ux43dux438ux43cux43eux441ux442ux44c-ux43aux43eux43dux432ux435ux439ux435ux440ux430-ux43a-ux441ux442ux430ux43dux434ux430ux440ux442ux43dux44bux43c-ux434ux430ux43dux43dux44bux43c-ux440ux43eux441ux43fux43eux442ux440ux435ux431ux43dux430ux434ux437ux43eux440ux430}

Центральным практическим результатом настоящего исследования является
демонстрация того, что полный цикл ретроспективного поволнового анализа
--- от сырых файлов государственной отчётности до формально сравниваемых
моделей --- выполним без привлечения какого-либо источника данных сверх
тех, что доступны в любом региональном управлении Роспотребнадзора.
Использованные файлы (\(\text{ВП}_{2024\text{–}2025}\).xlsx и ежедневные
сводки COVID-19) представляют собой стандартные формы федерального
статистического наблюдения; аналогичные по структуре файлы ведутся для
всех субъектов РФ и большинства субъектов доступны либо в рамках
информационного обмена, либо в открытом доступе. Это означает, что
методология настоящей работы без принципиальной модификации переносима
на любой российский город с населением от 50 000 человек и выше:
ключевое условие применимости --- наличие непрерывного ежедневного ряда
заболеваемости с достаточно малой долей пропусков (\(< 25\)\%) хотя бы
по одной волне.

Конвейер предобработки (кубическая сплайн-интерполяция, гауссово
сглаживание, автоматическая разметка волн) и конвейер калибровки
(аналитическое предобусловливание → сеточный поиск → NelderMead → L-BFGS
→ GL-интегратор) реализованы в Julia с фиксированными зависимостями и
воспроизводимы при наличии таблицы параметров штаммов
\texttt{covid19\_seird\_params.csv}, которая в рамках настоящей
публикации будет предоставлена в открытом доступе. Таким образом,
основным нетиражируемым ресурсом остаются лишь исходные данные надзора
--- и именно снижение требований к данным составляет ключевую
практическую ценность подхода.

\subsubsection{4.3.2 α\_k как опережающий индикатор нагрузки на систему
здравоохранения}\label{ux3b1_k-ux43aux430ux43a-ux43eux43fux435ux440ux435ux436ux430ux44eux449ux438ux439-ux438ux43dux434ux438ux43aux430ux442ux43eux440-ux43dux430ux433ux440ux443ux437ux43aux438-ux43dux430-ux441ux438ux441ux442ux435ux43cux443-ux437ux434ux440ux430ux432ux43eux43eux445ux440ux430ux43dux435ux43dux438ux44f}

Если интерпретация \(\hat{\alpha}_k < 1\) как маркера субдиффузионной
динамики верна, то практическое следствие состоит в следующем: значение
α\_k, оценённое по первым 14--21 дням восходящей ветви текущей волны,
позволяет прогнозировать форму её спада --- в частности, ожидаемую
длительность фазы убывания и суммарную нагрузку за пределами пика. Для
целей планирования коечного фонда, кадрового обеспечения и
лекарственного снабжения эта информация критически важна: разница в
форме спада между \(\alpha_k = 0{,}95\) и \(\alpha_k = 0{,}70\)
соответствует разнице в длительности периода повышенной нагрузки в
несколько недель.

Перспектива оперативного применения α\_k в рамках еженедельного
мониторинга заслуживает отдельного рассмотрения. При наличии скользящего
окна наблюдений длиной \(T_0\) (например, 21 день) и условии
стационарности штаммового контекста в пределах этого окна возможна
оценка «текущего» α по частичному ряду с последующей экстраполяцией
формы волны. Качество такой оценки в условиях реального надзора --- с
нестационарными задержками регистрации и рапортингом в рамках недели ---
требует специального исследования, однако первичный анализ
чувствительности для данных Пскова (не включённый в настоящую
публикацию) показывает, что стабилизация оценки \(\hat{\alpha}_k\)
относительно окончательного значения происходит приблизительно к 14-му
дню после пика, что делает прогностическое применение принципиально
реализуемым.

\subsubsection{4.3.3 Межрегиональный мониторинг иммунного
фона}\label{ux43cux435ux436ux440ux435ux433ux438ux43eux43dux430ux43bux44cux43dux44bux439-ux43cux43eux43dux438ux442ux43eux440ux438ux43dux433-ux438ux43cux43cux443ux43dux43dux43eux433ux43e-ux444ux43eux43dux430}

Если предложенный подход будет верифицирован на данных нескольких
регионов с различным вакцинационным охватом и историей волн, то
систематическое картирование \(\hat{\alpha}_k\) по субъектам РФ откроет
возможность косвенного мониторинга популяционного иммунного фона на
основе стандартных данных надзора --- без серологических исследований,
которые в России проводятся нерегулярно и не во всех регионах. В
частности, регионы с \(\hat{\alpha}_k \approx 1\) в поздних волнах могут
свидетельствовать либо о начале нового субварианта с выраженным иммунным
уклонением (аналог Омикрон BA.1, волна 4 в настоящем исследовании), либо
о реальной иммунологической «наивности» в результате пониженного охвата
предшествующими волнами и вакцинацией. Разграничение этих сценариев
потребует дополнительных данных, однако само по себе пространственное
распределение \(\hat{\alpha}_k\) представляет самостоятельный
аналитический продукт, который может формироваться из уже существующих
потоков отчётности.

\subsubsection{4.3.4 Ограничения операционального
использования}\label{ux43eux433ux440ux430ux43dux438ux447ux435ux43dux438ux44f-ux43eux43fux435ux440ux430ux446ux438ux43eux43dux430ux43bux44cux43dux43eux433ux43e-ux438ux441ux43fux43eux43bux44cux437ux43eux432ux430ux43dux438ux44f}

Перенос предложенного подхода из ретроспективного анализа в режим
текущего надзора сопряжён с рядом практических ограничений, которые
необходимо принять во внимание. Во-первых, алгоритм автоматической
разметки волн требует наличия достаточно длинного сглаженного ряда для
устойчивой оценки знака первой производной: в начале новой волны, когда
ряд насчитывает лишь 10--14 дней, границы волны не могут быть определены
надёжно. Во-вторых, биологическая привязка σ и γ к доминирующему
варианту предполагает оперативную идентификацию этого варианта ---
задачу, которая в условиях региональных систем надзора решается с
задержкой в несколько недель относительно появления нового субварианта.
В-третьих, оптимизационный конвейер с 15 стартами NelderMead и до 10 000
итерациями L-BFGS требует автоматизированного контроля сходимости,
поскольку в оперативном режиме нет возможности ручной верификации
каждого результата.

Эти ограничения не являются принципиально непреодолимыми: для каждого из
них существуют технические решения (ранняя остановка разметки волн по
порогу абсолютного числа случаев; использование «предположительного»
варианта на основе общегеографического календаря доминирования;
автоматизированный флаг проверки сходимости по кратности прогонов с
независимыми стартами). Разработка оперативной версии конвейера с этими
расширениями представляется реалистичной задачей в рамках одного-двух
человеколет прикладной разработки.

\subsubsection{4.3.5 Включение в систему поддержки принятия
решений}\label{ux432ux43aux43bux44eux447ux435ux43dux438ux435-ux432-ux441ux438ux441ux442ux435ux43cux443-ux43fux43eux434ux434ux435ux440ux436ux43aux438-ux43fux440ux438ux43dux44fux442ux438ux44f-ux440ux435ux448ux435ux43dux438ux439}

В более широком контексте государственного санитарно-эпидемиологического
надзора предложенная методология занимает нишу инструментов
среднесрочного анализа --- между оперативными дашбордами суточной
заболеваемости и долгосрочными демографическими прогнозами. Её
характерное время --- от одной до нескольких недель после идентификации
пика волны --- соответствует горизонту принятия решений о режиме работы
первичного звена здравоохранения, введении или отмене дополнительных
ставок в стационарах и приоритизации вакцинационных кампаний.
Формализованный вывод о значимом \(\hat{\alpha}_k < 1\) может служить
количественным подтверждением тезиса о том, что текущая волна
развивается в условиях выраженного популяционного иммунитета --- тезиса,
влияющего как на оценку ожидаемой тяжести, так и на интерпретацию
наблюдаемых клинических тенденций в рамках регионального
здравоохранения.

Следует, тем не менее, подчеркнуть: любое операциональное использование
\(\hat{\alpha}_k\) должно сопровождаться явным указанием на
феноменологический характер данного параметра и ограниченность его
интерпретируемости в отсутствие независимых иммунологических данных
(раздел 3.5.5). Включение модельных оценок в документы органов
здравоохранения без надлежащей фиксации неопределённости и без системы
последующей верификации несёт риск ложного ощущения количественной
определённости там, где её нет.

\subsection{4.4 Ограничения исследования и направления будущих
исследований}\label{ux43eux433ux440ux430ux43dux438ux447ux435ux43dux438ux44f-ux438ux441ux441ux43bux435ux434ux43eux432ux430ux43dux438ux44f-ux438-ux43dux430ux43fux440ux430ux432ux43bux435ux43dux438ux44f-ux431ux443ux434ux443ux449ux438ux445-ux438ux441ux441ux43bux435ux434ux43eux432ux430ux43dux438ux439}

\subsubsection{4.4.1 Ограничения данных: занижение регистрируемой
заболеваемости}\label{ux43eux433ux440ux430ux43dux438ux447ux435ux43dux438ux44f-ux434ux430ux43dux43dux44bux445-ux437ux430ux43dux438ux436ux435ux43dux438ux435-ux440ux435ux433ux438ux441ux442ux440ux438ux440ux443ux435ux43cux43eux439-ux437ux430ux431ux43eux43bux435ux432ux430ux435ux43cux43eux441ux442ux438}

Фундаментальным ограничением настоящего исследования является системный
характер занижения регистрируемой заболеваемости, присущий данным
рутинного санэпиднадзора. По оценкам Кобака {[}17{]}, основанным на
анализе избыточной смертности, реальная заболеваемость COVID-19 в России
превышала официально регистрируемую в 5--10 раз, причём коэффициент
занижения существенно варьировал между волнами в зависимости от
доступности тестирования и тяжести течения доминирующего варианта.
Следствием этого является то, что абсолютные значения подобранных
параметров (\(\hat{\beta}_k\), \(\hat{\mu}_k\),
\(\hat{N}_{\text{eff},k}\)) не могут интерпретироваться как оценки
истинных биологических параметров передачи: они отражают кинетику
регистрируемых случаев, а не реального инфицирования. В частности,
эффективный размер популяции \(\hat{N}_{\text{eff},k}\) является нижней
оценкой числа людей, реально подверженных риску в данной волне.

Применительно к основному параметру интереса --- \(\hat{\alpha}_k\) ---
данное ограничение проявляется менее остро, поскольку субдиффузионный
характер динамики кодируется в форме эпидемической кривой (соотношении
длительностей фаз подъёма и спада), а не в абсолютном масштабе
заболеваемости. Тем не менее систематически заниженный масштаб кривых
мог привести к смещению \(\hat{\alpha}_k\) в том случае, если
коэффициент занижения сам претерпевал динамические изменения внутри
волны --- например, нарастал в фазе подъёма вследствие временного
переполнения тест-центров. Верификация этого эффекта в отсутствие
независимых данных невозможна, поэтому все содержательные выводы о
\(\hat{\alpha}_k\) корректны лишь условно при гипотезе о примерно
постоянном коэффициенте регистрации внутри каждой отдельной волны.

\subsubsection{4.4.2 Ограничения модельного класса: идентифицируемость и
коллинеарность}\label{ux43eux433ux440ux430ux43dux438ux447ux435ux43dux438ux44f-ux43cux43eux434ux435ux43bux44cux43dux43eux433ux43e-ux43aux43bux430ux441ux441ux430-ux438ux434ux435ux43dux442ux438ux444ux438ux446ux438ux440ux443ux435ux43cux43eux441ux442ux44c-ux438-ux43aux43eux43bux43bux438ux43dux435ux430ux440ux43dux43eux441ux442ux44c}

Проблема идентифицируемости SEIRD-модели, детально проанализированная
Королёвым {[}3{]}, в полной мере применима к настоящему исследованию.
Параметры \(\hat{\beta}_k\), \(\hat{N}_{\text{eff},k}\) и
\(\hat{\alpha}_k\) образуют коллинеарную тройку: при фиксированной форме
кривой \(I(t)\) уменьшение \(\beta_k\) может компенсироваться
уменьшением \(N_{\text{eff},k}\) с одновременным изменением \(\alpha_k\)
таким образом, чтобы функция потерь оставалась практически неизменной.
Применённые меры по снижению коллинеарности --- аналитическое
предобусловливание начальных условий из уравнения финального размера
эпидемии и фиксация \(\sigma\), \(\gamma\) из штаммовых данных --- не
устраняют проблему полностью, но существенно ограничивают пространство
решений. Тем не менее профили функции потерь по \(\alpha_k\) в ряде
коротких волн (T \textless{} 60 дней) обнаруживали уплощённые минимумы с
относительно широкими долинами, что свидетельствует о неполной
идентифицируемости на малых выборках.

В этой связи 95\%-е доверительные интервалы для \(\hat{\alpha}_k\),
полученные из профиля функции потерь (раздел 3.3.1), следует
интерпретировать как нижнюю оценку реальной неопределённости: они
отражают кривизну функции потерь в окрестности найденного минимума, но
не учитывают неопределённость вследствие коллинеарности с
\(\hat{N}_{\text{eff},k}\) и \(\hat{\beta}_k\). Байесовский вывод с
явными совместными априорными распределениями на
\((\beta, N_{\text{eff}}, \alpha)\) позволил бы получить более
корректную характеристику неопределённости, однако потребовал бы
существенного усложнения вычислительного конвейера.

\subsubsection{4.4.3 Вычислительные ограничения
GL-схемы}\label{ux432ux44bux447ux438ux441ux43bux438ux442ux435ux43bux44cux43dux44bux435-ux43eux433ux440ux430ux43dux438ux447ux435ux43dux438ux44f-gl-ux441ux445ux435ux43cux44b}

Явная GL-схема с шагом \(h = 1\) день имеет вычислительную сложность
\(O(T^2)\), поскольку на каждом шаге суммируются взвешенные состояния
системы за все предшествующие дни. Для типичных длин волн
\(T \approx 30\)--\(150\) дней и числа волн \(K = 10\) это не
представляет практического препятствия: полная оптимизация по одной
волне выполнялась за {[}t\_opt{]} секунд на стандартном вычислительном
узле. Однако при масштабировании метода на многорегиональные наборы
данных (все субъекты РФ, N ≈ 85 регионов, 10 волн каждый) суммарная
вычислительная нагрузка возрастёт пропорционально и может потребовать
параллельной реализации или применения ускорённых алгоритмов
суммирования GL-весов (быстрое преобразование Фурье, схемы с усечённой
памятью). Разработка вычислительно эффективной реализации для
многорегионального анализа является актуальной технической задачей.

\subsubsection{4.4.4 Отсутствие пространственного и демографического
измерения}\label{ux43eux442ux441ux443ux442ux441ux442ux432ux438ux435-ux43fux440ux43eux441ux442ux440ux430ux43dux441ux442ux432ux435ux43dux43dux43eux433ux43e-ux438-ux434ux435ux43cux43eux433ux440ux430ux444ux438ux447ux435ux441ux43aux43eux433ux43e-ux438ux437ux43cux435ux440ux435ux43dux438ux44f}

Модель предполагает гомогенное перемешивание популяции Пскова, не
разделяя её по возрасту, профессиональным группам, географическому
распределению внутри города или степени иммунизации. Известно, что этот
класс упрощений систематически смещает оценки \(R_0\) вверх, поскольку
реальная контактная структура снижает эффективную скорость передачи
относительно теоретической {[}2{]}. Структурированные по возрасту
SEIRD-модели (MSEIRD) частично снимают это ограничение, однако требуют
возрастной разбивки данных о заболеваемости, которая в используемых
источниках недоступна на уровне города. Отсутствие пространственного
измерения означает также, что модель не способна воспроизвести динамику,
при которой отдельные микрорайоны или учреждения выступали локальными
очагами с собственной кинетикой распространения.

\subsubsection{4.4.5 Невозможность независимой валидации
α\_k}\label{ux43dux435ux432ux43eux437ux43cux43eux436ux43dux43eux441ux442ux44c-ux43dux435ux437ux430ux432ux438ux441ux438ux43cux43eux439-ux432ux430ux43bux438ux434ux430ux446ux438ux438-ux3b1_k}

Важнейшим нерешённым ограничением является отсутствие независимого
внешнего критерия для верификации предложенной интерпретации
\(\hat{\alpha}_k\) как маркера популяционного иммунитета. Прямая
верификация потребовала бы данных серологических обследований населения
Пскова, синхронизированных по времени с границами эпидемических волн,
--- данных, которые в открытом доступе для данного города не
представлены. Косвенные свидетельства --- корреляция \(\hat{\alpha}_k\)
с кумулятивной заражённостью, \(\hat{N}_{\text{eff},k}\) и индексом
асимметрии кривой (раздел 3.5) --- согласуются с иммунологической
гипотезой, однако не являются её строгим доказательством. Альтернативные
объяснения убывания \(\hat{\alpha}_k\) --- нарастающая поведенческая
реакция населения, сезонные изменения восприимчивости, артефакты
регистрации --- не могут быть исключены в рамках настоящего дизайна.

\subsubsection{4.4.6 Направления будущих
исследований}\label{ux43dux430ux43fux440ux430ux432ux43bux435ux43dux438ux44f-ux431ux443ux434ux443ux449ux438ux445-ux438ux441ux441ux43bux435ux434ux43eux432ux430ux43dux438ux439}

Полученные результаты открывают несколько взаимосвязанных перспектив.
Первое направление --- пространственное обобщение метода: применение
поволнового дробного SEIRD-анализа к данным по всем или репрезентативной
выборке субъектов РФ позволит проверить, является ли наблюдаемая
закономерность убывания \(\hat{\alpha}_k\) универсальной для российских
городов различного размера и географического положения, или специфична
для северо-западного региона. Второе направление --- байесовская
реализация модели с полным апостериорным распределением по
\((\beta, N_{\text{eff}}, \alpha)\), которая обеспечит корректную
совместную характеристику неопределённости параметров и позволит
тестировать гипотезы об изменении \(\alpha_k\) между волнами с учётом
коллинеарности.

Третье направление --- включение данных о вакцинации. Официальная
статистика охвата вакцинацией по субъектам РФ доступна с конца 2020 года
и позволяет ввести в модель явный компартмент вакцинированных (V) с
варьирующей по времени и штаммам эффективностью. Это даст возможность
частично разделить вклады вакцинального и постинфекционного иммунитета в
наблюдаемое убывание \(\hat{\alpha}_k\) --- что является ключевым
условием для превращения \(\alpha_k\) из феноменологического маркера в
верифицируемую иммунологическую характеристику. Наконец, прогностическое
применение \(\hat{\alpha}_k\) --- оценка данного параметра по ранней
(первые 14--21 день) восходящей ветви текущей волны с целью предсказания
формы её спада и совокупной нагрузки на систему здравоохранения ---
представляет значительный практический интерес и является логичным
следующим шагом на пути от ретроспективного к оперативному
аналитическому применению разработанного подхода.

\section{5.
Заключение}\label{ux437ux430ux43aux43bux44eux447ux435ux43dux438ux435}

Настоящее исследование демонстрирует, что полный цикл ретроспективного
поволнового анализа пандемии COVID-19 --- от сырых файлов
государственной эпидемиологической отчётности до формально сравниваемых
компартментных моделей --- выполним на основе исключительно стандартных
данных рутинного санитарно-эпидемиологического надзора, без привлечения
геномного мониторинга, серологических регистров или специализированных
баз данных. Воспроизводимый конвейер, реализованный на языке Julia и
апробированный на 1750-дневном ряду заболеваемости COVID-19 в г. Пскове
(2020--2025 гг.), охватывает все этапы от импутации пропусков и
автоматической разметки волн до поволновой калибровки конкурирующих
моделей и их формальной селекции, и принципиально переносим на любой
российский город с населением от 50 000 человек при наличии стандартных
форм федерального статистического наблюдения.

Основным статистическим результатом является следующее. Дробная
SEIRD-модель с производной Капуто порядка \(\alpha_k\), откалиброванная
независимо для каждой из десяти идентифицированных волн при
зафиксированных из штаммовых данных параметрах \(\sigma\) и \(\gamma\),
статистически значимо превосходит целочисленную SEIRD-модель в пяти
волнах (3, 7, 8, 9, 10) по совокупности LRT с поправкой Бонферрони, AIC
и BIC. Для ранних волн (1--2) оба класса моделей неразличимы, что
содержательно интерпретируется как отсутствие иммунологической «памяти»
в полностью восприимчивой популяции, а не как недостаточная мощность
критерия. Монотонное нарастание \(\Delta\text{AIC}\) в волнах 7--10
(поздние субварианты Омикрон) достигает значений, соответствующих
критерию «исключительно сильного свидетельства» по шкале
Бёрнхэма--Андерсона. Сопоставление дробной GL-модели с высокоточным
ОДУ-решателем (Verner--Rodas5P) исключает объяснение наблюдаемого
преимущества артефактами GL-дискретизации: разница AIC между
ОДУ-решателем и дробной GL-моделью воспроизводит те же паттерны, что и
LRT между целочисленной и дробной GL-схемами.

Совокупность четырёх эмпирических критериев --- хронологического
(монотонное убывание \(\hat{\alpha}_k\) с номером волны), биологического
(корреляция с кумулятивной заражённостью), модельного (совместное
убывание \(\hat{\alpha}_k\) и \(\hat{N}_{\text{eff},k}\)) и
диагностического (отрицательная корреляция \(\hat{\alpha}_k\) с индексом
асимметрии кривой) --- взаимно согласована и поддерживает интерпретацию
\(\hat{\alpha}_k < 1\) как косвенного феноменологического маркера
субдиффузионного режима, обусловленного нарастанием популяционного
иммунного фона. Локальное «возвратное» значение \(\hat{\alpha}_4\) при
волне Омикрон BA.1 не опровергает эту гипотезу, а уточняет её через
механизм частичного иммунного уклонения нового субварианта от
предсуществующего иммунитета. Временны́е масштабы ядра памяти GL при
характерных значениях \(\hat{\alpha}_k\) поздних волн (60--90 дней)
совпадают с горизонтами клинически значимого гуморального ответа, что
усиливает правдоподобность иммунологической интерпретации. Вместе с тем
данная интерпретация остаётся феноменологической: GL-оператор агрегирует
разнородные «памятные» механизмы в единый скаляр, и в отсутствие
серологических данных по Пскову механистическое разложение суммарного
вклада остаётся недостижимым.

В части научной новизны вклад настоящей работы состоит \textbf{не в
применении дробного исчисления к задачам эпидемиологии как таковом, а в
одновременной реализации трёх методологических решений, которые в
цитируемой литературе не встречаются совместно: поволновая оценка
\(\alpha_k\) с биологически обоснованными \(\sigma\), \(\gamma\) в роли
информированного prior, и формальная попарная селекция с коррекцией на
множественные сравнения.} Именно это сочетание позволяет не просто
констатировать глобальное преимущество дробной модели, но и локализовать
его во времени --- связать с конкретными волнами и конкретным
иммунологическим контекстом.

Среди ограничений исследования наиболее существенными остаются
\textbf{системное занижение регистрируемой заболеваемости,
коллинеарность параметров
\((\hat{\beta}_k, \hat{N}_{\text{eff},k}, \hat{\alpha}_k)\) и
невозможность независимой серологической верификации
\(\hat{\alpha}_k\).} Тем не менее именно в условиях, когда
серологические и геномные данные недоступны, а единственным источником
информации служат формы стандартной отчётности, параметр \(\alpha_k\)
приобретает практическую ценность нового класса: он извлекаем из данных,
которые уже собираются повсеместно и без дополнительных затрат, и
открывает возможность систематического косвенного мониторинга
популяционного иммунного фона в масштабе всей страны --- там, где
серология недоступна, недостаточна или опаздывает.

\end{document}